% This LaTeX document needs to be compiled with XeLaTeX.
\documentclass[10pt]{article}
\usepackage[utf8]{inputenc}
\usepackage{amsmath}
\usepackage{amsfonts}
\usepackage{amssymb}
\usepackage[version=4]{mhchem}
\usepackage{stmaryrd}
\usepackage{graphicx}
\usepackage[export]{adjustbox}
\graphicspath{ {./images/} }
\usepackage{bbold}
\usepackage[fallback]{xeCJK}
\usepackage{polyglossia}
\usepackage{fontspec}
\usepackage{newunicodechar}
\IfFontExistsTF{Noto Serif CJK SC}
{\setCJKmainfont{Noto Serif CJK SC}}
{\IfFontExistsTF{STSong}
  {\setCJKmainfont{STSong}}
  {\IfFontExistsTF{Droid Sans Fallback}
    {\setCJKmainfont{Droid Sans Fallback}}
    {\setCJKmainfont{SimSun}}
}}

\setmainlanguage{english}
\IfFontExistsTF{CMU Serif}
{\setmainfont{CMU Serif}}
{\IfFontExistsTF{DejaVu Sans}
  {\setmainfont{DejaVu Sans}}
  {\setmainfont{Georgia}}
}

\author{Q.E.D.}
\date{}


\newunicodechar{⇒}{\ifmmode\Rightarrow\else{$\Rightarrow$}\fi}
\newunicodechar{α}{\ifmmode\alpha\else{$\alpha$}\fi}
\

\begin{document}
\title{GENERALIZED WHITTAKER VECTORS AND REPRESENTATION THEORY}
\maketitle
GENERALIZED WHITTAKER VECTORS AND REPRESENTATION THEORY by

Thomas Emile Lynch\\
B.S., University of Georgia\\
(1973)

SUBMITTED IN PARTIAL FULFILLMENT\\
OF THE REQUIREMENTS FOR THE\\
DEGREE OF

DOCTOR OF PHILOSOPHY\\
at the

MASSACHUSETTS INSTITUTE OF TECHNOLOGY\\
June 1979

Signature of Author\\
Department of Mathematics, May 4, 1979

Certified by...erac........................

\begin{itemize}
  \item Thesis Supervisor
\end{itemize}

Accepted by\\
Chairman, Departmental Committee\\
on Graduate Students

\section{ARCHIVES}
MASSACHISETTS INSTITUTE\\
OF TECHNOLOGY\\
JUN 7 19̂79\\
LIBRARIES

GENERALIZED WHITTAKER VECTORS AND REPRESENTATION THEORY by

Thomas Emile Lynch\\
Submitted to the Department of Mathematics on May 4, 1979 in partial fulfillment of the requirements for the Degree of Doctor of Philosophy.

\section{ABSTRACT}
Let $U$ be the enveloping algebra of a semi-simple Lie algebra, $f$. Let $p$ be a parabolic subalgebra of of and let $\cap$ be its nilradical. In this paper, we study those $U$-modules, $V$, containing vectors which transform as certain (admissible) characters when acted upon by $\varkappa$. We call such vectors Whittaker vectors. If V is generated by Whittaker vectors, then V is determined by its Whittaker vectors viewed as a $U(\bar{\rho})^{\mathrm{N}}$-module, where $U(\bar{p})^{N}$ is a certain subalgebra of $U$.

Now, modules generated by Whittaker vectors have certain nice algebraic properties and we use these properties to calculate the dimensions of the spaces of Whittaker vectors in generalized Verma modules and the duals of HarishChandra modules arising from induction on minimal parabolics.

Thesis Supervisor: Bertram Kostant Title: Professor of Mathematics

\section{Contents}
Page\\
Abstract ..... 2\\
Contents ..... 3\\
Introduction\\
0.1 Some background and acknowledgments ..... 6\\
0.2 Statement of results in Chapters 1 and 2 ..... 7\\
0.3 Statement of results in Chapters 3 and 4 ..... 11\\
0.4 Statement of results in Chapters 5-8 ..... 13\\
Chapter 1 A decomposition theorem for $S(\bar{\rho})$\\
1.1 Some preliminaries on parabolic subalgebras ..... 16\\
and the element $\mathrm{x}_{0}$\\
1.2 Admissible nilpotents and the isomorphism ..... 19\\
$N \times(f+s) \rightarrow f+f$\\
1.3 Some definitions due to Kostant and Kazhdan ..... 24\\
1.4 The graded action of $n$ on $S(\bar{\rho})$ ..... 26\\
1.5 The graded decomposition $S(\bar{\rho})=A \otimes S(\bar{\rho})^{N}$ ..... 30\\
1.6 The structure of A ..... 35\\
1.7 The structure of $S(\bar{\rho})^{N}$ ..... 38\\
1.8 The restriction of $\mathrm{S}(\bar{p})^{\mathrm{N}}$ to $\mathrm{S}(m)$ ..... 43\\
1.9 Comments ..... 45\\
Chapter 2 A decomposition theorem for ..... $U(\bar{p})$\\
2.1 The $\mathrm{x}_{0}$-filtration of $\mathcal{U}(\bar{\rho})$ and $\operatorname{Gr} U(\bar{\rho})$ ..... 46\\
2.2 The $n$-reduced action of $n$ or $U(\bar{\rho})$ ..... 50\\
page\\
2.3 The structure of $U(\bar{p})^{\mathrm{N}}$ ..... 54\\
2.4 The decomposition $U(\bar{p})=\tilde{\mathrm{A}} \otimes l(\bar{p})^{\mathrm{N}}$ ..... 64\\
2.5 The isomorphism $A \rightarrow A$ of $N$-modules ..... 66\\
2.6 The image of in $U(\bar{p})^{\mathrm{N}}$ ..... 70\\
2.7 Comments ..... 72\\
Chapter 3 Whittaker modules\\
3.1 The annihilator of a Whittaker vector is ..... 75$U U_{w}(\bar{p})^{N}+U \dot{U}(n)$\\
3.2 One has $W h \mathrm{v}=2(\bar{p})^{\mathrm{N}_{\mathrm{w}}}$ for a cyclic $w \varepsilon \mathrm{v}$ ..... 80\\
3.3 Whittaker modules and cyclic $U(\bar{\rho})^{\mathrm{N}}$-modules ..... 81\\
3.4 The irreducibility of Whittaker modules ..... 82\\
Chapter 4 Tnesor products and $\eta$-finite modules\\
4.1 The $n$-isomorphism $\mathrm{V} \rightarrow ?(n)^{*} \otimes$ Wh(V) ..... 85\\
4.2 Tensor products and Whittaker vectors ..... 89\\
4.3 $\eta$-finite modules ..... 92\\
4.4 Composition series and Whittaker modules ..... 97\\
4.5 Submodules generated by Whittaker modules ..... 99\\
4.6 A formula for dim Wh (M') ..... 102\\
Chapter 5 Whittaker vectors and Verma modules\\
5.1 The completions of generalized Verma modules 108\\
5.2 The annihilators of Whittaker modules ..... 11.5\\
5.3 Comments ..... 116\\
page\\
\includegraphics[max width=\textwidth, alt={}, center]{dc0495e0-61cd-4445-9972-03b906b7ff05-005_54_1234_451_319}\\
6 . 1 Preiiminaizies ..... 19\\
\includegraphics[max width=\textwidth, alt={}]{dc0495e0-61cd-4445-9972-03b906b7ff05-005_82_1132_609_442} ..... 122\\
6. 3 COROIOIL OONEINIALEIOIN OI CIAIOOTEIS ..... 126\\
\includegraphics[max width=\textwidth, alt={}]{dc0495e0-61cd-4445-9972-03b906b7ff05-005_93_1055_770_438} ..... 130\\
\includegraphics[max width=\textwidth, alt={}]{dc0495e0-61cd-4445-9972-03b906b7ff05-005_54_1234_877_315} Q1OMONES\\
7.1 SOINO OBMOROI I LIQS ..... 134\\
\includegraphics[max width=\textwidth, alt={}]{dc0495e0-61cd-4445-9972-03b906b7ff05-005_50_1136_1132_438} ..... 138\\
\includegraphics[max width=\textwidth, alt={}]{dc0495e0-61cd-4445-9972-03b906b7ff05-005_54_985_1218_562} AlgOBIOS\\
7.3 OPMOILIOIS FOR LIO EXISLOIOO OI IOOI ..... 146\\[0pt]
AOMISSIBIQ NINDOLOILS FOP ElO OIOSSIOAI. A]CBbIAQ\\
7.4 SOINO OOMMONLO OI TOAI AOINISSIVIO ..... 752\\
\includegraphics[max width=\textwidth, alt={}, center]{dc0495e0-61cd-4445-9972-03b906b7ff05-005_54_1064_1729_557}\\[0pt]
Ghapter 8 Hxame]es\\
\includegraphics[max width=\textwidth, alt={}]{dc0495e0-61cd-4445-9972-03b906b7ff05-005_58_1153_1895_434} ..... 154\\
\includegraphics[max width=\textwidth, alt={}]{dc0495e0-61cd-4445-9972-03b906b7ff05-005_54_1059_1980_434} ..... $9 x(n \cdot K)] 59$\\[0pt]
8. 3 Minima] DOPEBOI C SIVEIOWOIORS ..... 161\\
\includegraphics[max width=\textwidth, alt={}, center]{dc0495e0-61cd-4445-9972-03b906b7ff05-005_54_715_2155_557}\\
RECOTOROQ ..... 163\\
Biogidginioal note ..... 260

\section{Introduction}
0.1. Some background and acknowledgments.

Let $o f$ be a complex, semi-simple Lie algebra and let $p=m \oplus \Omega \oplus n$ be a Langlands decomposition of a parabolic subalgebra of $\sigma$. Let 放 denote the enveloping algebra of $O \gamma$. Then, generalized whittaker theory is the study of (-modules, $V$, for which there exists $w \varepsilon V$ such that there is an admissible homomorphism (see §2.2) $\eta: \eta \rightarrow c$ for which $x w=\eta(x) w$ for all $x \in \eta$. Such vectors are called whittaker vectors in analogy to the theory developed by Kostant in $[9]$. Under the assumption of admissibility for $\eta$, modules containing Whittaker vectors are shown to have very tractable properties that will be developed in this paper.

My interest in generalized Whittaker theory arose from an attempt to understand Kostant's paper [9]. In particular, I wished to remove the dependence on results concerning orbits of regular elements in a Lie algebra and the centers of enveloping algebras from the proofs in that paper. The first result was a new proof of Theorem 1.2 (in both papers) which was substantially simpler than Kostant's. To compensate, the proof of Theorem 2.3 (2.4.1 in [9]) becomes extremely complicated.

With the new proof of Theorem 1.2, it became clear that Whittaker theory was valid for many other parabolic subalgebras than just the Borel subalgebra as was developed in [9]. Probably the chief result in this paper is the definition of a subalgebra $U(\bar{p})^{\mathrm{N}}$ which has the properties associated to the (Whittaker model) $W$ in $[9]$.

Many of the proofs and results in this paper are straightforward generalizations of the material in Kostant's paper [9]. Much of the exposition is simply paraphrased from that paper since any attempt to change it would not positively affect lucidity.

I wish to thank Bertram Kostant, Alvany Rocha-Carida, Dave Vogan, Don King, Dale Peterson and Joe Johnson for helpful conversations. I also wish to thank Phyllis Ruby for being so patient with my chronic inability to meet deadlines in typing this paper. During my graduate career at M.I.T., I was supported by a National Science Foundation Fellowship for two years and a research assistantship this past year. For this, I am eternally grateful.\\
0.2 . Statement of results in Chapters 1 and 2.

The representation theory of Chapters 3-6 are based on certain decomposition theorems developed in

Chapters 1 and 2. If $\bar{p}=m \oplus \sigma \oplus \bar{n}$ is the opposite parabolic subalgebra to $p$, then corresponding to an admissible nilpotent element $f \varepsilon \bar{n}$ (see §1.2) there is an action of N (the unipotent subgroup of $\operatorname{Ad}(\mathrm{og})$ corresponding to $N$ ) on the symmetric algebra $S(\bar{\rho})$ of $\bar{\rho}$. If $S(\bar{\rho})^{N}$ is the algebra of invariants for this action, then by Theorem 1.2 (a cross-section theorem for the orbits of $N$ in $p$ ) one shows that (Theorem 1.4)


\begin{equation*}
S(\bar{p})=A \otimes S(\bar{p})^{N} \tag{0.2.1}
\end{equation*}


where $A$ transforms under $N$ like the affine algebra, $A(N)$, of $N$ under left translation. Under a grading due to Kazhdan in the case where $p$ is a Borel subalgebra, (0.2.1) is shown to be a graded isomorphism. In §l.7, $S(\bar{\rho})^{N}$ is shown to be isomorphic to a polynomial ring on $\operatorname{dim}(m)$ generators and its Poincaré polynomial is calculated. The final result of this section is that the restriction $S(\bar{p})^{N} \rightarrow S(m)$ is injective and that the corresponding complex variety is just $\{0\}$.

The proof of Theorem 1.2 is a considerable simplification of the corresponding proof in [9]. In fact,\\[0pt]
one can use Theorem 1.2 to prove some results concerning the orbits of regular elements. However, the resulting proofs are not much simpler than those in [7] and will be omitted.

Corresponding to the gradation of $S(\bar{p})$, there is a filtration of the enveloping algebra of $\bar{p}, \Downarrow(\bar{p})$, also due to Kazhdan. (See §2.1). Given an admissible nilpotent $f \varepsilon \bar{n}$, we may associate to $f$ a character $\eta: \eta \rightarrow C$ via the Killing form. We call such a character admissible. This character leads to an action of $n$ on $U(\bar{p})$ and we let $U(\bar{\rho})^{N}$ Denote the space of $n$-invariants in $U(\bar{\rho})$. It is shown that the action of $N$ on $S(\bar{p})$ corresponds to the action of $n$ on $\mathcal{U}(\bar{p})$ and that we have short exact sequences, $k=1,2, \ldots$

$$
\text { (0.2.2) } 0 \rightarrow U_{(k-1)}(\bar{p})^{N} \xrightarrow{\text { inj }} U_{(k)}(\bar{p})^{N} \xrightarrow{\tau}(k) s_{(k)}(\bar{p})^{N} \rightarrow 0
$$

since we can show that $U(\bar{\rho})^{\mathrm{N}}$ is a graded subalgebra of $U(\bar{p})$. In particular, $U(\bar{p})^{N}$ has a Birkhoff-Witt basis consisting of $\operatorname{dim}(m)$ generators and the restriction map $\mathcal{U}(\bar{p})^{\mathrm{N}} \rightarrow U(m)$ is injective. (Here we let $p=m \oplus \pi$ for convenience.) Showing that $\tau_{(k)}$ in (0.2.2) is onto (Theorem 2.3) is by far the most difficult proof in this paper and graphically illustrates the\\
difficulties arising from the non-commutativity of $U(\bar{\rho})^{\mathrm{N}}$ in the general case. (In fact, $U(\bar{\rho})^{N}$ is commutative if and only if $\hat{\rho}$ is a Borel subalgebra.)

From (0.2.2) it is easy to show (Theorem 2.4) that


\begin{equation*}
U(\bar{p}) \cong \tilde{A} \otimes U(\bar{p})^{N} \tag{0.2.3}
\end{equation*}


for some graded subalgebra $\tilde{A} \subseteq U(\bar{\rho})$ which maps isomorphically onto $A$ via ${ }^{\tau}(k)$. (See §2.4.) Further, if $\eta$ and $f$ correspond to one another, we have (Theorem 2.2.)\\
(0.2.4)

$$
\operatorname{Gr} U(\bar{p}) \cong s(\bar{p})
$$

as $n$-modules.\\
It is interesting to note that we have a commutative diagram:

$$
\begin{aligned}
(0.2 .5) & 0 \rightarrow \gamma_{k-1} \rightarrow \sum_{k} \rightarrow \underset{\downarrow}{k} \rightarrow s_{k}(\sigma)^{G} \rightarrow 0 \\
& 0 \rightarrow U_{(k-1)}(\bar{p})^{N} \rightarrow U_{(k)}(\bar{p})^{N} \rightarrow s_{(k)}(\bar{p})^{N} \rightarrow 0 \\
& +\underset{\downarrow}{0} \rightarrow U_{k-1}(m) \rightarrow U_{k}(m) \rightarrow s_{k}(m) \rightarrow 0
\end{aligned}
$$

where the horizontal lines are exact sequences and the vertical lines are injections. (See §2.7.). We are thus led to consider the map $U_{(\bar{p})^{N} \rightarrow U(m) \text { as a (factorization }}$ of a) generalized Harish-Chandra homomorphism. This map wll be very important in the study of Whittaker vectors in the completions of generalized Verma modules.\\
0.3. Statement of results in Chapters 3 and 4.

Let $\eta: \eta \rightarrow C$ be a fixed admissible admissible homomorphism. If $V$ is a $U$-module, we say that $w \in V$ is a Whittaker vector if $\mathrm{xw}=\eta(\mathrm{x}) \mathrm{w}$ for all $\mathrm{x} \varepsilon \lambda$ and that V is a Whittaker module if $\ell \mathrm{w}=\mathrm{V}$. Uis the enveloping algebra of $\left.g_{.}\right)$Let $U_{\eta}(\eta)$ be the kernel of the extension $n: \hat{U}(n) \rightarrow c$. Now, if $\emptyset_{w}$ is the annihilator in Q of a Whittaker vector $0 \neq \mathrm{w} \varepsilon \mathrm{V}$, then we have (Theorem 3.1).


\begin{equation*}
U_{w}=U U_{w}(\bar{p})^{N}+U U_{n}(r) \tag{0.3.1}
\end{equation*}


where $U_{w}(\bar{p})^{N}$ is the annihilator of $w$ in $U(\bar{\rho})^{N}$. From this, it follows that (Theorem 3.2) the space of Whittaker vectors in V is a $U(\bar{\rho})^{\mathrm{N}}$-module and that the collection of Whittaker modules is naturally parameterized by the collection of cyclic $U(\bar{p})^{\mathrm{N}}$-modules\\
(Theorem 3.3). Further, (Theorem 2.4) a Whittaker module is irreducible if and only if its Whittaker vectors form an irreducible $U(j)^{N}$-module.

In Chapter 4 we study the algebraic properties of U-modules containing Whittaker vectors. First, if V is any $U$-module and $F$ is any finite-dimensional $U$-module, then the space of Whittaker vectors in $V \otimes F$ is naturally parameterized by elements of the form $\mathrm{w} \otimes \mathrm{z}$ where $w$ is a Whittaker vector in $V$ and $z \varepsilon F$. (See Theorem 4.2). In particular, $\operatorname{dim} W h(V \otimes F)<\infty$ if and only if $\operatorname{dim}$ Wh $(V)<\infty$ and we have


\begin{equation*}
\operatorname{dim} W h(V \otimes F)=(\operatorname{dim} W h(V))(\operatorname{dim} F) . \tag{0.3.2}
\end{equation*}


Next, we show that if $V$ is of the form $V=W h(V)$, then $V$ has a (finite) composition series of length $k$ as a U-module if and only if Wh (V) has a (finite) composition series of length $k$ as a $U(\bar{\rho})^{N}$-module (Proposition 4.5.) Moreover, in such a case, if $0=\mathrm{v}_{0} \subseteq \ldots \subseteq \mathrm{v}_{\mathrm{i}-1} \subseteq \mathrm{v}_{\mathrm{i}} \subseteq \ldots \subseteq \mathrm{v}_{\mathrm{k}}=\mathrm{v}$ is a composition series for $V$, then (Theorem 4.4.)


\begin{equation*}
\operatorname{dim} W h(V)=\sum_{i=1}^{k} \operatorname{dim} W h\left(V_{i} / V_{i-1}\right) . \tag{0.3.3}
\end{equation*}


Lastly, we have a similar result to (0.3.3) (Theorem 4.6) relating the dimension of the space of Whittaker vectors of a module with the dimensions of the spaces of Whittaker vectors in the subquotients of a composition series for the module. This result is rather unsatisfactory since it depends on the finiteness of the space of Whittaker vectors and cannot immediately be generalized.\\
0.4. Statement of results in Chapters 5-8.

In Chapter 5 we study the existence of Whittaker vectors in the completions of generalized Verma modules. In particular, if $\mathrm{V}=\mathrm{V}_{\mathrm{H}}$ is an irreducible generalized Verma module induced from a finite-dimensional $U(M)$-module $H$, then (Theorem 5.1)\\
(0.4.1)

$$
\operatorname{dim} W h(\bar{V})=\operatorname{dim}(H)
$$

where $\bar{V}$ is the completion of $V$. (See §5.1.) Further, $W h(\bar{V})$ is isomorphic to $H$ as a $U(\bar{p})^{N}$-module where $U(\bar{p})^{\mathrm{N}}$ acts on H via the generalized Harish-Chandra homomorphism $U(\bar{\rho})^{N} \rightarrow U(m) \subseteq U(o)$. Further, if $w$ is a Whittaker submodule of $\bar{V}$ as above, we can show that the annihilator of $W$ in $U$ is the same as the\\
annihilator of V in $U$. (Theorem 5.2.) Thus, Whittaker modules are in some sense "maximally faithful" with respect to $p \subseteq O \gamma$. There are many unsolved problems concerning Whittaker submodules of completions of Verma modules that we discuss in this chapter.

In Chapter 6 we apply the results of Chapter 3 and 4 to study the space of Whittaker vectors in the dual of Harish-Chandra modules induced from a minimal parabolic. In particular, we are able to show that (Theorem 6.3)


\begin{equation*}
\operatorname{dim} W h\left(X_{\sigma, v}^{\prime}\right)=|W| \operatorname{dim}(\sigma) \tag{0.4.2}
\end{equation*}


where $\sigma$ is a finite dimensional representation of and $|\mathrm{w}|$ is the cardinality of the "little" Weyl group $w(g, \Omega)$. This result relies strongly on Kostant's results concerning the $U(\pi)$-structure of the spherical principal series. The technqiue of the proof was suggested to me by Dave Vogan.

It seems likely that (0.4.2) can be generalized to cover other Harish-Chandra modules induced from finitedimensional representations of other parabolic subalgebras (in particular, the degenerate series) but this is not clear.

In Chapter 7, we show which of the parabolic subalgebras of the classical Lie algebras have admissible nilpotent elements and give examples of such elements. (In fact, we go a bit further and study those admissible nilpotents occuring in the various real forms.) In particular, all minimal parabolic subalgebras have (real) admissible nilpotent elements so that the results of Chapter 5 are non-trivial. Finally, we mention some results concerning admissible nilpotents in the exceptional Lie algebras.

Chapter 8 is devoted to some examples of parabolic subalgebras for which the structure of $U(\bar{p})^{N}$ can be understood reasonably. However, even in these simple cases the structure is not entirely clear and many questions of importance to the study of the irreducible representations of $U(\bar{p})^{N}$ remain unsolved.

\section{Chapter 1}
\section{A decomposition theorem for $S(\bar{p})$.}
Many of the results, definitions and proofs in this chapter are simple generalizations of material in Chapter 1 of Kostant [9]. Much of the exposition of this material will need only slight modifications. Material thus obtained will be denoted by [Kostant].\\
1.1 Some preliminaries on parabolic subalgebras and the element $\mathrm{x}_{0}$.

The material in this section is standard and may be found, in a slightly different form, in Warner [ 16].

Let $g_{0}$ be a real, semi-simple Lie algebra with Cartan involution $\theta$ and an associated Iwasawa decomposition $\sigma_{0}=k_{0} \oplus \sigma_{0} \oplus \pi_{0}$. Then, $\sigma_{0}$ is an abelian subalgebra of $\gamma_{0}$ whose elements, acting by the adjoint action on $y_{0}$, are semi-simple with real eigenvalues. We have that $\pi_{0}$ is a nilpotent subalgebra of $\mathcal{J}_{0}$ normalized by $\sigma_{0}$. Let $\Delta^{+}$denote the set of roots for the pair $\left(\pi_{0}, \pi_{0}\right)$. Then $0 \notin \Delta^{+}$and one knows that $\Delta^{+}$contains a fundamental system of roots, $\pi=\left\{\alpha_{1}, \ldots, \alpha_{\ell}\right\} \quad(\ell \quad$ is the real rank of $g_{0}$ ), so that every $\phi \varepsilon \Delta^{+}$may be written uniquely as a linear combination of elements of $\pi$ with non-negative integral coefficients. Further, $\pi$ is a basis for the\\
real dual of $\Omega_{0}$. We will denote the opposite subalgebra of $n_{0}$ by $\bar{n}_{0}=\theta\left(\pi_{0}\right)$. Then $\bar{\pi}_{0}$ is a nilpotent subalgebra of $y_{0}$, normalized by $\sigma_{0}$, and the set of roots for the pair $\left(\bar{\tau}_{0}, \sigma_{0}\right)$ is $\Delta^{-}=\left\{-\phi \mid \phi \varepsilon \Delta^{+}\right\}$. If we let $\pi_{0}$ denote the centralizer of $\Omega_{0}$ in $k_{\theta}$, then $y_{0}=m_{0} \oplus r_{0} \oplus r_{0} \oplus \bar{r}_{0}$ and $\rho_{0}=m_{0} \oplus r_{0} \oplus r_{0}$ is a minimal parabolic subalgebra of of 0 .

From the above, it follows that if $\Sigma \subseteq \pi$, then we may define $x_{0} \varepsilon q_{0}$ by\\
(1.1.1) $\phi\left(\mathrm{x}_{0}\right)=\left\{\begin{array}{lllll}1 & \text { if } & \phi \varepsilon \Sigma & \\ 0 & \text { if } & \phi \varepsilon \pi & \phi \notin \Sigma .\end{array}\right.$

It is then clear that for $\phi \varepsilon \Delta^{+}$\\
(1.1.2) $\quad \phi\left(\mathrm{x}_{0}\right)=\sum_{\alpha_{\mathrm{i}} \varepsilon \Sigma} \mathrm{n}_{\mathrm{i}}(\phi)$ if $\phi=\sum_{\alpha_{i} \varepsilon \pi} \mathrm{n}_{\mathrm{i}}(\phi) \alpha_{\mathrm{i}}, \mathrm{n}_{\mathrm{i}}(\phi) \varepsilon \mathrm{Z}_{+}$. We define the length of $\phi$ (with respect to $\Sigma \subseteq \pi$ ) by $0(\phi)=\phi\left(\mathrm{x}_{0}\right)$. It is clear that $x_{0}$ is a semi-simple element of $\sigma_{0}$ with integral eigenvalues, and if $d_{i}$ and $d_{j}$ are the eigenspaces in $g_{0}$ corresponding to the eigenvalues $i, j \varepsilon Z$ under ad $\mathrm{x}_{0}$, then\\
(1.1.3) $\quad\left[d_{i}, d_{j}\right] \subseteq d_{i+j}$.

It is convenient to think of the $d_{i}$ 's as the diagonals in $\sigma_{0}$ corresponding to $\Sigma \subseteq \pi$.

If we let $p_{0}(\Sigma)=\sum_{i \geq 0} d_{i}$, then $p_{0}(\Sigma)$ is a parabolic\\
subalgebra of $y_{0}$ since $\rho_{0}(\Sigma)$ clearly contains the minimal parabolic $\rho_{0}$ and is a subalgebra by (1.1.3). It follows from a theorem of Tits (see Warner [16], p. that every parabolic subalgebra of $\mathcal{J}_{0}$ is conjugate, under the adjoint group of $\gamma_{0}$, to a unique $\rho_{0}(\Sigma)$ for some $\Sigma \subseteq \pi$. We will henceforth fix an Iwasawa decomposition of $g_{0}$ and call the $\rho_{0}(\Sigma)$ 's the standard parabolic subalgebras of of $_{0}$.

Note that $\pi_{0}(\Sigma)=\sum_{i>0} d_{i}$ is the nilpotent radical\\
of $P_{0}(\Sigma)$ and that $m_{0}(\Sigma)=\mathrm{d}_{0}$ is reductive and normalizes $n_{0}(\Sigma)$. Note that $m_{0}(\Sigma) \oplus \pi_{0}(\Sigma)$ is not a Langlands decomposition of $\rho_{0}(\Sigma)$ since $m_{0}(\Sigma)$ is the full reductive component of $P_{0}(\Sigma)$. Further, note that $\mathcal{J}_{0}=m_{0}(\Sigma) \oplus n_{0}(\Sigma) \oplus \bar{n}_{0}(\Sigma)$ where $\bar{n}_{0}(\Sigma)=\theta\left(\pi_{0}(\Sigma)\right)=\sum_{\mathrm{i}<0} \mathrm{~d}_{\mathrm{i}}$. Then, if we set $\bar{\rho}_{0}(\Sigma)=\theta\left(\rho_{0}(\Sigma)\right)=\sum_{i \leq 0} \mathrm{~d}_{\mathrm{i}}, \bar{\rho}_{0}(\Sigma)$ is the opposite parabolic subalgebra of $P_{0}(\Sigma)$.

To simplify notation, we fix $\Sigma \subseteq \pi$ and (abandoning the notation we used for the minimal parabolic above) set $p_{0}=p_{0}(\Sigma), m_{0}=m_{0}(\Sigma)$, and $n_{0}=n_{0}(\Sigma)$.

Now, let of denote the complexification of $\sigma_{0}$ (as a Lie algebra.) We define $p, \bar{p}, m, n$, and $\bar{n}$ similarly. We will denote the complexification of $d_{i}$ by $d_{i}$ and hope that no confusion arises. (In general, we will denote the real $d_{i}$ by $d_{i} \cap \sigma_{0}$.) Note that the complexification of a parabolic subalgebra of $g_{0}$ is a parabolic subalgebra of of.

Next, let $G$ denote the adjoint group of $O f$ and let $N$ denote the connected unipotent subgroup of $G$ with Lie algebra $n$. Then if $P$ is the normalizer of $N$ in $G$, the Lie algebra of $P$ is $P$ and we may write $P$ as a semi-direct product $P=M N$ where the Lie algebra of $M$ is $M$.

If $a \varepsilon G, x \varepsilon O$, then $a x \varepsilon O f$ will denote the adjoint action of a on x .

For any $x \in O, r \subseteq O$, we will let $r^{x}$ denote the centralizer of $x$ in $r$.\\
1.2 Admissible nilpotent elements and the isomorphism

$$
N \times(f+\Delta) \rightarrow f+p
$$

Now, let $\rho$ be a (standard) parabolic subalgebra of $O$ and let $x_{0}$ be as above. Recall that we have a decomposition $p=m \oplus n$ of $p$ into its reductive and nilpotent parts. We say that an element $f \varepsilon \mathcal{y}$ is admissible if

\begin{enumerate}
  \item $f \varepsilon d_{-1}$ and\\
(1.2.1)
  \item $n^{f}=\{0\}$.
\end{enumerate}

We note that since $d_{-1} \subseteq \bar{\pi}, f$ is a nilpotent element of $\sigma$.

Note that if $\rho$ is a Borel subalgebra of $g$ (that is, if $P$ is the complexification of a minimal parabolic of a quasi-split Lie algebra) and if $f$ is a principal nilpotent element (in the notation of [6]), then $\mathrm{of}^{f} \subseteq \bar{n}$ and hence $f$ is admissible for $p$. Thus, admissible nilpotents should be thought of as analogues of principal nilpotent elements.

We will later classify those parabolics (except, perhaps, for some parabolic subalgebras of $\mathrm{E}_{6}, \mathrm{E}_{7}$, and $\mathrm{E}_{8}$ ) for which there exist admissible nilpotents. In passing, we mention that all minimal parabolics have admissible nilpotents and that for any nilpotent element $f$ in $\Delta \ell(n ; \mathbb{C})$, there is a parabolic subalgebra for which $f$\\
is admissible.

Since $f \varepsilon d_{-1}$, it is clear that $[7, f]$ is an ad $x_{0}$-stable subspace of $p$. Since $x_{0}$ is semi-simple, we may select an ad $x_{0}$-stable subspace $\Delta \subseteq \rho$ such that


\begin{equation*}
p=[\pi, f] \oplus D \tag{1.2.2}
\end{equation*}


is a direct sum. Fix such an $J$. [Kostant]

It is interesting to consider the case in which $\mathrm{x}_{0}$ is the semi-simple element of a TDS (three-dimensional simple subalgebra) of $O y$. (See Kostant [6], Dynkin [3], or Ozeki-Wakimoto [ 10 ] for the material in this paragraph.) We may then choose $e \varepsilon d_{1}$ and $f \varepsilon d_{-1}$ such that $\left\{e, x_{0}, f\right\}$ spans the TDS. By considering the decomposition of $O \gamma$ under the adjoint action of the TDS, one finds that $\bar{\pi}^{e}=n^{f}=\{0\}$ and hence $f$ is admissible. In this case, we may choose $\Delta=0 \gamma^{e}$. However, such a choice is not particularly useful. Curiously, it turns out that for any admissible nilpotent, $f, g^{f}=\bar{p}^{f}$ is naturally dual to $\Lambda$.

The following theorem forms the basis for what follows. Theorem 1.2 [Kostant] The map


\begin{equation*}
N \times(f+\Delta) \rightarrow f+p \tag{1.2.3}
\end{equation*}


given by $(\mathrm{a}, \mathrm{x}) \rightarrow \mathrm{ax}$ is an isomorphism of affine varieties.

Proof. If a $\varepsilon N$ and $x \varepsilon \Delta \subseteq p$, it is clear that $\mathrm{a}(\mathrm{f}+\mathrm{x})$ ɛ f+p.

To show the injectivity of (1.2.3), we need only show that if $a(f+x)=f+y$ for $a \varepsilon N, x, y \varepsilon \Delta$, then $a$ is the identity. Since $p, n$, and $s$ are ad $x_{0}$-stable, we may find elements $x_{i}, y_{i}, z_{i} \varepsilon d_{i}$ such that $x=\sum_{i \geq 0} x_{i}, y=\sum_{i \geq 0} y_{i}$ and $a=\exp \left(\sum_{i \geq 1} z_{i}\right)$. If a is not the identity, let $I$ be the least integer such that $\mathrm{z}_{\mathrm{I}} \neq 0$. Then,


\begin{equation*}
a(f+x)=f+x+\left[z_{I}, f\right]+\left(\text { elements in } \sum_{i \geq I} d_{i}\right) . \tag{1.2.4}
\end{equation*}


Thus $y_{I-1}-x_{I-1}=\left[z_{I}, f\right]$. Since $\Delta$ is ad $x_{0}$-stable, $y_{I-1}-x_{I-1} \varepsilon \Delta$. Hence, $\left[z_{I}, f\right] \varepsilon \Delta \bigcap[\Omega, f]=\{0\}$. But this is a contradiction since $\pi^{f}=\{0\}$. Thus, a is the identity and (1.2.3) is an injection.

To show that (1.2.3) is onto, we take $y \varepsilon \rho$ and write $y=\sum_{i \geq 0} y_{i}$ where $y_{i} \varepsilon d_{i}$. Clearly, $(f+y)-f \varepsilon \sum_{i \geq 0} d_{i}$ so we may proceed inductively. Suppose that we have found elements $a_{(I)} \varepsilon N, x_{(I)} \varepsilon D$ such that\\
$w=(f+y)-a_{(I)}\left(f+x_{(I)}\right) \varepsilon \sum_{i \geq I \geq 0} d_{i}$. Let $w_{I}$ be the componert of $w$ in $d_{I}$. By (1.2.2), there exist unique elements $z_{I+1} \varepsilon d_{I+1}$ and $x_{I} \varepsilon d_{I} \cap D$ such that $w_{I}=\left[z_{I+1}, f\right]+x_{I} . \quad$ Set $a_{(I+1)}=a_{(I)} \exp \left(z_{I+1}\right)$, $x_{(I+1)}=x_{(I)}+x_{I}$. Then,

$$
\begin{aligned}
(f+y) & -a_{(I+1)}\left(f+x_{(I+1)}\right)= \\
& =(f+y)-a_{(I)}(f+x(I))-\left[z_{I+1}, f\right]-x_{I} \\
& +\left(\text { elements in } \sum_{i>I} d_{i}\right) \\
& =w-w_{I}+\left(\text { elements in } \sum_{i>I} d_{i}\right) \varepsilon \sum_{i>I} d_{i} .
\end{aligned}
$$

Thus, we may sind elements $a_{(I+1)} \varepsilon N, x_{(I+1)} \varepsilon \Delta$ so that $(f+y)-a_{(I+1)}(f+x(I+1)) \varepsilon \sum_{i \geq I+1} d_{i}$. However, since $d_{i}=\{0\}$ for large enough $i$, by induction we have elements $a \varepsilon N, x \varepsilon \Delta$ such that $a(f+x)=f+y$. Thus, (1.2.3) is onto.

Hence, (1.2.3) is a bijection. It is also a morphism of nonsingular affine varieties. It is therefore an isomorphism of affine varieties by the Zariski main theorem. Q.E.D.\\
1.3 Some definitions due to Kostant and Kazhdan

If $V$ is a finite-dimensional complex vector space, $S(V)$ will denote the symmetric algebra over $V$ and $S_{k}(V)$ will denote the homogeneous terms of degree $k$. If $V^{\prime}$ is the dual space to $V$, then $S\left(V^{\prime}\right)$ will be regarded as the algebra of polynomial functions on V .

Now, the adjoint action of $G$ on $\mathcal{J}$ and the contragradient action on $y^{\prime}$ extend to automorphisms of $\mathrm{s}(\mathrm{o})$ and $S\left(\sigma^{\prime}\right)$. Let $B$ denote the Killing form. Then there is a unique algebra, G-module isomorphism $S(g) \rightarrow S($ of $)$, $v \rightarrow v^{\prime}$ where if $x, y \varepsilon$ of, then $x^{\prime}(y)=B(x, y)$. For convenience we will write $u(x)$ for $u^{\prime}(x)$ where $u \in S(g)$ and $x \varepsilon g$ and thus avoid considering $s\left(g^{\prime}\right)$.

Since $\bar{p}$ is dual to $p$ under $B$, the above map gives an isomorphism $S(\bar{p}) \rightarrow S\left(\bar{p}^{\prime}\right)$ and thus we may avoid considering $S\left(p^{\prime}\right)$.

It will be convenient to introduce a new grading on $S(\bar{p})$ (due to Kazhdan) as follows. For any $x \in O f$, the extension of ad $x$ to a derivation of $S(O \gamma)$ will again be denoted by ad $x$. Since the restriction of $-a d x_{0}$ to $\bar{p}$ has non-negative integral eigenvalues, the same is true of its restriction to $S(\bar{p})$. Let $\left(S_{k}(\bar{p})\right)_{j}$ be the eigenspace in $S_{k}(\bar{p})$ for -ad $x_{0}$ corresponding\\
to the eigenvalue $j$ and put\\
(1.3.1)

$$
S_{(i)}(\bar{\rho})=\underset{j+k=i}{\oplus}\left(S_{k}(\bar{p})\right)_{j}
$$

Note that this is a direct sum. Clearly, $S_{(i)}(\bar{p}) S_{(j)}(\bar{p}) \subseteq S_{(i+j)}(\bar{p})$, so that the $S_{(i)}(\bar{p})$ give $S(\bar{p})$ the structure of a graded algebra. We refer to this grading as the $x_{0}$-grading. For convenience, we set $S_{i}(\bar{p})=S_{(i)}(\bar{p})=0$ if $i<0$.

Since $f+p$ is stable under the adjoint action of $N$, we may define an affine (but not linear) action of $N$ on $\varphi$ where if $a \varepsilon N$ and $x \varepsilon p$, then


\begin{equation*}
a \cdot x=a(f+x)-f . \tag{1.3.2}
\end{equation*}


By Theorem 1.2, the map


\begin{equation*}
N \times \Delta \rightarrow p, \quad(a, x) \rightarrow a \cdot x \tag{1.3.3}
\end{equation*}


is an isomorphism. But now, since $S(\bar{\rho})$ is the affine algebra of $\rho$, this action of $N$ induces an $N$-module structure on $S(\bar{p})$ where if $a \varepsilon N, x \varepsilon \rho$ and $u \varepsilon S(\bar{p})$ one has


\begin{equation*}
(a \cdot u)(x)=u\left(a^{-1} \cdot x\right) . \tag{1.3.4}
\end{equation*}


Of course, $N$ operates as a group of automorphisms of $S(\bar{\rho})$ and hence the space of $N$-invariants, $S(\bar{\rho})^{N}$, is a subalgebra of $S(\bar{\rho})$. Furthermore, by the isomorphism (1.3.3) we see that the map $\left.u \rightarrow u\right|_{D}$ is an isomorphism of $S(\bar{p})^{N}$ and $S\left(D^{\prime}\right)$.\\
1.4 The graded action of $\pi$ on $S(\bar{\rho})$.

The material in this section is due to [Kostant].\\
The action (1.3.4) of $N$ on $S(\bar{\rho})$ is locally finite (i.e. $N \cdot V$ spans a finite dimensional space for any $v \varepsilon S(\bar{p})$ ) since it arises from an affine action of $N$ on $p$. Infinitesimally, one then obtains a representation of $\Omega$ as a Lie algebra of derivations of $S(\bar{\rho})$ where if $z \varepsilon \lambda, v \varepsilon S(\bar{p})$ one defines $z \cdot v \varepsilon S(\bar{p})$ by


\begin{equation*}
z \cdot v=\left.\frac{d}{d t}(\exp t \text { ad } z \cdot v)\right|_{t=0} . \tag{1.4.1}
\end{equation*}


Remark 1.4. Again, by local finiteness one notes that $z \cdot v=0$ for all $z \varepsilon \pi$, if and only if $v \varepsilon S(\bar{p})^{N}$.

Now clearly the derivation defined by $z$ is determined as soon as one knows $z^{\circ} x$ for $x \varepsilon \bar{p}$. Let $\pi: \mathcal{\gamma} \rightarrow \bar{p}$ be\\
the projection map with kernel $\pi$. If $l$ is the identity in $S(\bar{\rho})$, then $S_{(0)}(\bar{\rho})=\mathbb{C} 1$. We now observe that the derivation carries $\bar{p}$ into $\bar{p} \oplus s_{(0)}(\bar{p})$.

Proposition 1.4.1. For any $x \in \bar{p}, z \varepsilon \pi$ one has


\begin{equation*}
z \cdot x=\pi[z, x]+B([z, x], f)] . \tag{1.4.2}
\end{equation*}


Proof. Let $z(t)=\exp t z \varepsilon N$ and let $y \varepsilon p$. Then, $(z \cdot x)(y)=\left.\frac{d}{d t} x(z(-t) \cdot y)\right|_{t=0}=\left.\frac{d}{d t} B(x, z(-t) \cdot y)\right|_{t=0}$. But $z(-t) \cdot y=z(-t)(f+y)-f$ so that one has $\left.\frac{d}{d t}(z(-t) \cdot y)\right|_{t=0}=-[z, f+y]$. Thus $(z \cdot x)(y)=B(x,-[z, f+y])$ and by the invariance of the Killing form one has $(z \cdot x)(y)=B([z, x], f+y)=(\pi[z, x])(y)+B([z, x], f)]$. Thus $z \cdot x=\pi[z, x]+B([z, x], f)]$. Q.E.D.

Now, let $\left\{x_{i}\right\}$ be a basis for $\Omega$. Then, if $v \in S(\bar{p})$, there exist unique elements $v_{0}, v_{i} \varepsilon S(\bar{p})$ such that if $\mathbf{x} \in \Omega$,


\begin{equation*}
\operatorname{ad} x(v)=v_{0}+\sum_{i} v_{i} x_{i} . \tag{1.4.3}
\end{equation*}


The uniqueness is clear since $x_{i} \notin S(\bar{p})$. The existence\\
follows from the fact that l) ad x is a derivation and 2) (1.4.3) holds if $v \in \bar{p}$.

Proposition 1.4.2. Let $\left\{\mathrm{x}_{\mathrm{i}}\right\}, \mathrm{v}, \mathrm{v}_{\mathrm{i}}$ be as above. Then


\begin{equation*}
x^{\cdot} v=v_{0}+\sum_{i} v_{i} B\left(f, x_{i}\right) . \tag{1.4.4}
\end{equation*}


Proof. If $v=y \varepsilon \bar{p}$, then (in the notation of (1.4.3)) $y_{0}=\pi[x, y]$ and the $y_{i}$ 's are constants if $i>0$. Now, $B([x, y], f)=B\left(\sum y_{i} x_{i}, f\right)=\sum y_{i} B\left(f, x_{i}\right)$. Thus, $\mathrm{x} \cdot \mathrm{y}=\mathrm{y}_{0}+\sum \mathrm{y}_{\mathrm{i}} \mathrm{B}\left(\mathrm{f}, \mathrm{x}_{\mathrm{i}}\right)$ by Proposition (1.4.1). This proves (1.4.4) for $v \varepsilon \bar{p}$. Next, let $D_{x}$ be the operator defined so that if $v \varepsilon S(\bar{\rho}), D_{x} v=v_{0}+\sum_{i} v_{i} B\left(f, x_{i}\right)$. Then $D_{X} v=X \cdot v$ for $v \varepsilon \bar{p}$. Since $x$ operates as a derivation of $S(\bar{p})$, to prove (1.4.4) we only need to show that $D_{X}$ is a derivation. Now, if $u, v \in S(\bar{p})$ we have that ad $x(u v)=\left(u_{0} v+u v_{0}\right)+\sum_{i}\left(u_{i} v+u v_{i}\right) x_{i}$. Thus,

$$
\begin{aligned}
D_{x}(u v) & =u_{0} v+u v_{0}+\sum_{i}\left(u_{i} v+u v_{i}\right) B\left(f, x_{i}\right) \\
& =\left(D_{x} u\right) v+u\left(D_{x} v\right)
\end{aligned}
$$

and $D_{X}$ is a derivation.\\
Q.E.D.

We next observe that the action of $\pi$ on $S(\bar{p})$ is graded in the sense described in the next proposition.

Proposition 1.4.3. Let $x \in d_{j}, j \geq 1$ and $v \in S_{(k)}(\bar{p})$. Then $x \cdot v \in S_{(k-j)}(\bar{p})$.

Proof. Clearly, we may take $x \in d_{j}, j \geq 1$ and $v \varepsilon\left(S_{p}(\bar{p})_{q}\right.$ where $p+q=k$. Since $S_{p}(\bar{p})$ is stable under the adjoint action and since - ad $x_{0}(x)=-j x$, it follows from (1.4.3) that $v_{0} \varepsilon\left(S_{p}(\bar{p})\right)_{q-3}$ and $v_{i} \varepsilon\left(S_{p-1}(\bar{p})\right)_{q-j+0\left(x_{i}\right)}$. However, $B\left(f, x_{i}\right)=0$ unless $0\left(x_{i}\right)=1$ and it clearly follows from (1.4.4) that $x \cdot v \varepsilon S_{(k-j)}(\bar{p})$. Q.E.D.

Proposition 1.4.4. $S(\bar{p})^{N}$ is an $x_{0}$-graded subalgebra of $S(\bar{p})$.

Proof. If $v \varepsilon S(\bar{\rho})^{N}, v=\sum_{i \geq 0} v_{i}$ where $v_{i} \varepsilon S_{(i)}(\bar{\rho})$. Then, for any $x \in d_{j} \prime^{\prime} \geq 1$ we have that $x \cdot v=\sum_{i \geq 0}\left(x \cdot v_{i}\right)$.

But, by Proposition 1.4.3 $x \cdot v_{i} \varepsilon S_{(i-j)}(\bar{\rho})$. Since $S(\bar{p})=\underset{i}{\oplus} S_{(i)}(\bar{p})$ is a direct sum, this implies that $\mathbf{x} \cdot \mathbf{v}_{\mathbf{i}}=0$ for each $\mathbf{i}$. Thus, $\mathbf{v}_{\mathbf{i}} \varepsilon S(\bar{\rho})^{\mathbf{N}}$ and hence $S(\bar{\rho})^{\mathrm{N}}$ is an $x_{0}$-graded subalgebra of $S(\bar{\rho})$.\\
Q.E.D.\\
1.5. The graded decomposition $S(\bar{\rho})=A \otimes S(\bar{\rho})^{N}$

The material in this section is due to [Kostant].\\
Now, N is a unipotent algebraic group. Let $\mathrm{A}(\mathrm{N})$ be the affine algebra of all regular functions on $N$. That is, $A(N)$ is the algebra of representative functions on $N$. Then $P$ operates as a group of automorphisms of $A(N)$ where if $\gamma \varepsilon A(N), n, a \varepsilon N$, and $m \varepsilon M$, then one defines $(\mathrm{nm}) \cdot \gamma \varepsilon \mathrm{A}(\mathrm{N})$ so that


\begin{equation*}
((n m) \cdot \gamma)(a)=\gamma\left(\operatorname{Ad}\left(m^{-1}\right)\left(n^{-1} a\right)\right) . \tag{1.5.1}
\end{equation*}


Remark 1.5. If one identifies $N$ with $P / M$ by coset projection, then the action (1.5.1) of $P$ on $N$ is natural.

Now, recalling (1.2.3), note that the affine algebra of $N \times \Delta$ is $A(N) \otimes S\left(\Delta^{\prime}\right)$. Since (1.2.3) is an isomorphism, one necessarily has an algebra isomorphism


\begin{equation*}
A(N) \otimes S\left(S^{\prime}\right) \rightarrow S(\bar{p}) \tag{1.5.2}
\end{equation*}


In particular, for any $\gamma \in \mathrm{A}(\mathrm{N})$, there exists an element $v_{\gamma} \varepsilon S(\bar{\rho})$ so that $v_{\gamma}(a \cdot x)=\gamma(a)$ for any a $\varepsilon N$ and $x \in D$. Furthermore, if $A \subseteq S(\bar{p})$ is the space of all $v_{\gamma}, \gamma \varepsilon A(N)$, then $A$ is a subalgebra of $S(\bar{\rho})$ and one\\
has an algebra isomorphism


\begin{equation*}
A(N) \rightarrow A, \gamma \rightarrow v_{\gamma} . \tag{1.5.3}
\end{equation*}


Now, A inherits an N-module structure from $S(\bar{p})$ by (1.3.4). It remains to relate the N -module structures on $A$ and $A(N)$.

Proposition 1.5.1. A is an N-submodule of $S(\bar{\rho})$ and the map (1.5.3) is an isomorphism of N -modules.

Proof. Let $\mathrm{a}, \mathrm{b} \varepsilon \mathrm{N}, \gamma \varepsilon \mathrm{A}(\mathrm{N})$ and $\mathrm{x} \varepsilon \Delta$. Then, $\left(a \cdot v_{\gamma}\right)(b \cdot x)=v_{\gamma}\left(a^{-1} \cdot(b \cdot x)\right)=v_{\gamma}\left(a^{-1}(b \cdot x+f)-f\right. =v_{\gamma}\left(\left(a^{-1} b\right)(f+x)-f\right)=v_{\gamma}\left(\left(a^{-1} b\right) \cdot x\right)=\gamma\left(a^{-1} b\right)=(a \cdot \gamma)(b)$. Hence


\begin{equation*}
a \cdot v_{\gamma}=v_{a \cdot \gamma} \cdot \tag{1.5.4}
\end{equation*}


Q.E.D.

We next show that $A$ is a graded subalgebra of $S(\bar{\rho})$.\\
Let $x_{0}(t)=\exp \left(t x_{0}\right)$ for $t \varepsilon R$. Now, let $r(t)$, $t \varepsilon R$, be the one-parameter group of linear transformations of $\rho$ defined by


\begin{equation*}
r(t) z=e^{t} A d\left(x_{0}(t)\right) z \tag{1.5.5}
\end{equation*}


for $z \varepsilon p$. One then defines an action $r(t)$ on $S(\bar{p})$ by $(r(t) \cdot v)(z)=v(r(-t) z)$.

Proposition 1.5.2. For any $k \in \mathbb{Z}_{+}$and $v \in S_{(k)}(\bar{p})$ one has


\begin{equation*}
r(t) \cdot v=e^{-k t} v . \tag{1.5.6}
\end{equation*}


Furthermore, a subspace $V \subseteq S(\bar{p})$ is graded with respect to the $\mathrm{x}_{0}$-gradation if and only if it is stable under the action of the one-parameter group $r(t)$.

Proof. Recalling the definition in §1.1 of the diagonals $d_{i}$ of $o \gamma$, one clearly has $d_{-j}=S_{(j+1)}(\bar{p}) \cap \bar{p}$ for $j \varepsilon \mathbb{Z}$. Now, if $x \varepsilon d_{-j}$ and $y \varepsilon \rho$, we have that $(r(t) \cdot x)(y)=B(x, r(-t) y)=B\left(x, e^{-t} A d\left(x_{0}(-t)\right) y\right)= e^{-t} B\left(\operatorname{Ad}\left(x_{0}(t)\right) x, y\right)=e^{-t} e^{-j t} B(x, y)=\left(e^{-(j+1) t} x\right)(y)$. Thus, $r(t) \cdot x=e^{-(j+1) t} x$. This proves (1.5.6) for $v \varepsilon-\bar{p}$. Since $\bar{p}$ generates $S(\bar{p})$ and $r(t)$ operates as an automorphism, this proves (1.5.6) for all $v \in S_{(k)}(\bar{\rho})$. The last statement clearly follows from (1.5.6) and the linear independence of the functions $e^{-k t}, k \varepsilon \mathbb{Z}$.\\
Q.E.D.

Lemma 1.5.1. Let $a \varepsilon N, x \varepsilon \rho$. Then for any $t \varepsilon \mathbb{R}$\\
one has


\begin{equation*}
r(t)(a \cdot x)=\left(\operatorname{Ad}\left(x_{0}(t)\right) a\right) \cdot r(t) x . \tag{1.5.7}
\end{equation*}


Proof. Since $f \varepsilon d_{-1}$, one has $x_{0}(t) f=e^{-t} f$ and hence $r(t) f=f$. Thus, $r(t)(a \cdot x)=(r(t) a(f+x))-f =e^{t} \operatorname{Ad}\left(x_{0}(t)\right)(a(f+x))-f=\left(\operatorname{Ad}\left(x_{0}(t)\right) a\right)(r(t)(f+x))-f$. Thus, $x(t)(a \cdot x)=\left(x_{0}(t) \cdot a\right)(f+r(t) x)-f=$ (Ad $\left.\left(x_{0}(t)\right) a\right) \cdot r(t) x$.\\
Q.E.D.

Theorem 1.5. Let $A$ be the subalgebra of $S(\bar{\rho})$ defined by (1.5.3) so that $A$ is isomorphic to the affine algebra $A(N)$ of the unipotent group $N$. Then $A$ is graded with respect to the $x_{0}$-grading. In fact, for any $\gamma \varepsilon A(N)$ one has


\begin{equation*}
r(t) \cdot v_{\gamma}=v_{x_{0}}(t) \cdot \gamma \tag{1.5.8}
\end{equation*}


where $\left(x_{0}(t), \gamma\right)(a)=\gamma\left(\operatorname{Ad}\left(x_{0}(-t)\right) a\right)$ for any a $\varepsilon N$.\\
Moreover, if tensor product maps to multiplication, one has an isomorphism


\begin{equation*}
A \otimes S(\bar{p})^{N} \rightarrow S(\bar{p}) \tag{1.5.9}
\end{equation*}


of graded algebras with respect to the $x_{0}$-grading. In particular, the action of $N$ on $S(\bar{p})$ defined by (1.3.4) in effect reduces to left translation on $A(N)$ in that one has\\
(1.5.10) $\quad a \cdot\left(v_{\gamma} w\right)=v_{a \cdot \gamma} w$\\
for any a $\varepsilon N, \gamma \varepsilon A(N)$ and $w \varepsilon S(\bar{\rho})^{N}$.\\
Proof. Note that the subspace $D \subseteq \rho$ is ad $x_{0}$-stable and hence is $r(t)$-stable. Thus, let $a \varepsilon N, x \varepsilon \Delta$. Then, $\left(r(t) \cdot v_{\gamma}\right)(a \cdot x)=v_{\gamma}(r(-t)(a \cdot x))=v_{\gamma}\left(\left(\operatorname{Ad}\left(x_{0}(-t)\right) a\right) \cdot r(-t) x\right) =\gamma\left(\operatorname{Ad}\left(x_{0}(-t)\right) a\right)=\left(x_{0}(t) \cdot \gamma\right)(a)=v_{x_{0}}(t) \cdot \gamma(a \cdot x)$. Thus, $r(t) \cdot v_{\gamma}=v_{x_{0}}(t) \cdot \gamma$ which proves the first two statements of the theorem.

Now, $S(\bar{p})^{N}$ is an $x_{0}$-graded subalgebra by Proposition 1.4.4. Furthermore, $S\left(\Delta^{\prime}\right)$ clearly maps into $S(\bar{p})^{\mathrm{N}}$ by the isomorphism (1.5.2). However, as we noted after (1.3.4), $S(\bar{p})^{N}$ is isomorphic to $S\left(S^{\prime}\right)$ and hence (1.5.2) establishes (1.5.9) as an isomorphism. The final relation (1.5.10) follows from Proposition 1.5 .1 and the fact that $N$ operates as a group of automorphisms of $S(\bar{p})$.\\
Q.E.D.\\
1.6. The structure of A

By Theorem 1.5, $A$ is an $x_{0}$-stable subalgebra of $S(\bar{p})$ and hence $A$ is graded by $A_{(i)}=A \cap S_{(i)}(\bar{p})$. We may now describe the graded structure of $A$.

Theorem 1.6. [Kostant] A is isomorphic to a polynomial algebra on $\operatorname{dim} \pi$ generators


\begin{equation*}
A \cong C\left[u_{1}, \ldots, u_{n}\right] \tag{1.6.1}
\end{equation*}


Moreover, the $u_{i}$ 's may be chosen to be $x_{0}$-stable so that if $r_{j}$ is the number of $u_{i}$ 's such that $u_{i} \varepsilon A_{(j)}$, then $r_{j}=\operatorname{dim} d_{j}$. Thus, the Poincaré polynomial for $A$ is given by\\
(1.6.2) $\quad \sum_{k=0}^{\infty}\left(\operatorname{dim} A_{(k)}\right) t^{k}=\prod_{j \geq 1}\left(\frac{1}{1-t^{j}}\right)^{\operatorname{dim} d_{j}}$.

Proof. Recall the action of $P$ on $N$ given by (1.5.1). For any $x \in \rho$, let $\xi_{x}$ be the vector field on $N$ given by $\left(\xi_{x}\right)(a)=\left.\frac{d}{d t} \gamma(\exp (-t x) \cdot a)\right|_{t=0}$ for a $\varepsilon N$ and $\gamma \varepsilon A(N)$.

If $U(n) \subseteq U(\rho)$ are respectively the universal enveloping algebras of $\pi$ and $p$, the representation $x \rightarrow \xi_{x}$ of on $A(N)$ extends to a representation $u \rightarrow \xi_{u}$ of $U(p)$\\
on $A(N)$. Since $N$ operates as left translations on $N$, the subalgebra $U(T)$ maps bijectively onto the set of all right-invariant differential operators on $N$. Thus, one obtains a non-singular pairing of $A(N)$ and $U(\Omega)$ by putting


\begin{equation*}
\langle\gamma, u\rangle=\left(\xi_{u} \gamma\right)(1) \tag{1.6.3}
\end{equation*}


where 1 is the identity of $N, u \in U(\pi)$ and $\gamma \varepsilon A(N)$.\\
Since $x_{0} \varepsilon m$ normalizes $U(\pi)$, if $j \varepsilon \mathbb{Z}_{+}$let $U(T)_{j}=\left\{u \in U(\pi) \mid\left[x_{0}, u\right]=j u\right\}$. Since for $x \varepsilon d_{i}$, [ $\mathrm{x}_{0}, \mathrm{x}$ ] $=\mathrm{ix}$, it follows from the Birkhoff-Witt theorem that the $U(\pi)_{j}$ define the structure of a graded algebra on $U(\pi)$ whose Poincaré polynomial is given by the right side of (1.6.2). But now the vector field $\xi_{x_{0}}$ vanishes at 1 and hence $0=\left(\xi_{x_{0}} \xi_{u} \gamma\right)(1)=\left(\xi_{\left[x_{0}, u\right]} \gamma\right)(1)+ \left(\xi_{u} \xi_{x_{0}} \gamma\right)(1)$. Thus, for any $\gamma \varepsilon A(N), u \varepsilon U(\pi)$


\begin{equation*}
\left\langle\xi_{-x_{0}} \gamma, u\right\rangle=\left\langle\gamma,\left[x_{0}, u\right]\right\rangle . \tag{1.6.4}
\end{equation*}


Now, for any $j \varepsilon \mathbb{Z}_{+}$let $A(N)_{j}=\left\{\gamma \varepsilon A(N) \mid \xi_{-x_{0}} \gamma=j \gamma\right\}$. We assert that the $A(N)_{j}$ define the structure of a graded algebra on $A(N)$ whose Poincaré series is also\\
given by the right hand side of (1.6.2). If $U(\pi)$ ' is the dual space to $U(\pi)$, let $\rho: A(N) \rightarrow U(\pi)^{\prime}$ be the injective map defined by the pairing (1.6.3). The assertion then easily follows from (1.6.4) as soon as one notes that the image of $\rho$ is the set of all linear functionals on $U(\pi)$ which vanish on all but a finite number of the $U(\pi)_{j}$. But since $N$ is unipotent, one knows that the image of $\rho$ is the set of all linear functionais which vanish on a power of the augmentation ideal $\pi U(\pi)$. The observation, and hence the assertion, is then immediate.

If $x_{0}(t)=\exp \left(t\right.$ ad $\left.x_{0}\right) \varepsilon M$ as in Theorem 1.5, one has $x_{0}(t) \cdot \gamma=e^{-t k} \gamma$ for $\gamma \varepsilon A(N)_{k}$. This is clear since $\xi_{x_{0}} \gamma=-k \gamma$. But now recalling the isomorphism $A(N) \rightarrow A \subseteq S(\bar{p}), \gamma \rightarrow v_{\gamma^{\prime}}$ one has $r(t) \cdot v_{\gamma}=v_{x_{0}}(t) \cdot \gamma$ by (1.5.8). Thus, if $\gamma \varepsilon A(N)_{k}$ one has $r(t) \cdot v_{\gamma}=e^{-t k} v_{\gamma}$. But then $v_{\gamma} \varepsilon A_{(k)}$ by Proposition 1.5.2. Thus, the map $\gamma \rightarrow v_{\gamma}$ induces an isomorphism


\begin{equation*}
A(N)_{k} \rightarrow A_{(k)} \tag{1.6.5}
\end{equation*}


for any $k \in \mathbb{Z}_{+}$. Thus, $\operatorname{dim} A_{(k)}=\operatorname{dim} A(N)_{k}$ proving (1.6.2). The other statements of the theorem also\\
follow from (1.6.5).\\
Q.E.D.\\
1.7. The structure of $S(\bar{p})^{\mathrm{N}}$

Analogous to our definition of the $x_{0}$-grading in $S(\bar{p})$ we may define a grading in $S(p)$ as follows. As before, if $j, k \varepsilon \mathbb{Z}_{k}$ and $\left(S_{j}(p)\right)_{k}$ is the eigenspace in $S_{j}(p)$ corresponding to the eigenvalue $k$ of $\operatorname{ad} x_{0}$, then we set\\
(1.7.1)

$$
S_{(i)}(p)=\underset{j+k=i}{\oplus}\left(S_{j}(p)\right)_{k}
$$

and note that (1.7.1) is a direct sum. As before, the $S_{(i)}(p)$ define a grading on $S(p)$. Now, since $p$ and $\bar{p}$ are nonsingularly paired under the Killing form, we may extend this pairing in a natural way to a nonsingular pairing of $S(p)$ and $S(\bar{p})$. Under this pairing, $\left(S_{j}(p)\right)_{k}$ is nonsingularly paired to $\left(S_{j},(\bar{p})\right)_{k^{\prime}}$ if $j=j^{\prime}$ and $k=k^{\prime}$ and is orthogonal otherwise. Thus, $S_{(i)}(p)$ is nonsingularly paired to $S_{\left(i^{\prime}\right)}(\bar{p})$ if and only if $i=i$ '.

We will use the pairing of $S(p)$ and $S(\bar{p})$ defined above together with the decomposition $S(\bar{p})=A \otimes S(\bar{p})^{N}$\\
to determine the structure of $S(\bar{\rho})^{N}$. First, however, we need the following lemma.

Recall that $S_{(i)}(\bar{p})=\underset{j+k=i}{\oplus}\left(S_{j}(\bar{p})\right)_{k}$ (1.3.1) is a direct sum. Thus, if $u \in S_{(i)}(\bar{p})$, let $\pi_{1} u$ denote the component of $u$ in $\left(S_{1}(\bar{p})\right)_{i-1}$ under (1.3.1). Clearly, $\pi_{i} u \in d_{-i+1} \subseteq \bar{p}$.

Lemma 1.7. If $u \varepsilon S_{(i)}(\bar{p})^{N}$, then $\pi_{1} u \varepsilon \bar{p}^{f} \cap d_{-i+1}$ where $\bar{p}^{f}$ is the centralizer of $f$ in $\bar{p}$. (Note that since $f \varepsilon d_{-1}, p^{-f}$ is clearly ad $x_{0}$-stable.)

Proof. If $u \in S_{(i)}(\bar{p})^{N}$, then $u=\pi_{1} u+u^{\prime}$ where $u^{\prime} \varepsilon \underset{\substack{j+k=i \\ j>1}}{\oplus}\left(S_{j}(\bar{\rho})\right)_{k}$. It is then clear that $u^{\prime}$ may be\\
written as a sum of elements of the form $\mathrm{v}_{\mathrm{p}} \mathrm{v}_{\mathrm{q}}$ where $v_{p} \varepsilon S_{(p)}(\bar{p})$ and $v_{q} \varepsilon S_{(q)}(\bar{p}), p+q=i, l \leq p, q<i$. By Proposition 1.4.3 we have that if $x \varepsilon d_{i}$, then $x \cdot\left(v_{p} v_{q}\right)=\left(x \cdot v_{p}\right) v_{q}+v_{p}\left(x \cdot v_{q}\right)$. However, $x \cdot v_{p} \varepsilon S_{(p-i)}(\bar{p}), x \cdot v_{q} \varepsilon S_{(q-1)}(\bar{p})$ implies that $\mathrm{x} \cdot \mathrm{v}_{\mathrm{p}}=\mathrm{x} \cdot \mathrm{v}_{\mathrm{q}}=0$. Thus, $\mathrm{x} \cdot\left(\mathrm{v}_{\mathrm{p}} \mathrm{v}_{\mathrm{q}}\right)=0$ for all $\mathrm{x} \varepsilon \mathrm{d}_{\mathrm{i}}$ and hence $\mathrm{x} \cdot \mathrm{u}^{\prime}=0$. Thus, $\mathrm{x}^{\circ} \mathrm{u}=\mathrm{x}^{\circ}\left(\pi_{1} \mathrm{u}\right)$ for all $\mathrm{x} \varepsilon \mathrm{d}_{\mathrm{i}}$. By Proposition 1.4.1 we have that $x \cdot\left(\pi_{1} u\right)=B\left(f,\left[x, \pi_{1} u\right]\right)$\\
since $\left[x, \pi_{1} u\right] \varepsilon d_{i} \subseteq \pi$. By the invariance of the Killing form we have that if $x \in d_{i}$ and $u \in S_{(i)}(\bar{\rho})^{N}$, $0=\mathrm{x} \cdot \mathrm{u}=\mathrm{x} \cdot\left(\pi_{1} \mathrm{u}\right)=\mathrm{B}\left(\mathrm{f},\left[\mathrm{x}, \pi_{1} \mathrm{u}\right]\right)=-\mathrm{B}\left(\left[\mathrm{f}, \pi_{1} \mathrm{u}\right], \mathrm{x}\right)$. Now, $\left[f, \pi_{1} u\right] \varepsilon d_{-i}$ and since $d_{-i}$ and $d_{i}$ are nonsingularly paired under the Killing form, it follows that $\pi_{1} u \in p^{-f}$. Q.E.D.

We are now ready to prove the following theorem.

Theorem 1.7. $S(\bar{p})^{N}$ is isomorphic to a polynomial algebra on dim $m$ generators


\begin{equation*}
S(\bar{p})^{N} \cong\left[I_{1}, \ldots, I_{\operatorname{dim} m}\right] \tag{1.7.2}
\end{equation*}


where the generators may be chosen to be consistent with the $x_{0}$-gradation. Further, if $m_{i}$ is the number of generators occurring in $S_{(i)}(\bar{\rho})^{N}$, then $m_{i}=\operatorname{dim}\left(\Delta \cap d_{i-1}\right) =\operatorname{dim} d_{i-1}-\operatorname{dim} d_{i}$. Thus, the Poincaré polynomial of $S(\bar{p})^{\mathrm{N}}$ is given by\\
(1.7.3) $\quad \sum_{k=0}^{\infty}\left(\operatorname{dim} S_{(k)}(\bar{p})^{N}\right) t^{k}=\prod_{i \in \mathbb{Z}_{+}}\left(\frac{1}{1-t^{i}}\right)^{m}$.

Proof. To prove the theorem, it is only necessary to verify (1.7.3). For $k=0$, we have that $S_{(0)}(\bar{p})^{N}=\mathbb{C} 1$\\
and this is trivial. Suppose we have shown that the coefficients $\operatorname{dim} S_{(j)}(\bar{p})^{N}$ satisfy (1.7.3) for all $j<k$.

Now, $S_{(k)}(\bar{p})$ is nonsingularly paired to $S_{(k)}(p)$. We may write $S_{(k)}(\rho)=d_{k-1} \oplus S_{(k)}^{\prime}(\rho)$ where $S_{(k)}^{\prime}(p)=\underset{\substack{i+j=k \\ i>1}}{\oplus}\left(S_{i}(p)\right)_{j}$ is the subspace of elements in $S_{(k)}(p)$ spanned by products of lower order elements. Similarly, let $A_{(k)}^{\prime}$ and $S_{(k)}^{\prime}(\bar{\rho})^{N}$ denote the subspaces of $A_{(k)}$ and $S_{(k)}(\bar{\rho})^{N}$ spanned by products of lower order elements from $A$ and $S(\bar{p})^{N}$ respectively. Let $A_{(k)}^{*}$ and $S_{(k)}^{*}(\bar{\rho})^{N}$ be complementary subspaces to $A_{(k)}^{\prime}$ and $S_{(k)}^{\prime}(\bar{\rho})^{N}$ in $A_{(k)}$ and $S_{(k)}(\bar{\rho})^{N}$ respectively. Let $S_{(k)}^{\prime}(\bar{p})=A_{(k)}^{\prime} \oplus S_{(k)}^{\prime}(\bar{p})^{N} \oplus \underset{\substack{i+j=k \\ j, i \geq 1}}{\oplus}\left(A_{(i)} \otimes S_{(j)}(\bar{p})^{N}\right)$ so under the graded isomorphism $S(\bar{p})=A \otimes S(\bar{p})^{N}, S_{(k)}^{\prime}(\bar{p})$ represents those terms in $S(\overline{-p})$ which are spanned by products of lower order terms under this isomorphism.

Since $S_{(i)}(\bar{p})$ is nonsingularly paired to $S_{(i)}(p)$ for all $i$ it follows that $S_{(k)}^{\prime}(\bar{p})$ is nonsingularly paired to $S_{(k)}^{\prime}(p)$. Thus, we have that $A_{(k)}^{*} \oplus S_{(k)}^{*}(\bar{p})^{N}$ is nonsingularly paired to $d_{k-1}$. Now, if. $v \varepsilon S_{(k)}^{*}(\bar{p})^{N}$, then $\pi_{1} v \neq 0$ since otherwise $v$ would be orthogonal to $\mathrm{d}_{\mathrm{k}-1}$. On the other hand, $\pi_{1} \mathrm{v} \varepsilon p^{-f} \cap \mathrm{~d}_{-\mathrm{k}+1}$ by

Lemma 1.7. However, $d_{k-1}=\left[d_{k}, f\right] \oplus\left(\Delta \cap d_{k-1}\right)$ by (1.2.2). If $x \varepsilon d_{k}, v \varepsilon S_{(k)}^{*}(\bar{p})^{N}$, then $\langle v,[x, f]\rangle=\left\langle\pi_{1} v,[x, f]\right\rangle =B\left(\pi_{1} v,[x, f]\right)=B\left(\left[f, \pi_{1} v\right], x\right)=0$. It follows that $S_{(k)}^{*}(\bar{p})^{N}$ is orthogonal to $\left[d_{k}, f\right]$. Now, by Theorem 1.6 and since $A_{(k)}^{*}$ consists of those elements not spanned by products of lower order elements, it follows that ${ }^{A^{*}}(k)$ is nonsingularly paired to $\left[d_{k}, f\right]$. Thus, $S_{(k)}^{*}(\bar{p})^{N}$ is nonsingularly paired to $\Delta \bigwedge \mathrm{d}_{\mathrm{k}-1}$.

Going back to (1.7.3) we see that $\operatorname{dim} S_{(k)}(\bar{p})^{N}= \operatorname{dim} S_{(k)}^{\prime}(\bar{p})^{N}+\operatorname{dim} S_{(k)}^{*}(\bar{p})^{N}$ is correctly given by (1.7.3) if and only if $m_{k}=\operatorname{dim}\left(\Delta \cap d_{k-1}\right)$. Q.E.D.

This theorem has a number of important consequences. Note that the map $u \rightarrow \pi_{1} u$ of $S(\bar{p}) \rightarrow \bar{p}$ is the same as the operation of taking differentials.

Corollary 1.7.1. Let $v_{i}, i=1, \ldots, \operatorname{dim} \eta$ be a set of generators for $A$ as in Theorem l.6. Then $\pi_{1} v_{i}, i=1, \ldots, \operatorname{dim} \eta$ defines a coordinate system on $[\pi, f]$.

Clearly, the generators of $A$ appearing in $A_{(k)}$ form a basis for a complementary subspace $A_{(k)}^{*}$ as above. However, $A_{(k)}^{*}$ is nonsingularly paired to $\left[d_{k}, f\right]$. Thus, $\pi_{1} A_{(k)}^{*}$ is a coordinate system on $\left[d_{k}, f\right]$.\\
Q.E.D.

Corollary 1.7.2. Let $I_{1}, \ldots, I_{\operatorname{dim}} m$ be a basis for $S(\bar{p})^{N}$ as in Theorem 1.7. Then $\pi_{1} I_{i}, i=1, \ldots, \operatorname{dim} m$ forms a coordinate system on $\Delta$.

Proof. This is clear from the proof of Theorem 1.7.\\
Q.E.D.

Corollary 1.7.3. If $f$ is admissible, then $\operatorname{dim} \bar{p}^{f}=\operatorname{dim} o^{f}=\operatorname{dim} m=\operatorname{dim} p-\operatorname{dim} n=\operatorname{dim} \Delta$.

Proof. Since $f$ is admissible, $\pi^{f}=\{0\}$ and hence $\gamma^{f}=\bar{p}^{f}$. By the nonsingularity of the Killing form, $y \varepsilon p^{-f}$ if and only if $B(x,[f, y])=0$ for all $x \in \eta$. However, since $[\bar{p}, f] \subseteq \bar{\pi}$, we need only consider $x \varepsilon \pi$. By the invariance of the Killing form, it follows that $y \in p^{-f} \Leftrightarrow B([f, x], y)=0$ for all $x \in \pi$. Thus, $\bar{p}^{-f}$ is the orthocomplement to $[\pi, f]$ in $\bar{p}$ and since $r^{f}=\{0\}$ the corollary follows.\\
Q.E.D.\\
1.8. The restriction of $S(\bar{p})^{N}$ to $S(m)$.

Recall that for any $i, S_{(i)}(\bar{p})=\underset{j+k=i}{\oplus}\left(S_{j}(\bar{p})\right)_{k}$ is a\\
direct sum. For any $u \in S_{(i)}(\bar{p})$, let $\pi_{i} u$ be the component of $u$ in $\left(S_{i}(\bar{p})\right)_{0}=S_{i}(m)$. Clearly, if $u \varepsilon S_{(i)}(\bar{p})$ and $v \varepsilon S_{(j)}(\bar{p}), \pi_{i+j}(u v)=\pi_{i}(u) \pi_{j}(v)$ so that the $\pi_{i}$ 's\\
define a graded map $\pi: S(\bar{p}) \rightarrow S(\pi)$ where $S(\bar{p})$ has the $x_{0}$-gradation and $S(\pi)$ has the standard gradation. Note that this map is in some sense "opposite" to the map $\pi_{1}$ defined in §1.7.

We now show that under this map $S(\bar{p})^{\mathrm{N}}$ may be considered as a submodule of $S(m)$.

Theorem 1.8. The map $S(\bar{p}) \rightarrow S(m)$ defines a graded isomorphism of $S(\bar{p})^{N}$ with a subalgebra of $S(m)$.

Proof. Clearly, we need only show that $S(\bar{\rho})^{N} \rightarrow S(M)$ is injective. Thus, we need only show that if $u \varepsilon S(\bar{p})^{N}$, then there is an $x \in M$ such that $u(x) \neq 0$. Now, $u \varepsilon S(\bar{p})^{N}$ if and only if $u(a \cdot x)=u(x)$ for all a $\varepsilon N$, $x \in p$. Thus, it suffices to show that $N \cdot M$ contains an open subset of $p$. Now, if $x \varepsilon m=d_{0}$ and $y \varepsilon d_{i}, i \geq 1$, then $[x, y] \varepsilon d_{i}$ and $[f, y] \varepsilon d_{i-1}$. From this, it quickly follows that $\lambda^{f+x}=\{0\}$ for all $x \in \mathcal{M}$. (In fact, it is clear that $\eta^{f+x}=\{0\}$ for all $x \in \rho$.) If $y \varepsilon \eta$, then $y \rightarrow[y,(f+x)]$ represents the differential of the map $N \cdot m \rightarrow p$ at the point $x \varepsilon \eta n$. However since $n^{f+x}=\{0\}$, the differential has rank $\operatorname{dim} n$ at each $x \varepsilon M$. Thus, $N \cdot M$ contains an open subset of $p$ containing $m$.\\
Q.E.D.

\subsection{Comments}
The prime motivation of this paper was the desire to prove the results in Kostant's paper [ 9 ] without reference to $S(G)^{G}$ and $\gamma$. The proof cf Theorem 2.1 was a major step ir this direction and allowed a signoficant generalization. It is interesting to note that we used very little information about $\pi$ and $\rho$ to prove this theorem. In particular, taking any $f \varepsilon d_{-1}$ and an $x_{0}$-stable $r \subseteq \pi$ such that $r^{f}=\{0\}$, Theorem 1.2 holds if we replace $N$ by $R=\exp (r)$. (On the other hand, Theorem 1.2 holds even if $o f$ is real.) Then the results of §1.3-1.5 hold if we take (1.5.1) to hold for $r \varepsilon R$ only (so as to keep $A(R)$ stable). Since $r$ is ad $x_{0}$-stable, $R$ is Ad exp ( $t x_{0}$ )-stable and we can define the necessary gradations. Thus, we have (as in 1.5.9) a graded isomorphism $A \otimes S(\bar{p})^{R} \rightarrow S(\bar{p})$ where $A$ now is the isomorph of $A(R)$. Also, Theorems 1.6 and 1.7 generalize whereas Theorem 1.8 probably does not. I have not written the exposition in this generality since the condition of admissible $\pounds$ seems to be precisely what is needed for representation theory.

\section{Chapter 2}
\section{A decomposition theorem for $U(\bar{p})$.}
As in Chapter 1, much of the material in this chapter comes from material in Chapter 2 of Kostant [9]. Material thus obtained will be denoted by [Kostant].\\
2.1. The $\mathrm{x}_{0}$-filtration of $U(\bar{p})$ and $\operatorname{Gr} U(\bar{p})$.

The material in this section is due to [Kostant] and Kazhdan.

If $\Omega \subseteq O \gamma$ is any Lie subalgebra, then $U(\Omega)$ will denote the universal enveloping algebra of $\Omega$. Of course, we will always regard $U(\sigma) \subseteq U(\sigma)$. We set $U=U(\sigma)$. For any $k \in \mathbb{H}_{+}$, let $U_{k}(\Omega) \subseteq U(\Omega)$ denote the subspace generated by all terms of the form $x_{1} \cdot \ldots \cdot x_{i}$ where $i \leq k$ and $x_{j} \in \Omega$. We refer to $\ldots \subseteq U_{k-1}(\Omega) \subseteq U_{k}(\Omega) \subseteq \ldots$ as the standard filtration of $U(\Omega)$. For convenience, we set $U_{k}(\pi)=\{0\}$ if $k<0$.

Corresponding to the $x_{0}$-gradation in $S(\bar{p})$, there is an analogous filtration in $U(\bar{\rho})$. As before, we note that if ad $x_{0}$ denotes the adjoint action of $x_{0}$ on $U(\bar{p})$, then $-a d x_{0}$ acts semi-simply with non-negative integral eigenvalues on $U(\bar{\rho})$. Clearly $U_{k}(\bar{\rho})$ is stable under

\begin{itemize}
  \item ad $x_{0}$. For $j \varepsilon \mathbb{Z}_{+}$, let $\left(U_{k}(\bar{p})\right)_{j}$ be the eigenspace for the eigenvalue $j$ of - ad $x_{0}$ restricted to $U_{k}(\bar{p})$ and put\\
(2.1.1)
\end{itemize}

$$
U_{(i)}(\bar{p})=\underset{j+k=i}{\oplus}\left(U_{k}(\bar{p})\right)_{j} .
$$

Clearly, $U_{(i)}(\bar{p}) \subseteq U_{(i+1)}(\bar{p})$ and $U_{(i)}(\bar{p}) U_{(j)}(\bar{p}) \subseteq U_{(i+j)}(\bar{p})$. Also, the union of the $U_{(i)}(\bar{p})$ equals $U_{(\bar{p})}$ so they define a filtration of $U(\bar{p})$. We refer to it as the $x_{0}$-filtration of $U(\bar{p})$.

Again, if $\Omega \subseteq O \gamma$ is any Lie algebra, then by the Birkhoff-Witt theory we have a series of exact sequences, $k=1,2, \ldots$\\
(2.1.2) $\quad 0 \rightarrow U_{k-1}(\sigma) \xrightarrow{\text { inj }} U_{k}(\sigma) \xrightarrow{\tau_{k} \mid \sigma} s_{k}(\sigma) \rightarrow 0$\\
such that $\tau_{1} l_{\Omega}$ is the identity map and if $u \varepsilon U_{i}(\Omega)$, $v \varepsilon U_{j}(a)$ then\\
(2.1.3)

$$
\tau_{i+j}(u v)=\tau_{i}(u) \tau_{j}(v) .
$$

Since $\tau_{k}$ commutes with the adjoint action of $\sigma$, we have the exact sequences\\
(2.1.4) $0 \rightarrow\left(U_{k-1}(\bar{\rho})\right)_{j} \xrightarrow{i n j}\left(U_{k}(\bar{\rho})\right)_{j} \xrightarrow{\tau_{k}}\left(S_{k}(\bar{\rho})\right)_{j} \rightarrow 0$.

Since the right hand side of (2.1.1) is clearly a direct sum, given $u \varepsilon U_{(i)}(\bar{p})$, we may write $u$ uniquely as $u=\sum_{j=0}^{i} u_{j}$ where $u_{j} \varepsilon\left(U_{i-j}(\bar{p})\right)_{j}$. But then we have a well defined map

$$
(2.1 .5) \quad \tau_{(i)}: U_{(i)}(\bar{p}) \rightarrow S_{(i)}(\bar{p})
$$

where $\tau_{(i)}(u)=\sum_{j=0}^{i} \tau_{i-j} u_{j}$.\\
Proposition 2.1.1. For any $i \varepsilon \mathbb{Z}_{+}$one has an exact sequence\\
(2.1.6) $\quad 0 \rightarrow U_{(i-1)}(\bar{p}) \xrightarrow{i n j} U_{(i)}(\bar{p}) \xrightarrow{\tau(i)} S_{(i)}(\bar{p}) \rightarrow 0$.

Furthermore, if $u \varepsilon U_{(i)}(\bar{p})$ and $v \varepsilon U_{(j)}(\bar{p})$, then\\
(2.1.7) $\tau_{(i+j)}(u v)=\tau_{(i)}(u) \tau_{(j)}(v)$.

Proof. Since $S_{(i)}(\bar{p})$ is a direct sum of - ad $x_{0}$ eigenspaces as is $U_{(i)}(\bar{p})$, (2.1.6) is an exact\\
sequence by (2.1.4).

Now, let $u_{p} \varepsilon\left(U_{i-p}(\bar{p})\right)_{p}, v_{q} \varepsilon\left(U_{j-q}(\bar{\rho})\right)_{q}$ be such that $u=\Sigma u_{p}, v=\Sigma u_{q}$. Then clearly $u v=\sum_{r}(u v)_{r}$ where $(u v)_{r} \varepsilon\left(U_{(i+j)-r}(\bar{p})\right)_{r}$ and $(u v)_{r}=\sum_{p+q=r} u_{p} v_{q}$. Thus, $\tau_{(i+j)}(u v)=\sum_{r} \tau_{i+j-r}\left((u v)_{r}\right) . \quad$ But by (2.1.3), $\tau_{i+j-r}\left((u v)_{r}\right)=\sum_{p+q=r} \tau_{i-p}\left(u_{p}\right) \tau_{j-q}\left(v_{q}\right)$. By summing over $r$ one clearly has (2.1.7).\\
Q.E.D.

Let $\operatorname{Gr} \|(\bar{p})$ be the graded algebra associated to the $\mathrm{x}_{0}$-filtration. Thus,\\
(2.1.8) $\quad \operatorname{Gr} U(\bar{\rho})={\underset{i=0}{\infty}}_{\operatorname{m}}^{\operatorname{Gr}}(i) U(\bar{\rho})$\\
where\\
(2.1.9) $\quad \operatorname{Gr}_{(i)} U(\bar{p})=U_{(i)}(\bar{p}) U_{(i-1)}(\bar{p})$.

Since one has the commutation relations $\left[\left(U_{p}(\bar{p})\right)_{q},\left(U_{r}(\bar{p})\right)_{s}\right] \subseteq\left(U_{p+r-1}(\bar{p})\right)_{q+s}$ it follows that\\
(2.1.10) $\quad\left[U_{(i)}(\bar{p}), U_{(j)}(\bar{p})\right] \subseteq U_{(i+j-1)}(\bar{p})$\\
and hence $\operatorname{Gr} U(\bar{\rho})$ is a commutative algebra.\\
Proposition 2.1.2. Let\\
(2.1.11) $\quad{ }^{\tau} \mathrm{x}_{0}: \operatorname{Gr} U(\bar{\rho}) \rightarrow S(\bar{\rho})$\\
be the map defined so that its restriction to $\operatorname{Gr}_{(i)} U_{(\bar{\rho})}$ is the map\\
(2.1.12) $\quad U_{(i)}(\bar{p}) / U_{(i-1)}(\bar{p}) \rightarrow s_{(i)}(\bar{p})$\\
induced by ${ }^{\tau}(\mathrm{i})$. Then $\tau_{\mathrm{x}_{0}}$ is an isomorphism of graded commutative algebras.

Proof. This is obvious from Proposition 2.1.1. Q.E.D.\\
2.2. The $n$-reduced action of $n$ on $U(\bar{p})$.

The material in this section is due to [Kostant].\\
It is useful to note that there is a natural bijection between $d_{-1}$ and the characters (one-dimensional representations) of $\Omega$ so that if $f \varepsilon d_{-1}$, then $f$ corresponds to $n: n \rightarrow \mathbb{C}$ if\\
(2.2.1)

$$
B(f, x)=\eta(x)
$$

for all $x \varepsilon \pi$. Indeed, since any character must vanish on $[\pi, \pi]=\sum_{i \geq 2} d_{i}$, it is completely determined by its restriction to $d_{1}$. On the other hand, since $d_{1}$ and $d_{-1}$ are dual under the Killing form, and since $d_{-1}$ is orthogonal to $[\pi, \pi]$, the correspondence follows.

We say that $f \varepsilon d_{-1}$ corresponds to $\eta: \eta \rightarrow \mathbb{C}$ if (2.2.1) holds. If $f$ is admissible, then we say that $\eta$ is admissible.

We now define an action of $n$ on $U(\bar{\rho})$ related to the --action of $n$ on $S(\bar{p})$.

Let $n: U(n) \rightarrow C$ be the homomorphism induced by $\eta: \eta \rightarrow \mathbb{C}$ and let $U_{\eta}(n)$ be the kernel of this homomophism. Then, $U(\eta)=\mathbb{C} 1 \otimes U_{n}(n)$ and since $U(o f)=U(\bar{f}) \otimes U(n)$, we have the decomposition $U(\sigma)=U(\bar{p}) \oplus U(\sigma) U_{n}(n)$. Let $p_{n}$ denote the projection onto the first summand. Further, if $u \varepsilon U(O \gamma)$, we set $\rho_{\eta}(u)=u^{\eta}$.

Now, if $x \in \Omega$ and $u \in U(\bar{\rho})$, we put


\begin{equation*}
x \cdot u=(x u)^{\eta}-\eta(x) u \varepsilon U(\bar{\rho}) . \tag{2.2.2}
\end{equation*}


We refer to this action as the $\eta$-reduced action of $\eta$ on $U(\bar{p})$. It is an easy calculation to show that under the\\
$\eta$-reduced action $U(\bar{\rho})$ is an $\eta$-module.\\
Lemma 2.2.1. Let $u \varepsilon U$ and $x \in \lambda$. Then\\
(2.2.3) $\quad x \cdot u^{n}=[x, u]^{n}$

Proof. This is clear since $x \cdot u=\left(x u^{\eta}\right)^{\eta}-\eta(x) u^{\eta}= \left[x, u^{\eta}\right]^{\eta}+\left(u^{\eta} x\right)^{\eta}-\eta(x) u^{\eta}=\left[x, u^{\eta}\right]^{\eta}$ and the fact that if $u=u^{n}+v$ where $v \varepsilon U(O) U_{n}(n)$, then since $U_{n}(n)$ is ad $\eta$-stable $[x, v]^{\eta}=0$.\\
Q.E.D.

We now can relate the --action of $\lambda$ on $S(\bar{p})$ and the $\eta$-reduced action of $\eta$ on $U(\bar{\rho})$ if $f$ corresponds to $n$.

Theorem 2.2. Let $\mathrm{p}, \mathrm{q} \varepsilon \mathbb{Z}_{+}, \mathrm{p}>0, \mathrm{x} \varepsilon \mathrm{d}_{\mathrm{p}} \subseteq \sqcap$ and $u \in U_{(q)}(\bar{p})$. then $x \cdot u \in U_{(q-p)}(\bar{p})$ so that the $n$-reduced action of $n$ induces a graded action of $n$ on $\operatorname{Gr} U(\bar{p})$. Moreover, recalling (2.1.5) one has\\
(2.2.4)

$$
{ }^{\tau}(q-p)^{x \cdot u}=x^{\cdot \tau}(q)^{u}
$$

and hence the map (2.1.11).


\begin{equation*}
\operatorname{Gr} U(\bar{p}) \rightarrow s(\bar{p}) \tag{2.2.5}
\end{equation*}


is an isomorphism of graded $\lambda$-modules. (We say the action of $n$ on a graded vector space $V$ is graded if $x \cdot v \varepsilon V_{q-p}$ for $v \varepsilon V_{q}$ and $\left.x \varepsilon d_{p} \cdot\right)$

Proof. Since $S(\bar{p})$ and $U(\bar{p})$ are $\lambda$-modules and since $d_{1}$ generates $n$, it suffices to take $x \varepsilon d_{1}$. Let $\left\{x_{i}\right\}$ be a basis for $d_{1}$. Then, if $u \varepsilon U(\bar{p})$ we assert that there exist unique elements $u_{0}, u_{i} \varepsilon U(\bar{\rho})$ such that


\begin{equation*}
[x, u]=u_{0}+\Sigma u_{i} x_{i} \tag{2.2.6}
\end{equation*}


Indeed, $U(\bar{p})=U(\bar{n}) U(m)$ and $[x, U(\bar{n})] \subseteq U(\bar{p})$, $\left[x, U(m) \subseteq \cup U(m) d_{1}\right.$. Since ad $x$ is a derivation, this establishes (2.2.6). The uniqueness follows from the isomorphism $U \cong U(\bar{p}) \otimes U(n)$. To prove the first statement, it suffices (recalling (2.2.1)) to take $u \varepsilon\left(U_{k}(\bar{p})\right)_{j}$ where $j+k=q$. But then $[x, u] \varepsilon U_{k}(q)$. Also, - $\operatorname{ad} x_{0}[x, u]=(j-1)[x, u]$. However, since $U(\bar{p})$ and $\mathrm{d}_{1}$ are normalized by ad $\mathrm{x}_{0}$ and since\\
$U_{k} \cong \underset{r+s=k}{\oplus} U_{r}(\bar{p}) \otimes U_{s}(n)$ it follows from the uniqueness of (2.2.6) that $u_{0} \varepsilon\left(U_{k}(\bar{p})\right)_{j-1}$ and $u_{i} \varepsilon\left(U_{k-1}(\bar{p})\right)_{j}$. Thus $u_{1}, u_{2} \varepsilon U_{(q-1)}(\bar{\rho})$. From this, it follows that\\
(2.2.7)

$$
x \cdot u=u_{0}+\sum u_{i} \eta\left(x_{i}\right)
$$

and hence $x \cdot u \varepsilon U_{(q-1)}(\bar{p})$. This proves the first statement.\\
Again, since $d_{1}$ generates $n$, and since the $\left(U_{k}(\bar{\rho})\right)_{j}$ where $k+j=q$ span $\mathcal{U}_{(q)}(\bar{p})$ to prove the theorem it suffices to show that for $x \in d_{1}, u \varepsilon U_{(q)}(\bar{p})$\\
(2.2.8)

$$
{ }^{\tau}(q-1)(x \cdot u)=x \cdot \tau(q)^{u} .
$$

But by the definition (2.1.5) of $\tau_{(q-1)}$ and (2.2.7)\\
(2.2.9)

$$
\tau_{(q-1)}(x \cdot u)=\tau_{k} u_{0}+\sum_{i} \eta\left(x_{i}\right) \tau_{k-1} u_{i} .
$$

On the other $r^{\prime}$ and, since $\tau_{k}$ commutes with the adjoint action, on applying $\tau_{k}$ to (2.2.6) one has ad $x\left(\tau_{k} u\right)=\tau_{k} u_{0}+\sum\left(\tau_{k-1} u_{i}\right) x_{i}$ by (2.1.3). But then by Proposition 1.4.2 one has $x \cdot \tau_{k} u=\tau_{k} u_{0}+\Sigma B\left(f, x_{i}\right) \tau_{k-1} u_{i}$. However $\tau_{k} u=\tau_{(q)} u$ since $u \in\left(U_{k}(\bar{p})\right)_{j}$. This together with (2.2.9) proves (2.2.8).\\
Q.E.D.\\
2.3. The structure of $U(\bar{p})^{N}$.

Let $U(\bar{\rho})^{N}$ denote the space of $n$-invariants in $U(\bar{\rho})$\\
for the $\eta$-reduced action. That $U(\bar{\rho})^{N}$ is a subalgebra of $U(\bar{p})$ follows from the following lemma.

Lemma 2.3 [Kostant]. Let $x \in \pi, u \in U(\bar{p})$ and $v \varepsilon U(\bar{p})^{N}$. Then,


\begin{equation*}
x \cdot(u v)=(x \cdot u) v . \tag{2.3.1}
\end{equation*}


In particular, if $u, v \varepsilon U(\bar{\rho})^{N}$, then for any $x \varepsilon n$, $x \cdot(u v)=0$ and hence $U(\bar{p})^{N}$ is a subalgebra of $U(\bar{p})$.

Proof. We have that $x \cdot(u v)=[x, u v]^{\eta}$ for $x \in \Omega$, $u \in U(\bar{q})$ and $v \in U(\bar{\rho})^{N}$ by Lemma 2.2.1. Thus, $x \cdot(u v)=([x, u] v)^{\eta}+(u[x, v])^{\eta}=([x, u] v)^{\eta}+u([x, v])^{\eta}= ([x, u] v)^{\eta}+u(x \cdot v)=([x, u] v)^{\eta}$ since $v \varepsilon U(\bar{\rho})^{N}$. Next, by an argument similar to that leading to (2.2.6), we have that if $\left\{x_{i}\right\}$ is a basis for $n$, then there exist unique elements $u_{0}, u_{i} \varepsilon \bigcup(\bar{p})$ such that $[x, u]= \mathrm{u}_{0}+\sum_{\mathrm{i}} \mathrm{u}_{\mathrm{i}} \mathrm{x}_{\mathrm{i}}$. Thus,

$$
\begin{aligned}
x \cdot(u v) & =\left(\left(u_{0}+\sum_{i} u_{i} x_{i}\right) v\right)^{\eta}=u_{0} v+\sum_{i} u_{i}\left(x_{i} v\right)^{\eta} \\
& =u_{0} v+\sum_{i} u_{i}\left(\left[x_{i}, v\right]^{\eta}+\left(v x_{i}\right)^{\eta}\right) \\
& =u_{0} v+\sum_{i} u_{i}\left(x_{i} \cdot v+v^{\eta}\left(x_{i}\right)\right) \\
& =u_{0} v+\sum_{i} \eta\left(x_{i}\right) u_{i} v \\
& =(x \cdot u) v .
\end{aligned}
$$

The second statement of the lemma is obvious.\\
Q.E.D.

For any $k \in \mathbb{Z}_{+}$, let $U_{(k)}(\bar{\rho})^{N}=U(\bar{\rho})^{N} \cap U_{(k)}(\bar{\rho})$. Since $U(\bar{\rho})^{N}$ is an algebra by Lemma 2.3, the $U_{(k)}(\bar{\rho})^{N}$ define a filtration on $U(\bar{\rho})^{N}$.

The next theorem is crucial to all that follows. In contrast to the corresponding theorem of Kostant [ 9 , Theorem 2.4.1] the fact that $U(\bar{p})^{N}$ is not the projection onto $U(\bar{p})$ of the center of $U(y)$ renders the proof rather difficult.

Theorem 2.3. For any $k \in \mathbb{Z}_{+}$, the map ${ }^{\tau}(k)$ carries $U_{(k)}(\bar{p})^{N}$ into $S_{(k)}(\bar{p})^{N}$. In fact, $\tau_{(k)}$ is onto and we have an exact sequence


\begin{equation*}
0 \rightarrow U_{(k-1)}(\bar{\rho})^{N} \xrightarrow{\text { inj }}>U_{(k)}(\bar{\rho})^{N} \xrightarrow{\tau}(k)>s_{(k)}(\bar{\rho})^{N} \rightarrow 0 . \tag{2.3.2}
\end{equation*}


Proof. If $u \varepsilon \|_{(k)}(\bar{\rho})^{N}$ and $x \varepsilon d_{i}, i>0$, then by Theorem 2.2 we have $0=\tau_{(k-i)}(x \cdot u)=x^{\cdot} \tau_{(k)} u$ so that ${ }_{(k)}$ carries $U_{(k)}(\bar{p})^{N}$ into $s_{(k)}(\bar{p})^{N}$. By the exact sequence (2.1.6) it is clear that the kernel of this map is $U_{(k-1)}(\bar{p})^{N}$. Thus, to prove the theorem, we need only show that $\tau_{(k)}$ is onto.

Let $v \in S_{(k)}(\bar{\rho})^{N}$. Again, by the exact sequence (2.1.6), we may find an element $u_{1} \varepsilon U_{(k)}(\bar{p})$ such that $\tau_{(k)} u_{1}=v$. But then, by Theorem 2.2 we have that if $x \varepsilon d_{i}, i>0$, then $\tau_{(k-i)}\left(x \cdot u_{1}\right)=x \cdot \tau(k) u_{1}=x \cdot v=0$. Thus, for any $x \varepsilon d_{i}$ we have that $x \cdot u_{1} \varepsilon U_{(k-i-1)}(\bar{p})$. Consider the following condition. An element $u_{r, I} \varepsilon U_{(k)}(\bar{\rho})$ is said to satisfy Condition $(\mathrm{r}, \mathrm{I})$ if $\mathrm{r} \geq \mathrm{l}$,

$$
\text { 1) } \tau_{(k)} u_{r, I}=v \text {, and }
$$

\[
\text { 2) } x \cdot u_{r, I} \varepsilon\left\{\begin{array}{l}
U_{(k-i-r-1)}(\bar{\rho}) \text { if } x \varepsilon d_{i} \text { with } i>I  \tag{2.3.3}\\
U_{(k-i-r)}(\bar{\rho}) \text { if } x \varepsilon d_{i} \text { with } i \leq I .
\end{array}\right.
\]

We will show that if there exists an element $u_{r, I}$ satisfying Condition ( $r, I$ ), then there is an element $u_{r, I-l}$ satisfying Condition ( $r, I-1$ ). But then, Condition ( $r, 0$ ) implies that $x \cdot u_{r, 0} \varepsilon U_{(k-i-r-1)}(\bar{p})$ for all $x \varepsilon d_{i}$ so that if we take $I$ large enough so that $d_{I}=\{0\}$, then Condition ( $r, 0$ ) is the same as Condition ( $r+1, I$ ). Thus, by induction we will be able to make $r$ arbitrarily large. However, if $r=k$ then if $x \in d_{i}$ and if $u_{k}, I$ satisfies Condition $(k, I)$ then $x \cdot u_{k, I}=0$ and\\
$u_{k, I} \varepsilon U_{(k)}(\bar{p})^{N}$ is such that $\tau_{(k)} u_{k, I}=v$. Thus, it only remains to find the element $u_{r, I-1}$.

Let $\left\{x_{\alpha}\right\}_{\alpha=1}^{n}$ be a basis for $d_{I}$ and let $\left\{y_{\beta}\right\}$ be a collection of elements in $d_{-I+1}$ such that $B\left(f,\left[x_{\alpha}, y_{\beta}\right]\right)=\delta_{\alpha \beta}$. Let $\Lambda=\left(\mathbb{Z}_{+}\right)^{n}$ and for any $J=\left(j_{1}, \ldots, j_{n}\right) \varepsilon \Lambda$ set $x_{J}=x_{1}^{j_{1}} x_{2}^{j_{2}} \ldots x_{n}^{j_{n}} \varepsilon U(n)$ and $y_{J}=y_{1}^{j_{1}} y_{2}^{j_{2}} \cdots y_{n}^{j_{n}} \varepsilon U(\bar{p})$. We set $|J|=j_{1}+\ldots+j_{n}$ and $v(J)=\left(j_{1}!\right) \cdot \ldots \cdot\left(j_{n}!\right)$. Lastly, we set $x_{J} \cdot u_{r, I}=x_{1} \cdot x_{1} \cdot \ldots x_{1} \cdot x_{2} \cdot \ldots \cdot x_{n} \cdot \ldots \cdot x_{n} \cdot u_{r, I}$, where if $j_{1}$ times $\quad j_{n}$ times $J=(0, \ldots, 0)$ we set $x_{J} \cdot u_{r, I}=u_{r, I}$ and $y_{J}=1$. We show that $u_{r, I-1}=\sum_{J \varepsilon \Lambda}\left(-\frac{1}{N}\right)(v(J))^{-1} y_{J}\left(x_{J} \cdot u_{r, I}\right)$ satisfies Condition ( $\mathrm{r}, \mathrm{I}-1$ ). We do this by stages.\\
I) $\tau_{(k)} u_{r, I-1}=v$

Proof. By definition, the term in the sum defining $u_{r, I-1}$ for $J=(0, \ldots, 0)$ is $u_{r, I}$ and by assumption $\tau_{(k)} u_{r, I}=v$. Thus, we need to show that if $J \neq(0, \ldots, 0)$, then the term in the sum corresponding to $J$ is in $U_{(k-1)}(\bar{p})$. Now, $y_{\alpha} \varepsilon U_{(I)}(\bar{p})$ so that $y_{J} \varepsilon U_{(I|J|)}(\bar{p})$.

On the other hand, since $x_{\alpha} \cdot u_{r, I} \varepsilon U_{(k-I-r)}(\bar{\rho})$, then by Theorem 2.2 and induction, $x_{J} \cdot u_{r, I} \varepsilon U_{(k-I|J|-r)}(\bar{p})$ so that $y_{J}\left(x_{J} \cdot u_{r, I}\right) \varepsilon U_{(k-r)}(\bar{\rho})$. However, $r \geq 1$ and we have shown (I).\\
(II) If $x \in d_{i}$ with $i<I$, then $x \cdot u_{r, I-1} \varepsilon U_{(k-i-r)}(\bar{\rho})$.

As above, if $J=(0, \ldots, 0)$, then $u_{r, I}$ is the term corresponding to $J$ and $x \cdot u_{r, I} \varepsilon U_{(k-i-r)}(\bar{\rho})$ by assumption. Again, if $J \neq(0, \ldots, 0)$, then the term $y_{J}\left(x_{J} \cdot u_{r, I}\right) \varepsilon U_{(k-r)}(\bar{p})$ as we showed above. But then by Theorem 2.2, $x \cdot\left(y_{J}\left(x_{J} \cdot u_{r, I}\right)\right) \varepsilon U_{(k-i-r)}(\bar{\rho})$.\\
III) If $x \in d_{i}$ with $i>I$, then $x \cdot u_{r, I-1} \varepsilon U_{(k-i-r-1)}(\bar{p})$.

As above, if $J=(0, \ldots, 0)$, then $u_{r, I}$ is the term corresponding to $J$ and $x \cdot u_{r, I} \varepsilon U_{(k-i-r-1)}(\bar{p})$ by assumption. Now, if $J \neq(0, \ldots, 0)$, then $\left.x \cdot\left(y_{J} x_{J} \cdot u_{r, I}\right)\right)= \left(\left[x, y_{J}\right]\left(x_{J} \cdot u_{r, I}\right)\right)^{\eta}+\left(y_{J}\left[x, x_{J} \cdot u_{r, I}\right]\right)^{\eta}=\left(\left[x, y_{J}\right]\left(x_{J} u_{r, I}\right)\right)^{\eta}+ y_{J}\left(x \cdot\left(x_{J} \cdot u_{r, I}\right)\right)$. We consider these two terms separately. Write $y_{J}=y_{\alpha_{1}} \cdot \ldots \cdot y_{\alpha_{|J|}}$. Then

$$
\begin{aligned}
& {\left[x, y_{J}\right]=\sum_{\ell=1}^{J} y_{\alpha_{1}} \ldots .\left[x, x_{\alpha_{\ell}}\right] \cdot \ldots \cdot y_{\alpha_{|J|}} \text {. However, since }} \\
& i>I,\left[x, x_{\alpha_{\ell}}\right] \varepsilon d_{j} \text { with } j>1 \text {. Thus, } \\
& \left(\left[x, y_{J}\right]\left(x_{J} \cdot u_{r, I}\right)\right)^{\eta}=\sum_{\ell=1}^{|J|}\left(y_{\alpha_{1}} \cdots y_{\alpha_{\ell-1}}\left[x_{,} y_{\alpha_{\ell}}\right] y_{\alpha_{\ell+1}} \cdots y_{\alpha_{|J|}}\left(x_{J} \cdot u_{r, I}\right)\right) \eta \\
& =\sum_{\ell=1}^{|J|} y_{\alpha_{1}} \cdots y_{\alpha_{\ell-1}}\left\{\left[\left[x, y_{\alpha_{\ell}}\right], y_{\alpha_{\ell+1}} \cdots y_{\alpha_{\mid J} \mid}\left(x_{J} \cdot u_{r, I}\right)\right]+\right. \\
& \left.y_{\alpha_{\ell+1}} \cdots y_{\alpha|J|}\left(x_{J} \cdot u_{r, I}\right)\left[x, y_{\alpha_{\ell}}\right]\right\}^{\eta} \\
& =\sum_{\ell=1}^{|J|} y_{\alpha_{1}} \cdots y_{\alpha_{\ell-1}}\left\{\left[x, y_{\alpha_{\ell}}\right] \cdot\left(y_{\alpha_{\ell+1}} \cdots y_{\alpha_{|J|}}\left(x_{J} \cdot u_{r, I}\right)\right)\right\} \\
& \left.\varepsilon \sum_{\ell=1}^{|J|} \hat{U}_{((\ell-1) I)}(\bar{p}) U_{((|J|-\ell)}+(k-|J| I-r)-(i-I+1)\right)(\bar{p}) \\
& \subseteq U_{(k-r-i-1)}(\bar{p}) \text { since } j=i-(I-1) \text { above. }
\end{aligned}
$$

Lastly, we need to show that $\left(x^{\circ}\left(x_{J} \cdot u_{r, I}\right)\right) \varepsilon U_{(k-I|J|-r-i-l)}(\bar{\rho})$.\\
We show this by induction. If $J=(0, \ldots, 0)$, then for any $x^{\prime} \varepsilon d_{i^{\prime}}$ with $i^{\prime}>I$ we have that $x_{J} u_{r, I}=u_{r, I}$ and $x^{\prime} \cdot u_{r, I} \varepsilon \mathcal{U}_{(k-r-1)}(\bar{p})$ by assumption. Suppose that we have shown for all $\left|J^{\prime}\right|<|J|$ and $x^{\prime} \varepsilon d_{i}$, with $i^{\prime}>I$ that $x^{\prime} \cdot\left(x_{J^{\prime}} \cdot u_{r, I}\right) \varepsilon U_{\left(k-I\left|J^{\prime}\right|-r-i^{\prime}-I\right)}(\bar{\rho})$.

Then, if we set $x_{J}=x_{\alpha_{1}} \cdots \cdot x_{\alpha_{|J|}}$


\begin{align*}
& x \cdot\left(x_{J} \cdot u_{r, I}\right)=x \cdot x_{\alpha_{1}} \cdot \ldots \cdot x_{\alpha|J|} \cdot u_{r, I} \\
&=\left[x_{,}\left[x_{\alpha_{1}}, x_{\alpha_{2}} \cdot \ldots \cdot x_{\alpha_{|J|}} \cdot u_{r, I}\right]\right]  \tag{by2.2.3}\\
&=\left[x_{\alpha_{1}},\left[x, x_{\alpha_{2}} \cdot \ldots \cdot x_{\alpha|J|} \cdot u_{r, I}\right]\right] \\
&+\left[\left[x_{\alpha_{1}}, x\right], x_{\alpha_{2}} \cdot \ldots \cdot x_{\alpha|J|} \cdot u_{r, I}\right] \\
&=x_{\alpha_{1}} \cdot x \cdot x_{\alpha_{2}} \cdot \ldots \cdot x_{\alpha|J|} \cdot u_{r, I} \\
&+\left[x_{\alpha_{1}}, x\right] \cdot x_{\alpha_{2}} \cdot \ldots \cdot x_{\alpha|J|} \cdot u_{r, I} \\
& \text { However } x \cdot x_{\alpha_{2}} \cdot \ldots \cdot x_{\alpha|J|} \cdot u_{r, I} \varepsilon U_{(k-I(J-1)-r-i-1)}(\bar{\rho}) \text { and } \\
& \begin{aligned}
{\left[x_{\alpha}, x\right] } & x_{\alpha_{2}} \cdot \ldots \cdot x_{\alpha|J|} \cdot u_{r, I} \varepsilon U_{(k-I(J-1)-r-(i+I)-1)}(\bar{\rho}) \text { by }
\end{aligned}
\end{align*}


the induction hypothesis. Thus, $\mathbf{x} \cdot\left(\mathrm{x}_{\mathrm{J}} \cdot \mathrm{u}_{\mathrm{r}, \mathrm{I}}\right) \varepsilon U_{(\mathrm{k}-\mathrm{I}|\mathrm{J}|-\mathrm{r}-\mathrm{i}-1)}(\overline{\mathrm{p}})$ and we have shown (III).\\
(IV) If $x \varepsilon d_{I}$, then $x \cdot u_{r, I-1} \varepsilon U_{(k-I-r-1)}(\bar{p})$.

Proof. As above, we have that

$$
\begin{aligned}
& x \cdot u_{r, I-1} \\
& =\sum_{j \varepsilon \Lambda}(-1)^{|J|}(v(J))^{-1}\left\{\left(\left[x, y_{J}\right]\left(x_{J} \cdot u_{r, I}\right)\right)+y_{J}\left(x \cdot\left(x_{J} u_{r, I}\right)\right)\right\}^{n} .
\end{aligned}
$$

Clearly, we may take $x=x_{\alpha}$ for some $\alpha$. Consider the term $\left(x_{\alpha} \cdot\left(x_{J} \cdot u_{r, I}\right)\right)$. Then, since $\left[x_{\alpha}, x_{\beta}\right] \varepsilon d_{2 I}$, we have that by the proof of the last portion of (III), that

$$
\begin{aligned}
x_{\alpha} \cdot x_{\alpha_{1}} \cdot \ldots \cdot x_{\alpha|J|} \cdot u_{r, I} & =x_{\alpha_{1}} \cdot x_{\alpha} \cdot x_{\alpha_{2}} \cdot \ldots \cdot x_{\alpha|J|} \cdot u_{r, I} \\
& +\left[x_{\alpha} \cdot x_{\alpha_{1}}\right] \cdot x_{\alpha_{2}} \cdot \ldots \cdot x_{\alpha|J|} \cdot u_{r, I} \cdot
\end{aligned}
$$

However, the second term lies in $U_{(k-(|J|-1) I-2|J|-r-1)}(\bar{p})$.\\
Thus, $x_{\alpha} \cdot x_{\alpha_{1}} \cdot \ldots \cdot x_{\alpha_{|J|}} \cdot u_{r, I}$ differs from a term $x_{J} \cdot u_{r, I}$ where $J^{\prime}$ is the same as $J$ except that $j_{\alpha}^{\prime}=j_{\alpha}+1$, by a term in $U_{(k-(|J|+1) I-r-1)}(\bar{p})$. Next consider the $\operatorname{term}\left(\left[x, y_{J}\right]\left(x_{J} \cdot u_{r, I}\right)\right)^{\eta}$. Then, if $y_{J}=y_{\alpha_{1}} \cdots y_{\alpha|J|}$

$$
\begin{aligned}
& \left(\left[x, y_{J}\right]\left(x_{J} \cdot u_{r, I}\right)\right)^{\eta} \\
& =\sum_{\ell=1}^{|J|} y_{\alpha_{1}} \cdots y_{\alpha_{\ell-1}}\left(\left[x_{,} y_{\ell}\right] y_{\alpha_{\ell+1}} \cdots y_{\alpha_{|J|}}\left(x_{J} \cdot u_{r, I}\right)\right)^{\eta} \\
& =\sum_{\ell=1}^{|J|} y_{\alpha_{1}} \cdots y_{\alpha_{\ell-1}}\left\{y_{\alpha_{\ell+1}} \cdots y_{\alpha_{|J|}}\left(x_{J} \cdot u_{r, I}\right)\left[x, y_{\ell}\right]{ }^{\eta}\right. \\
& \left.+\left[x, y_{\ell}\right] \cdot\left(y_{\alpha_{\ell+1}} \cdots y_{\alpha_{|J|}}\left(x_{J} \cdot u_{r, I}\right)\right)\right\}
\end{aligned}
$$

$$
\begin{aligned}
& =\sum_{\ell=1}^{|J|} y_{\alpha_{1}} \cdots y_{\alpha_{\ell-1}} y_{\alpha_{\ell+1}} \cdots y_{\alpha_{|J|}}\left(x_{J} \cdot u_{r, I}\right) n\left[x_{\alpha}, y_{\ell}\right] \\
& =\sum_{\ell=1}^{J \mid} y_{\alpha_{1}} \cdots y_{\alpha_{\ell-1}} y_{\alpha_{\ell+1}} \cdots y_{\alpha_{|J|}}\left(x_{J} \cdot u_{r, I}\right) \delta_{\alpha, \alpha_{\ell}}
\end{aligned}
$$

modulo terms in $U_{(k-I-r-1)}(\bar{p})$ as is easily seen. Let $J_{\alpha}$ denote the term in $\Lambda$ obtained from $J$ by replacing $j_{\alpha}$ by $j_{\alpha}-1$. Then, modulo $U_{(k-I-r-1)}(\bar{\rho})$ we have

$$
\begin{aligned}
& \sum_{j \varepsilon \Lambda}(-1)^{J}(v(J))^{-1}\left\{\left(\left[x_{\alpha} y_{J}\right]\left(x_{J} \cdot u_{r, I}\right)\right)^{\eta}+y_{J}\left(x_{\alpha} \cdot\left(x_{J} u_{r, I}\right)\right)\right\} \\
& =\sum_{j \varepsilon \Lambda}(-1)^{J}(v(J))^{-1}\left\{j_{\alpha} y_{J_{\alpha}}\left(x_{J} \cdot u_{r, I}\right)+y_{J}\left(x_{J} \cdot u_{r, I}\right)\right\}
\end{aligned}
$$

Thus, to prove IV, we need only note that if $J_{\alpha}=\tilde{J}$, then $\mathrm{J}=\tilde{\mathrm{J}}^{\prime},|\mathrm{J}|=|\tilde{\mathrm{J}}|+1$, and $\mathrm{j}_{\alpha}(\mathrm{v}(\mathrm{J}))^{-1}=\mathrm{v}(\tilde{\mathrm{J}})$. But this is obvious. Moreover, since $u_{r, I-1}$ satisfies (I-IV), it satisfies Condition ( $r, I-1$ ) and by the induction argument outlined above there is an element $u \varepsilon U_{(k)}(\bar{p})^{N}$ such that $\tau_{(k)} u=v$. This proves the theorem.\\
Q.E.D.

Corollary 2.3.1. $U(\bar{\rho})^{N}$ is a finitely generated subalgebra of $U(\bar{p})$.

Proof. Since $S(\bar{p})^{N}$ is finitely generated, this is clear from Theorem 2.3.

Corollary 2.3.2. Let $\tilde{\pi}$ be the projection of $U(\bar{p})$ onto $U(m)$ with kernel $n \|(\bar{p})$. Then $\tilde{\pi}$ restricted to $U(\bar{p})^{N}$ is injective. Moreover, if $\pi: S(\bar{p}) \rightarrow S(m)$ is defined as in §1.8, we have that\\
(2.3.4)

$$
\tau_{(i)} \tilde{\pi} u=\pi \tau(i) \mathrm{u}
$$

for any $u \in U_{(i)}(\bar{p})$.\\
Proof. Clearly, $\tau_{(i)} \tilde{\pi} u=\tau_{i} \tilde{\pi} u$ and (2.3.4) is just (2.1.2) for $\Omega=m$. The fact that $\tilde{\pi}$ restricted to $U(\bar{\rho})^{N}$ is injective follows from the fact that $\pi$ is injective on $S(\bar{p})^{N}$ and (2.3.2).\\
Q.E.D.

Nota that if $u \in U_{(i)}(\bar{p})$ and $v \varepsilon U_{(j)}(\bar{\rho})$, then $\tilde{\pi}(u v)=\tilde{\pi}(u) \tilde{\pi}(v)$ so that we may regard $U(\bar{\rho})^{N}$ as a submodule of $U(m)$.\\
2.4. The decomposition $U(\bar{p})=\tilde{A} \otimes U(\bar{p})^{N}$ [Kostant]

Now, let $f$ be an admissible nilpotent and let $\eta: \eta \rightarrow \mathbb{C}$ correspond to $f$.

Recall the decomposition (1.5.9) of $S(\bar{p})$ so that as a tensor product of graded algebras we can write $S(\bar{p})=A \otimes S(\bar{p})^{N}$. But now by (2.1.6) let $\tilde{A}_{(k)} \subseteq U_{(k)}(\bar{p})$ be any fixed subspace such that ${ }^{\tau}(k)$ induces a linear isomorphism\\
(2.4.1) $\quad \tilde{A}_{(k)} \rightarrow \mathrm{A}_{(k)}$.

Now put $\tilde{A}=\sum_{j=0}^{\infty} \tilde{A}(j)$. The sum is direct since $\tilde{A}_{(k)} \cap U_{(k-1)}(\bar{p})=\{0\}$ and $\tilde{A}_{(j)} \subseteq U_{(k-1)}(\bar{p})$ for $j<k$. Now regard $U(\bar{p})$ as a right $U(\bar{p})^{\mathrm{N}}$-module with respect to right multiplication. (One notes of course that $U(\bar{\rho})^{N}$ is not, in general, in the center of $U(\bar{\rho})$. See Chapter 8) It is convenient to choose a basis of $\tilde{A}, u_{i}, i=1,2, \ldots$ such that for each $i \varepsilon \mathbb{Z}_{+}$one has a number $\gamma(i) \varepsilon \mathbb{Z}_{+}$satisfying $u_{i} \varepsilon \tilde{A}_{\gamma(i)}$. One then notes that we may choose the $u_{i}$ so that if $v_{i}=\tau_{\gamma(i)} u_{i} \varepsilon A_{\gamma(i)}$, then the elements $v_{i}$ are all nonzero and form a basis for $A$.

Theorem 2.4. $U(\bar{p})$ is free as a right $U(\bar{p})^{\mathrm{N}}$ module. In fact, right multiplication induces an isomorphism


\begin{equation*}
\tilde{A} \otimes U(\bar{p})^{N} \rightarrow U(\bar{p}) \tag{2.4.2}
\end{equation*}


of right $U(\bar{\rho})^{\mathrm{N}}$ modules. Furthermore, for any $k \varepsilon \mathbb{Z}_{+}$ the map (2.4.2) induces an isomorphism\\
(2.4.3) $\underset{p+q=k}{\oplus}\left(\tilde{A}(p) \otimes U_{(q)}(\bar{p})^{N}\right) \rightarrow U_{(k)}(\bar{p})$.

Proof. Of course, $C_{1}=U_{(0)}(\bar{\rho})=\tilde{A}_{(0)} \otimes U_{(0)}(\bar{\rho})^{N}$. Assume inductively that $U_{(j)}(\bar{\rho})$ is the image of the map (2.4.3) where $j$ replaces $k$ and where $j<k$. Let $u \varepsilon \|_{(k)}(\bar{p})$ and $v=\tau_{(k)} u \varepsilon S_{(k)}(\bar{p})$. Then, by (1.5.9) we can find unique elements $z_{i} \varepsilon S_{(k-\gamma(i))}(\bar{p})^{N}$ such that $\gamma(i) \leq k$ and $v=\sum v_{i} z_{i}$. By (2.3.2) we can find $w_{i} \varepsilon U_{(k-\gamma(i))}(\bar{p})^{N}$ such that $\tau_{(k-\gamma(i))}\left(w_{i}\right)=z_{i}$. By (2.1.6) we have that $u-u^{\prime} \varepsilon U_{(k-1)}(\bar{p})$ where $u^{\prime}$ is the image of the map (2.4.3) where $u^{\prime}=\Sigma u_{i} w_{i}$. Thus $u$ and hence $U_{(k)}(\bar{\rho})$ is in the image of (2.4.3). Thus, (2.4.2) is surjective. Uniqueness is proven in the standard way. Any non-trivial relation $\sum u_{i} w_{i}=0$ for $w_{i} \varepsilon U(\bar{\rho})^{N}$ will imply, by considering highest degree terms, a non-trivial relation $\Sigma v_{i} z_{i}=0$ where $z_{i} \varepsilon S(\bar{\rho})^{N}$. Q.E.D.\\
2.5. The isomorphism $\tilde{A} \rightarrow A$ of $\pi$-modules [Kostant] By (2.4.2) and (2.3.1), the $\eta$-reduced action of $\eta$\\
is determined as soon as one knows $x \cdot v$ for $v \varepsilon \tilde{A}$.\\
Now let $j, k \in \mathbb{Z}_{+}, j \geq 1$ and let $v \varepsilon \tilde{A}_{(k)}$ and $x \varepsilon d_{j} \subseteq \eta$. Then $x \cdot v \varepsilon U_{(k-j)}(\bar{\rho})$ by Theorem 2.2. Thus, by (2.4.3) $\mathrm{x} \cdot \mathrm{v}$ can be uniquely written\\
(2.5.1)

$$
x \cdot v=\sum_{p=0} u_{p}
$$

where $u_{p} \varepsilon \tilde{A}_{(p)} \otimes l_{(q)}(\bar{p})^{N}$ and $p+q=k-j$. Let\\
$\mathrm{x}^{*} \mathrm{v}=\mathrm{u}_{\mathrm{k}-\mathrm{j}}$ so that\\
(2.5.2)

$$
x^{*} v \varepsilon \tilde{A}_{(k-j)} .
$$

Now recall that $A \subseteq S(\bar{p})$ is a $n$-submodule of $S(\bar{p})$. The following lemma asserts that "to first order" $\mathrm{x}{ }^{*} \mathrm{v}=\mathrm{x} \cdot \mathrm{v}$.

Lemma 2.5.1. Let the notation be as above and let $w \varepsilon U(\bar{\rho})$ be defined by $x \cdot v=x * v+w$. Then $w \varepsilon U_{(k-j-1)}(\bar{p})$. Furthermore,


\begin{equation*}
x \cdot \tau(k) v=\tau(k-j)^{x * v} . \tag{2.5.3}
\end{equation*}


Proof. By Theorem 2.2, one has $\tau_{(k-j)}(x \cdot v)=x^{\cdot} \tau(k) v$. But then by (2.5.1) one has $x \cdot \tau(k) v=\sum_{p=0} v_{p}$ where\\
$v_{p}=\tau_{(k-j)} u_{p}$. But then $v_{p} \varepsilon A_{(p)} \otimes S_{(q)}(\bar{p})^{N}$ by (2.1.6) where $p+q=k-j$. However $x^{\cdot} \tau_{(k)} v \varepsilon A_{(k-j)}$ by 1.5.1 and 1.4.3. Thus, $v_{p}=0$ for $p \neq k-j$ and hence $x^{\cdot} \tau_{(k)} v=v_{k-j}$. But $v_{k-j}=\tau_{(k-j)} u_{k-j}=\tau_{(k-j)} x^{*} v$. This proves (2.5.3). However, the right hand side of (2.5.3) is just ${ }^{\tau}(k-j) \mathrm{x} \cdot v$. Since $\tau_{(k-j)}\left(x \cdot v-x^{*} v\right)=0$, we see that $w \varepsilon \dot{u}_{(k-j-1)}(\bar{p})$ by (2.1.6).\\
Q.E.D.

Remark 2.5. Note that Lemma 2.7.1 implies that $\mathrm{x} \cdot \mathrm{v}=0$ if and only if $\mathrm{x}^{*} \mathrm{v}=0$ and $\mathrm{w}=0$.

Now, by linearity extend the definition of $x^{*} v \in \tilde{A}$ so that it is defined for all $x \in \Omega$ and $v \varepsilon \tilde{A}$.

Proposition 2.5. $\tilde{A}$ is an $n$-module with respect to the action $x^{*} v, x \varepsilon \eta, v \varepsilon \tilde{A}$. Furthermore, if\\
(2.5.4) $\quad \tilde{A} \rightarrow A$\\
is the linear isomorphism of graded vector spaces defined by the maps (2.4.1) then (2.5.4) is an isomorphism of几-modules.

Proof. This is immediate from (2.5.3) and Remark 1.4.\\
Q.E.D.

Lemma 2.5.2. Let $j \varepsilon \mathbb{Z}_{+}$and let $v \varepsilon \tilde{A}_{j}$. Then, if $\left\{x_{i}\right\}$ is a basis for $d_{1}, x_{i}{ }^{*} v \varepsilon \tilde{A}_{j-1}$. Furthermore, if $j \geq 1$ and $v \neq 0$, there exists $i$ such that $x_{i}{ }^{*} v \neq 0$. In such a case one also has $\mathrm{x}_{\mathrm{i}} \cdot \mathrm{v} \neq 0$.

Proof. The first statement follows from (2.5.2). The next statement follows from Proposition 2.5 and Remark 1.4. The final statement follows from Remark 2.5.\\
Q.E.D.

Another useful lemma is

Lemma 2.5.3. Let $k \varepsilon \mathbb{Z}_{+}$and $u \varepsilon \tilde{A}_{(k)}$. Also let $1 \leq i \leq \ell$ so that $x_{i}{ }^{*} u \in \tilde{A}_{(k-1)}$ by Lemma 2.5.2. Now for any $v \varepsilon U(\bar{p})^{N}$ let $s \varepsilon U(\bar{p})$ be defined so that


\begin{equation*}
x_{i} \cdot(u v)=\left(x_{i}{ }^{*} u\right) v+s \tag{2.5.5}
\end{equation*}


then $s \varepsilon \Sigma \tilde{A}_{j} \otimes U(\bar{\rho})^{N}$ where the sum only goes from $j=0$ to $j=k-2$.

Proof. One has $x_{i} \cdot(u v)=\left(x_{i} \cdot u\right) v$ by (2.3.1). But then $x_{i} \cdot(u v)=\left(x_{i}{ }^{*} u\right) v+w v$ by Lemma 2.5.1 where $w \varepsilon U_{(k-2)}(\bar{p})$. But then the result follows from (2.4.3). Q.E.D.\\
2.6. The image of $\xi$ in $U(\bar{\rho})^{N}$

Let $\delta$ be the center of $U(\sigma)$. The standard filtration on $U(\sigma)$ induces a filtration $\xi_{k}$ in $\xi$. On the other hand, from the complete reducibility of ad of on $U$ one clearly has the well-known exact sequence\\
(2.6.1) $\quad 0 \rightarrow z_{k-1} \xrightarrow{\text { inj }}>z_{k} \xrightarrow{{ }^{\tau} k} s_{k}(\sigma)^{G} \rightarrow 0$\\
for any $k \varepsilon \mathbb{Z}_{+}$. It is well known that there exist elements $I_{j} \varepsilon S_{m_{j}+1}(j)^{G}, j=1, \ldots, \ell^{\prime}=$ complex rank of of, $m_{j} \varepsilon \mathbb{Z}+$ such that\\
(2.6.2)

$$
S(g)^{G}=\left[I_{1}, \ldots, I_{\ell},\right] .
$$

Thus, by (2.6.1), there exist elements $\tilde{I}_{j} \varepsilon U_{m_{j}}$ such that


\begin{equation*}
\gamma=\mathbb{C}\left[\tilde{I}_{1}, \ldots, \tilde{I}_{\ell},\right] \tag{2.6.3}
\end{equation*}


is a polynomial algebra in the $\tilde{\mathrm{I}}_{\mathrm{j}}$.\\
Recall that $\rho_{\eta}: U(\sigma) \rightarrow U(\bar{\rho})$ is the projection map defined in §2.2.

Let $\xi \rightarrow U(m), \mathrm{u} \rightarrow \mathrm{u}^{0}$ be the map induced by the decomposition $U=U(\bar{\rho}) \oplus U n$ where we note that since\\
elements in $\mathcal{Z}$ are $\theta$-invariant the component of $u$ under this decomposition is the same as under the decomposition $U=U(\bar{p}) \oplus \bar{\pi} U$. See Lemma 7.4.2 in [2].

We have the following proposition\\
Proposition 2.6. We have $z^{n} \subseteq U(\bar{\rho})^{N}$. In fact, $z^{n}$ lies in the center of $U(\bar{p})^{N}$. Moreover, for any $k \in \mathbb{Z}_{+}$we have an isomorphism\\
(2.6.4) $\quad \gamma_{k} \rightarrow\left(\gamma_{k}\right)^{n} \subseteq U_{(k)}(\bar{p})^{N}$.

Further, for any $u \varepsilon \xi, \tilde{\pi} u^{\eta}=u^{0}$ where $\tilde{\pi}$ is the projection map defined in Corollary 2.3.2.

Proof. If $u \varepsilon \gamma$ and $x \varepsilon \eta$, then $x \cdot u^{\eta}=[x, u]^{\eta}=0$ by (2.3.1). Further, if $v \varepsilon U(\bar{\rho})^{\mathrm{N}}$, then $(\mathrm{uv})^{\eta}=(\mathrm{vu})^{\eta}+[\mathrm{v}, \mathrm{u}]^{\eta}=\mathrm{v} \mathrm{u}^{\eta}$. This proves the first two statements.

Now, one has the direct sum $U(\bar{p})=\bar{n} U(\bar{p}) \oplus U(m)$ and one notes that for any $\mathrm{v} \varepsilon U(\bar{p})^{\mathrm{N}}, \pi \mathrm{v}$ is the component of $v$ in $U(m)$ under this decomposition. Thus, for any $v \varepsilon U, \pi v^{n}$ is the component of $v$ in $U(m)$ under the decomposition $U=\left(\bar{n} U+U U_{\eta}(n)\right) \oplus U(m)$ where this sum is direct since $U \stackrel{\sim}{=} U(\bar{n}) \otimes U(m) \otimes U(n)$. But then if $u \varepsilon\}$, it is clear that $\pi u^{\eta}=u^{0}$. However, since $\pi$\\
contains a Cartan subalgebra of $O \gamma$, the map $z \rightarrow U(m)$ is injective. Thus, for any $k \varepsilon \mathbb{Z}_{+}, \xi_{k} \rightarrow U_{k}(m)$ is injective. But now, $\pi U_{(k)}(\bar{p})^{N} \rightarrow U_{(k)}(m)=U_{k}(m)$ is injective by Corollary 2.3.2. From this it follows that\\
$\gamma_{k} \xrightarrow{\rho_{\eta}}\left(\gamma_{k}\right)^{\eta} \xrightarrow{\pi} U_{k}(m)$ is injective and thus (2.6.4) is an isomorphism.\\
Q.E.D.

\subsection{Comments}
The map $\tilde{\pi}: U(\bar{\rho})^{N} \rightarrow U(m)$ defined in §2.3 is the analogue of the Harish-Chandra homomorphism $\xi \rightarrow U(h)^{\mathrm{w}}$ since in the case of a Borel subalgebra we have $Z \xrightarrow{\rho_{\eta}} U(\bar{p})^{N} \xrightarrow{\tilde{\pi}} U(h)$ is just the Harish-Chandra homomorphism if $m=h$. See §2.6. The chief difficulty of this theory is that there is no natural (invariant of the choice of $f$ and $\eta$ ) preimage of $U(\bar{p})^{N}$ in $U(\sigma)$ as we will see later by example. In fact, there is a natural preimage if and only if $p$ is a Borel subalgebra. We will see later, when constructing Whittaker modules as submodules of completions of generalized Verme modules, that the representation theory of $U(\bar{p})^{N}$ and that of $U(m)$ are intimately related via the map $\tilde{\pi}_{\text {. }}$

In fact, if $\rho f: S(\mathcal{J}) \rightarrow S(\bar{\rho})$ is the map defined\\
by $\left(\rho_{f} u\right)(x)=u(f+x), u \varepsilon S(g), x \varepsilon p$, it is an easy exercise to show that if $u \varepsilon S(0 j)^{G}$ then $\left(\rho_{f} u\right) \varepsilon S(\bar{p}){ }^{N}$. It then follows by an argument analogous to that of Theorem 2.4.1 in Kostant's paper [9] that\\
(2.7.1)

$$
\begin{array}{ccc}
0 \rightarrow \gamma_{k-1} & \xrightarrow{\text { inj }} \gamma_{k} \xrightarrow{\tau_{k}} & s_{k}(o j)^{G} \rightarrow 0 \\
0 \rightarrow U_{(k-1)}(\bar{\rho})^{N} \xrightarrow{\text { inj }} U_{(k)}(\bar{\rho})^{N} \xrightarrow{\tau(k)} & \downarrow \rho f \\
& \stackrel{\tau}{(k)}(\bar{\rho})^{N} \rightarrow 0
\end{array}
$$

is a commutative diagram. However, we already know (Theorem 2.3) that the lower sequence is exact. For Borel subalgebras $\rho \eta$ and $\rho_{f}$ are isomorphisms above. In general, they are just injections. Although we won't be able to use (2.7.1) in the general case, the continuation of (2.7.1) by the Harish-Chandra homomorphism, ( $\pi, \pi$ ), yields

$$
\begin{array}{lllll}
0 \rightarrow \gamma_{k-1} & \xrightarrow{\text { inj }} & \gamma_{k} & \xrightarrow{\tau_{k}} & s_{k}(q)^{G} \rightarrow 0 \\
0 \rightarrow U_{(k-1)}(\bar{\rho})^{N} \xrightarrow{\text { inj }} U_{(k)}(\bar{\rho})^{N} \xrightarrow{\tau(k)} & s_{(k)}(\bar{\rho})^{N} \rightarrow 0 \\
& \downarrow \tilde{\pi} & & \downarrow \tilde{\pi} & \\
0 \rightarrow & W_{k-1} & \xrightarrow{\text { inj }} & \tilde{w}_{k} & \xrightarrow{\tau_{k}}
\end{array} w_{k} \rightarrow 0
$$

where $\tilde{W}_{k}$ is the image of $U_{(k)}(\bar{p})^{N}$ in $U_{k}(m)$ and $W_{k}$ is the image of $S_{(k)}(\bar{\rho})^{N}$ in $S_{k}(m)$. Now, (2.7.2) is also a commutative diagram by Corollary 2.3 .2 and $\tilde{\pi}$ ( $\pi$ ) are isomorphisms of filtered (graded) algebras. The lower two exact sequences figure prominently in the representation theory of Whittaker modules.

\section{Chapter 3}
\section{Whittaker Modules}
3.1 The annihilator of a Whittaker vector is

$$
U U_{w}(\bar{p})^{N}+U U_{q}(\pi) .
$$

Recall that of is a fixed complex semi-simple Lie algebra with a fixed parabolic subalgebra $p=\pi \oplus \pi$. For any subalgebra $\Omega \subseteq O \gamma, U(\Omega)$ will denote its universal enveloping algebra. We have also fixed an admissible homomorphism $\eta: \pi \rightarrow \mathbb{C}$ which we extend to a homomorphism $U(\pi) \rightarrow C$ which we continue to denote by $n$. As usual, $U=U(g)$.

Now, if $V$ is any $\pi$-module, the action will be denoted by $u v \varepsilon v$ for $u \varepsilon$ 认, $v \varepsilon v$. A vector $w \varepsilon v$ will be called a Whittaker vector (with respect to $\eta$ and $p$ ) if\\
(3.1.1)

$$
x w=\eta(x) w
$$

for all $\mathrm{x} \varepsilon \Omega$. A Whittaker vector will be called a cyclic Whittaker vector (for V ) if $U_{\mathrm{w}}=\mathrm{V}$. A $U$-module V will be called a Whittaker module if it contains a cyclic Whittaker vector.

If $w$ is a cyclic whittaker vector for $V$ and if\\
$q \subseteq o \gamma$, then $U_{w}(00$ will denote the annihilator of $w$ in $U(0)$. Of course, $U_{w}=U_{w}(o f)$ is a left ideal in $U(\sigma)$ and one has a $U$-module isomorphism $v \cong U_{/} U_{w}$ so that V is determined by the left ideal $U_{\mathrm{w}}$. Since $U=U(\bar{p}) \otimes U(\pi)$ we may decompose $U$ as $U=U(\bar{p}) \oplus U U_{n}(\pi)$, where $U_{\eta}(\pi)=\operatorname{ker} \eta$ in $U(\pi)$. But now, one clearly has $U U_{\eta}(\pi) \leq U_{w}$ and hence $v$ is determined by $U_{w}(\bar{p})$. On the other hand, the decomposition $U(\bar{p})=\tilde{A} \otimes U(\bar{p})^{N}$ yields the following theorem:

Theorem 3.1. Let $V$ be any U-module which admits a cyclic whittaker vector $w$ and let $U_{w}$ be the annihilator of $w$. Then


\begin{equation*}
U_{w}=U U_{w}(\bar{p})^{N}+U U_{n}(\pi) \tag{3.1.2}
\end{equation*}


where $U_{w}(\bar{p})^{N}$ is the annihilator of $w$ in $U(\bar{p})^{N}$.\\
By the above discussion, we need only show that $U_{w}(\bar{p})=U(\bar{p}) U_{w}(\bar{p})^{N}=\tilde{A} \otimes U_{w}(\bar{p})^{N}$. Thus, we need only prove the following lemma.

Lemma 3.1. Let $x=\{v \varepsilon U(\bar{p}) \mid(x \cdot v) w=0$ for all $x \varepsilon \eta\}$. Then,


\begin{equation*}
X=\left(\tilde{A} \otimes U_{W}(\bar{\rho})^{N}\right)+U(\bar{p})^{N} \tag{3.1.3}
\end{equation*}


where $x \cdot v$ denotes the $\eta$-reduced action of $\eta$ on $U(\bar{\rho})$. Furthermore, $U_{W}(\bar{p}) \subseteq X$ and in fact


\begin{equation*}
U_{w}(\bar{p})=\tilde{A} \otimes U_{w}(\bar{p})^{N} . \tag{3.1.4}
\end{equation*}


Proof. If $\mathrm{v} \in U(\bar{p})^{\mathrm{N}}$, then $\mathrm{x} \cdot \mathrm{v}=0$ for all $\mathrm{x} \in \pi$ by definition and hence $U(\bar{\rho})^{N} \subseteq x$. Further, if $u \varepsilon U_{w}(\bar{\rho})^{N}$ and if $v \in U(\bar{p})$, then for any $x \in \pi$ one has $\mathrm{x} \cdot(\mathrm{vu})=(\mathrm{x} \cdot \mathrm{v}) \mathrm{u}$ by Proposition 2.3. Thus, $\mathrm{vu} \varepsilon \mathrm{X}$ and hence $\tilde{A} \otimes U_{w}(\bar{\rho})^{N} \subseteq X$. If we let $Y$ denote the right hand side of (3.1.4), then we have shown that $Y \subseteq X$. We wish to show that $X \subseteq Y$.

First, we observe that $U_{w}(\bar{p}) \subseteq x$. But, if $u \varepsilon U_{w}$ and $x \in \pi$, then xuw $=0$ and uxw $=\eta(x) u w=0$. Thus $[x, u] \varepsilon U_{w}$ and hence $[x, u]^{\eta}=x \cdot u \varepsilon U_{w}(\bar{p})$ by (2.2.3). Hence, if $u \varepsilon U_{w}(\bar{p})$, then $u \varepsilon x$.

Next, we embark on a standard approximation argument. Let W be any fixed complementary subspace of $U_{\mathrm{w}}(\bar{\rho})^{\mathrm{N}}$ in $U(\bar{p})^{N}$. If we set $M_{j}=\tilde{A}_{j} \otimes W$ and $M=\sum_{j=1}^{\infty} M_{j}$, then


\begin{equation*}
U(\bar{p})=\mathrm{M} \oplus \mathrm{Y} . \tag{3.1.5}
\end{equation*}


To prove (3.1.3) we need only show that $X \cap M=\{0\}$.

Let $\mathrm{M}_{[k]}=\sum_{1 \leq j \leq k} \mathrm{M}_{j}$ so that the $\mathrm{M}_{[k]}$ are a filtration of M . Assume $\mathrm{X} \cap \mathrm{M} \neq\{0\}$ and let k be minimal so that there exists $0 \neq y \varepsilon X \cap_{[k]}$. We may then uniquely write $y=\sum_{j=1}^{k} y_{j}$ where $y_{j} \varepsilon \tilde{A}_{j}$ ® $w$. Furthermore, we know that $y_{k} \neq 0$ by the definition of $k$. Recall that under the *-action, $\tilde{A}$ is an $\pi$-module. See §2.5. If $\left\{x_{i}\right\}$ is a basis for $d_{l}$, let $C_{i}=\left\{v \varepsilon A_{(k)} \mid x_{i}{ }^{*} v=0\right\}$. By Lemma 2.5.2 one has $\bigcap_{i} c_{i}=\{0\}$. Hence, there exists an $i$ such that $y_{k} \notin C_{i} \otimes W$.

For any $q \in \mathbb{Z}_{+}$, let


\begin{equation*}
\tilde{A}_{[q]}=\sum_{p \leq q} \tilde{A}_{(p)} \tag{3.1.7}
\end{equation*}


We now assert that $x_{i} \cdot y$ can be written as\\
(3.1.7)

$$
x_{i} \cdot y=z_{k-1}+s
$$

where $s \varepsilon \tilde{A}_{[k-2]} \otimes U(\bar{\rho})^{N}$ and\\
(3.1.8)

$$
0 \neq z_{k-1} \varepsilon \tilde{A}_{(k-1)} \otimes w=M_{k-1}
$$

Indeed, by Lemma 2.5.3 one clearly has $x_{i} \cdot y_{j}=\tilde{A}_{[k-1]} \otimes U(\bar{p})^{N}$\\
for any $1 \leq j \leq k-1$ so that we need only consider $x_{i} \cdot y_{k}$. But if $D_{i}$ is a linear complement to $C_{i}$ in $\tilde{A}_{(k)}$ we may write $\mathrm{y}_{\mathrm{k}}=\mathrm{u}_{\mathrm{k}}+\mathrm{v}_{\mathrm{k}}$ where $\mathrm{u}_{\mathrm{k}} \varepsilon \mathrm{C}_{\mathrm{i}} \otimes \mathrm{W}$ and $v_{k} \varepsilon D_{i} \otimes W$ and $v_{k} \neq 0$. But now $x_{i} \cdot u_{k} \varepsilon \tilde{A}_{[k-2]} \otimes U(\bar{\rho})^{N}$ by Lemma 2.5.3. On the other hand, the map $D_{i} \rightarrow \tilde{A}_{(k-1)}$ given by $u \rightarrow x_{i}{ }^{*} u$ is injective. Thus, by Lemma 2.5.3 one has $x_{i} \cdot v_{k}=z_{k-1}+s^{\prime}$ where $z_{k-1}$ satisfies (3.1.8) and $s^{\prime} \varepsilon \tilde{A}_{[k-2]} \otimes U(\bar{p})^{N}$. This proves (3.1.7) and (3.1.8)

Now, $x_{i} \cdot y \varepsilon \mathcal{U}_{w}(\bar{\rho})$ by the definition of $x$. But then if $k=1$, one has $s=0$ and hence $0 \neq z_{0}=x_{i} \cdot y \varepsilon W . \quad$ But $W \cap \mathcal{U}_{W}(\bar{p})=\{0\}$ by definition. Hence $k>1$. But then since $U_{w}(\bar{p}) \subseteq x$, one has $x_{i} \cdot y \varepsilon x$. Writing $U(\bar{p})^{N}=W \oplus U_{W}(\bar{p})^{N}$ we may write $s=s_{1}+s_{2}$ where $s_{1} \varepsilon M_{[k-2]}$ and $s_{2} \varepsilon X$ where we have put $M_{[0]}=\{0\}$. Thus, $x_{i} \cdot y=t_{2}+s_{2}$ where $t_{2}=z_{k-1}+s_{1}$. But then $t_{2} \varepsilon X$ since $x_{i} \cdot y$ and $s_{2} \varepsilon x$. However, $0 \neq t_{2} \varepsilon M_{[k-1]}$ and thus $0 \neq t_{2} \varepsilon x \bigcap_{[k-1]}$ which contradicts the minimality of $k$. Hence $X=Y$. Since $U_{w}(\bar{p}) \subseteq x,(3.1 .4)$ clearly follows from (3.1.3) Q.E.D.\\
3.2 One has $W h \mathrm{~V}=U(\bar{p}) \mathrm{N}_{w}$ for a cyclic $w \varepsilon \mathrm{~V}$.

If $V$ is any $U$-module, let $W h V$ denote the space of Whittaker vectors in V. We then have the following characterization of WhV if $V$ admits a cyclic Whittaker vector.

Theorem 3.2. Let $V$ be any U-module with a cyclic Whittaker vector $w \varepsilon V$. Then any $v \varepsilon V$ is a Whittaker vector if and only if $v$ is of the form $v=u w$ where $u \in U(\bar{p})^{N}$.

Proof. If $u \in U(\bar{p})^{N}$, then xuw $=[x, u] w+u x w= (x \cdot u) w+\eta(x) u w=\eta(x) u w$ so that $u w$ is a Whittaker vector. Conversely, let $\mathrm{v} \varepsilon \mathrm{v}$ be a Whittaker vector. Write $v=u w$ for some $u \varepsilon U$. Then clearly $u w=u^{\eta} w$ so we may assume that $u \varepsilon U(\bar{p})$. But now, if $x \varepsilon \pi$, then xuw $=\eta(x)$ uw. But also, uxw $=\eta(x)$ uw. Thus $[x, u] w=0$ and hence $(x \cdot u) w=[x, u] \eta_{w}=0$ by Lemma 2.2.1. Thus, in the notation of Lemma 3.1 we can write $u=u_{1}+u_{2}$ where $u_{1} \varepsilon \tilde{A} \otimes U_{w}(\bar{\rho})^{N}$ and $u_{2} \varepsilon U(\bar{\rho})^{N}$. But then $u_{1} w=0$. Thus, $v=u_{2} w$ which proves the theorem.\\
Q.E.D.

Remark 3.2. Note that Theorem 3.2 asserts that the space of all Whittaker vectors in V is a cyclic $U(\bar{p})^{\mathrm{N}}$ module which is isomorphic to $U(\bar{p})^{N} / U_{w}(\bar{p})^{N}$.\\
3.3 Whittaker modules and cyclic $U(\bar{p})^{\mathrm{N}}$-modules.

We can now determine, up to equivalence, the collection of all Whittaker modules. They are naturally parameterized by the cyclic $U(\bar{p})^{N}$-modules.

Theorem 3.3. Let $V$ be any Whittaker module for $U$, the universal enveloping algebra of a semi-simple Lie algebra of. See §3.1. Let $W h \mathrm{~V}$ be the cyclic $U(\bar{\rho})^{\mathrm{N}}-$ module of all Whittaker vectors in $V$. (See Remark 3.2). Then the correspondence

$$
(3.3 .1) \quad \mathrm{V} \rightarrow \mathrm{WhV}
$$

is a bijection between the equivalence classes of Whittaker modules and cyclic $U(\bar{p})^{\mathrm{N}}$-modules.

Proof. Let $V_{1}$ and $V_{2}$ be Whittaker modules with cyclic Whittaker vectors $w_{1}$ and $w_{2}$ respectively. If $\mathrm{WhV}_{1}$ is isomorphic to $\mathrm{WhV}_{2}$, we may take $\mathrm{WhV}_{1}=\mathrm{WhV}_{2}$. Since, by Remark 3.2, $w_{1}$ and $w_{2}$ are cyclic $U(\bar{p})^{\mathrm{N}}$-vectors, we can find $u, u^{\prime} \varepsilon U(\bar{\rho})^{N}$ such that $u w_{1}=w_{2}, u^{\prime} w_{2}=w_{1}$. But then, $\mathrm{v}_{1}=U \mathrm{w}_{1}=U \mathrm{u}^{\prime} \mathrm{w}_{2} \subseteq \mathrm{v}_{2}$ and $\mathrm{v}_{2}=U \mathrm{w}_{2}=U u \mathrm{w}_{1} \subseteq \mathrm{v}_{1}$ and thus $\mathrm{V}_{1}=\mathrm{V}_{2}$. Thus, (3.3.1) is injective.

Conversely, let $W$ be any $U(\bar{p})^{N}$-module with a cyclic vector $w$. Let $U_{w}^{\prime}(\bar{p})^{N}$ denote the annihilator of $w$ in\\
$U(\bar{\rho})^{N}$. Now, if $L=\lambda U_{\eta}(n)+U U_{w}^{\prime}(\bar{p})^{N}$, then $V=/ L$ is clearly a Whittaker module where, in the notation of §3.1, $U_{w}=L$. However, since $U_{\eta}(n) \cap U(\bar{p})^{N}=\{0\}$, it follows that $U_{w}(\bar{p})^{N}$ (the annihilator of $w \varepsilon V$ ) is the same as $U_{w}^{\prime}(\bar{p})^{N}$ and the result follows.\\
Q.E.D.\\
3.4. The irreducibility of Whittaker modules

Let $\left(U(\bar{\rho})^{N}\right)^{\wedge}$ denote the collection of equivalence classes of irreducible $U(\bar{p})^{\mathrm{N}}$-modules. If $\quad \xi \in\left(l(\bar{p})^{\mathrm{N}}\right)^{\wedge}$, we let $Y_{\xi, \eta}$ denote the Whittaker module associated with $\xi$ by Theorem 3.3. We now show that the $Y_{\xi, \eta}, \xi \varepsilon\left(\mathcal{U}(\bar{p})^{N}\right)^{\wedge}$ classifies the irreducible Whittaker modules up to equivalence.

Theorem 3.4. [Kostant]. Let $V$ be any Whittaker module for $u$, the universal enveloping algebra of a semi-simple Lie algebra. Then the following are equivalent.\\
(1) V is an irreducible U-module.\\
(2) WhV is an irreducible $U(\bar{p})^{\mathrm{N}}$-module.\\
(3) All (non-zero) Whittaker vectors are cyclic vectors for $V$.\\
(4) $V$ is isomorphic to $Y_{\xi, \eta}$ for some $\xi \varepsilon\left(U(\bar{p})^{N}\right)^{\wedge}$.

Proof. The equivalence of (2), (3), and (4) follows from §'s 3.2 and 3.3. On the other hand, if WhV is reducible, then any non-cyclic (for WhV) Whittaker vector generates a U-submodule of V by Theorem 3.3. Thus (1) ⇒ (2). To prove the theorem, we need only show that (2) => (1).

We now show that any non-trivial $U$-submodule $D \subseteq V$ contains a non-trivial Whittaker vector. This will prove the theorem since then (2) ⇒ (1).

For any $\mathrm{v} \varepsilon \mathrm{V}$ and $\mathrm{x} \varepsilon \Uparrow$, let\\
(3.4.1)

$$
x \cdot v=x v-\eta(x) v .
$$

One easily checks that this defines an action of $\Omega$ on V since we are effectively tensoring with a 1 -dimensional representation of $\pi$. As usual, we refer to this as the $\eta$-reduced action of $\Omega$ on $V$.

Now, if $u \varepsilon U$ and $x \varepsilon \Omega$ note that if $w$ is a cyclic Whittaker vector


\begin{equation*}
x \cdot(u w)=[x, u] w . \tag{3.4.2}
\end{equation*}


Indeed,

$$
x \cdot(u w)=x u w-\eta(x) u w=x u w-u x w . \quad \text { But now if }
$$


\begin{equation*}
U \rightarrow v \tag{3.4.3}
\end{equation*}


is the map $u \rightarrow u w$, then (3.4.3) is a map of $\Omega$-modules where we regard $U$ as an $\mathcal{n}$-module with respect to the adjoint action and $V$ is an $\Omega$-module with respect to the $n$-reduced action.

Now, if $0 \neq \mathrm{v} \varepsilon \mathrm{D}$, write $\mathrm{v}=\mathrm{uw}$ for all $\mathrm{u} \varepsilon U_{\text {. }}$ Then, since $U$ is locally finite with respect to the adjoint action, the ad $\pi$-submodule $U_{0} \subseteq U$ generated by $u$ is finite dimensional. Thus, the image $V_{0}$ of $U_{0}$ in $V$ by (3.4.3) is finite dimensional. Furthermore, since ad $\pi$ operates as a Lie algebra of nilpotent elements on $U_{0}$, the same is true for $V_{0}$. Now, by (3.4.1) $D$ is clearly stable under the $\eta$-reduced action. Thus, $\mathrm{V}_{0} \subseteq \mathrm{D}$. Since $\mathrm{D} \neq\{0\}(\mathrm{v} \varepsilon \mathrm{D})$, Engel's theorem implies that there exists $0 \neq v_{0} \varepsilon v_{0}$ such that $x \cdot v_{0}=0$ for all $x \in \pi$. But then $v_{0}$ is a Whittaker vector.\\
Q.E.D.

Remark 3.4. If V is a U -module and W is an irreducible $U(\bar{p})^{\mathrm{N}}$-submodule of the Whittaker vectors of $V$, then $U W$ is an irreducible $U$-submodule of $V$.

\section{Chapter 4}
\section{Tensor products and $\eta$-finite modules}
Much of the material in this chapter originated in Chapter 4 of Kostant [9].\\
4.1. The $n$-isomorphism $\mathrm{V} \rightarrow \mathcal{U}(n)^{*} \otimes \mathrm{~Wh}(\mathrm{~V})$.

For any $U$-module $v$, the $\eta$-reduced action of $\pi$ on V is defined by putting $\mathrm{x} \cdot \mathrm{v}=\mathrm{xv}-\eta(\mathrm{x}) \mathrm{v}$ for $\mathrm{x} \varepsilon \eta, \mathrm{v} \varepsilon \mathrm{V}$. We extend this to an action of $U(n)$ on $v$.

Let V be a Whittaker module and let $\mathrm{V}=\mathrm{WhV} \oplus \mathrm{V}^{\prime}$ be a direct sum decomposition of $V$. Let $\ell: V \rightarrow$ WhV be the projection onto the first component under this decomposition. Now, if $\mathrm{v} \varepsilon \mathrm{v}$, let $\beta_{\mathrm{v}} \varepsilon U(N)$ ' Whv (the dual to $U(n)$ with values in $W h V$ ) be defined by $\left\langle\beta_{v, u}\right\rangle=\ell(u \cdot v)$. Here $u \rightarrow \check{u}$ is the linear isomorphism of $U(n)$ with itself characterized by $x=-x$. if $x \in \pi$ and $\left(u_{1} u_{2}\right)^{v}=\check{u}_{2} \check{u}_{1}$. Note that $U(n)$ is stable under this map.

Theorem 4.1. [Kostant]. Let $U(n){ }^{*}$ be the restricted dual of $U(\lambda)$. That is $U(n)$ * consists of those elements of $U(\lambda)$ ' which vanish on a power of the augmentation ideal $n U_{(n)}$. Let the notation be as above so that V is a Whittaker module. Then, for any $\mathrm{v} \varepsilon \mathrm{V}, \beta_{\mathrm{V}} \varepsilon U(n) * \otimes W h \mathrm{~V}$.

In fact, the map


\begin{equation*}
\mathrm{V} \rightarrow U(n) * \otimes \mathrm{Whv} \quad \mathrm{v} \rightarrow \dot{\beta}_{\mathrm{v}} \tag{4.1.1}
\end{equation*}


is an isomorphism of $\pi$-modules with respect to the $\eta$-reduced action of $n$ on $V$ and the left translation action of $n$ on $U(n) * \otimes W h \mathrm{~V}$ defined by $(x \cdot(\gamma \otimes w))(u)=-\gamma(x u) \otimes w$ where $x \varepsilon n, \gamma \varepsilon U(n) *, u \varepsilon U(n)$, and $w \varepsilon$ Whv.

Proof. Let $w$ be a fixed cyclic Whittaker vector for $V$ and let $U \rightarrow V$ be given by $y \rightarrow y w$. Then, if $n$ operates on $V$ by the $\eta$-reduced action, we have that this is a map of $n$-modules as in (3.4.2). Now, if $x \in U(n)$ and $y \varepsilon U$, let ad $x(y)$ be the action of $x$ on $y$ (i.e., iteration of commutation). Then $x \cdot(y w)=(\operatorname{ad} x(y)) w$ for all $y \varepsilon U, x \varepsilon U(n)$. Thus, if $v \varepsilon v$ and $v=y w$ for some $y \varepsilon U$, since $x \cdot v=(\operatorname{ad} x(y)) w$ it follows that $x \cdot v=0$ whenever $x \varepsilon(n U(n))^{k}$ for large k . Hence, $\beta_{\mathrm{v}} \varepsilon U(\lambda) * \otimes W h \mathrm{~V}$. Now, if $\mathrm{v} \varepsilon \mathrm{V}, \mathrm{x} \varepsilon \eta$, and $u \in U\left(N\right.$, we have $\left(x \cdot \beta_{v}\right)(u)=-\beta_{v}(x u)=-\ell((x u) \cdot v)= \ell((u x) \cdot v)=\ell(u \cdot(x \cdot v))=\beta_{x \cdot v}(u)$ and thus (4.1.1) is a map of $n$-modules. It remains to show that (4.1.1) is a linear isomorphism.

Now, in the notation of §3.1, let $U_{w}$ denote the\\
annihilator of $w$ in $U$ and let $U_{w}(\bar{p})^{N}$ denote the annihilator of w in $U(\bar{p})^{\mathrm{N}}$. Then, under the map $U \rightarrow \mathrm{~V}, \mathrm{y} \rightarrow \mathrm{yw}$ we have a linear isomorphism\\
(4.1.2)

$$
\tilde{A} \otimes D \rightarrow V
$$

where $D$ is any complementary subspace to $U_{w}(\bar{\rho})^{N}$ in $U(\bar{p})^{N}$ since $\tilde{A} \otimes D \oplus U_{W}=V$ by Theorem 3.1. It is clear (see Theorem 3.2) that under this map $D$ maps isomorphically onto WhV. Let


\begin{equation*}
\tilde{A} \otimes D \rightarrow U(x) * \otimes W h v, y \rightarrow \alpha y \tag{4.1.3}
\end{equation*}


be the composite of (4.1.2) and (4.1.1). To prove the theorem, it suffices to show that (4.1.3) is a linear isomorphism.

Now, if $y \varepsilon A \otimes D$, then $\alpha_{y}$ is given by the relation $\alpha y(x)=\ell(x \cdot y w)$ for $x \in U(n)$. Recalling the $\eta$-reduced action of $\lambda$ on $S(\bar{p})$, one has $(x \cdot y) w=(x y)^{\eta_{w}}-\eta(x) y w-x \cdot(y w)$ for any $x \in U(n)$ and hence $x \cdot(y w)=(x \cdot y) w$ for any $x \in U^{\prime}(n)$. Thus,


\begin{equation*}
\alpha y(x)=\ell((x \cdot y) w) \tag{4.1.4}
\end{equation*}


Now, for $j \varepsilon \mathbb{Z}_{+}$let $U(n)_{j}$ be the $j$-eigenspace for the action of ad $x_{0}$ on $U(n)$. These define a gradation of $U(N)$ By Theorem 1.6, we have that $A(N)_{j}$ is orthogonal to $U_{(n)_{i}}$ for $i \neq j$ and is nonsingularly paired to $U(n)_{j}$ by (1.6.3). Here $A(N)_{j}$ is as in the proof of Theorem 1.6. Let $U(n)_{j}^{*}$ be the image of $A(N)_{j}$ under this pairing. By Lemma 2.5.3, if $\tilde{A}_{[k]}=\sum_{j=0}^{k} \tilde{A}_{(j)}$, then for $y \varepsilon \tilde{A}_{[k]}, x \varepsilon U(n)_{i}$ we have that $x \cdot y=0$ if $i>k$. Since $U_{(n)_{i}}$ is clearly stable under the map $x \rightarrow \dot{x}$ it follows that we have a map

$$
\tilde{A}_{[k]} \otimes D \rightarrow U_{(n)_{[k]}^{*}} \otimes W h V
$$

where $U(n)_{[k]}^{*}=\sum_{j=0}^{k} U(n)_{j}^{*}$. It suffices to show that this is an isomorphism for all $k$. For $k=0$ this is obvious. But then by induction it suffices only to show that (4.1.4) defines a non-singular pairing of $\tilde{A}_{(k)} \otimes D$ and $U(n)_{k}$ with values in WhV. Now, for $x \in U(n)_{k}$ and $y \varepsilon \tilde{A}_{(k)}$ we have that $x \cdot y \varepsilon U_{0}(\bar{p})=c \cdot 1$ by Theorem 2.2. Further, if $z \varepsilon D \subseteq U(\bar{\rho})^{N}$, then $x \cdot(y z)=(x \cdot y) z$ by (2.3.1). Since $D w=$ WhV, it is thus sufficient to show that $\tilde{\mathrm{A}}_{(k)} \otimes 1$ is nonsingularly paired to $U(n)_{k}$ with values in $C w$.

Finally, if $x \in U(n)_{k}$ and $y \varepsilon \tilde{A}_{(k)}$ then $\alpha y(x)=\ell((x \cdot(y \otimes 1)) w)=\ell((x \cdot y) w)=(x \cdot y) w$ since $x \cdot y \varepsilon \mathbb{C} l$. However, $x \cdot y=x^{\circ} \tau(k)^{\mathrm{y}}$ by 2.2.4. Thus, the pairing of $\tilde{A}_{(k)}$ with $U_{(n)_{k}}$ is equivalent to the analogous pairing of $\mathrm{A}_{(\mathrm{k})}$ with $U_{(n)_{\mathrm{k}}}$. However $\mathrm{A}_{(\mathrm{k})}$ is isomorphic to $A(N)_{k}$ as an $n$-module by (1.6.5). But now, $A(N)_{k}$ is nonsingularly priced to $U(n)_{k}$ by (1.6.3). Q.E.D.\\
4.2. Tensor products and Whittaker vectors

As a consequence of Theorem 4.1 we have the following theorem. The proof is somewhat simpler than the corresponding proof in Kostant's paper [9] and the assumptions are much less stringent.

Theorem 4.2. Let $V$ and $L$ be modules for the enveloping algebra $U$ of a complex semi-simple Lie algebra. Assume that $L$ is finite dimensional and let $T=V \otimes L$ be the tensor product module. Let Wh(V), Wh(T) denote the space of Whittaker vectors in $V$ and $T$ respectively. Then, there is a natural linear isomorphism\\
(4.2.1) $\quad \mathrm{Wh}(\mathrm{T}) \rightarrow \mathrm{Wh}(\mathrm{V}) \otimes_{\mathbb{C}} \mathrm{L}$\\
so that if $\mathrm{Wh}(\mathrm{V})$ is finite dimensional, then so is $\mathrm{Wh}(\mathrm{T})$ and


\begin{equation*}
\operatorname{dim} W h(T)=(\operatorname{dim} W h(V))(\operatorname{dim} L) . \tag{4.2.2}
\end{equation*}


Proof. For any $0 \neq z \varepsilon L$, let $L_{z}=U(n) z$ so that $L_{z}$ is a cyclic $U(n)$-module. In fact, if $J_{z}=\{u \varepsilon U(n) \mid u z=0\}$, then $L_{z} \cong U(n) / J_{z}$. Let $L_{z}^{\prime}$ be the contragradient module. Now, if $\left\{\rho_{i}\right\}$ is a basis for $L_{z}$ and $\left\{\rho_{i}^{\prime}\right\}$ is the dual basis of $L_{z}^{\prime}$, then the element $\Sigma \rho_{i}^{\prime} \otimes \rho_{i}$ of $L_{z}^{\prime} \otimes L_{z}$ is a $U(n)$-invariant. But now, since $J_{z}$ clearly contains a power of the augmentation ideal $n U(x)$, we can take $L_{z}^{\prime}$ to be the orthocomplement of $J_{z}$ in $U(n)^{*}$. Now, let $w \varepsilon W h V$. Then $w$ is a Whittaker module and Theorem 4.1 holds. Thus, let $r_{i} w$ correspond to $\rho_{i}^{\prime} \otimes w$ in $U(n)^{*} \otimes W h(w)$ under (4.1.3) where $r_{i} \varepsilon \tilde{A}$. Then $w_{z}=\sum r_{i} w \otimes \rho_{i} \varepsilon V \otimes L=T$ is an $n$-invariant where $n$ operates by the $\eta$-reduced action on $V$. Thus, $w_{z} \varepsilon$ WhT. Since $L_{z}$ is cyclic, $L_{z}=C z \oplus L_{z}$. Choose the basis $\left\{\rho_{i}\right\}$ so that $z=\rho_{1}$ and $\rho_{i} \varepsilon \eta U(n) z$ for all $i>1$. Now, $0=x \cdot w_{z}=\Sigma\left(\left(x \cdot r_{i} w\right) \otimes \rho_{i}+r_{i} w \otimes x \rho_{i}\right)$ clearly implies that $x \cdot(r, x)=0$ for any $x \in \lambda$. However,\\
$x \cdot(r, w)=\left(x \cdot r_{1}\right) w$ by (2.3.1). But this implies that $x \cdot r_{1}=0$ for all $x \varepsilon n$ so that $r_{1} \varepsilon \tilde{A} \cap U(\bar{p})^{N}=\mathbb{C} 1$.

Thus, $r_{1}$ is a nonzero constant (since $z \neq 0$ ) which we may take to be unity. Now, $w_{z}-w \otimes z \varepsilon v \otimes n \|(x) z$. Choose a basis for $L,\left\{z_{i}\right\}$, such that if $L_{j}=\sum_{i=1}^{j} C z_{i}$ then $L_{j} \subseteq L_{j-1}$. But then, if $z=z_{j}$, it follows that


\begin{equation*}
w_{z_{j}}-w \otimes z_{j} \varepsilon \tilde{A} w \otimes L_{j-1} . \tag{4.2.3}
\end{equation*}


But then the elements $w_{z_{j}}, w \varepsilon W h(V)$, are clearly linearly independent. In particular, we have an injection\\
$W h(V) \otimes_{C} L \rightarrow W h(T)$.\\
Conversely, let $\left\{z_{i}\right\}$ be a basis for $L$ as above. If $\mathrm{w} \varepsilon$ WhT, we may write $\mathrm{w}=\Sigma \mathrm{v}_{\mathrm{i}} \otimes \mathrm{z}_{\mathrm{i}}$ where $\mathrm{v}_{\mathrm{i}} \varepsilon \mathrm{V}$. Now, take I maximal such that $\mathrm{v}_{\mathrm{I}} \neq 0$. Then, as above, we must have that $x \cdot v_{I}=0$ for all $x \in \Omega$ and hence $\mathrm{v}_{I}$ is a Whittaker vector. However, we showed above that there is an element $\mathrm{w}_{Z_{I}} \varepsilon \mathrm{~Wh}(T)$ such that\\
$w_{z_{I}}-v_{I} \otimes z_{I} \varepsilon \tilde{A} v_{I} \otimes L_{I-1}$. Thus, $w-w_{z_{I}} \varepsilon V \otimes L_{I-1}$ and $w-w_{z_{I}} \varepsilon \mathrm{~Wh}(T)$. Proceeding by induction, it is clear\\
that we may write $w$ as a linear combination of the $w_{z_{i}}$\\
for $w \varepsilon$ WhV as above. Thus (4.2.1) is a bijection and (4.2.2) follows trivially.\\
Q.E.D.

Thus, Whittaker vectors are extremely well behaved with respect to tensor products. We will see in the next two chapters that their behavior actually improves if one takes the tensor product of $U(p)$-induced modules and finite dimensional modules.\\
4.3. n-finite modules

The remainder of the material in this chapter is borrowed, with few changes, from [Kostant].

Lemma 4.3. Let V be a Whittaker module for $U$. Then $H^{i}(n, V)=0$ for $i>0$ where $H(n, V)$ denotes the cohomology of $n$ on $V$ with respect to the $n$-reduced action.

Proof. Recall that we have the isomorphism (4.1.1) $\mathrm{V} \rightarrow U(n)^{*} \otimes \mathrm{Whv}$ of $n$-modules where $n$ acts by the $\eta$-reduced action on $V$ and by the left translation action on $U(n)^{*} \otimes W h$. Since under this action $n$ acts trivially on WhV, we have $\left.H(n, V) \cong H(n, U(n))^{*} \otimes W h V\right) \cong H\left(n, U(x)^{*}\right) \otimes W h V$. On the other hand $U(n)^{*}$ and $A(N)$ (the affine algebra of $N$ ) are isomorphic as ct-modules\\
where $x \in \Omega$ operates on $A(N)$ as the right invariant vector field $\xi_{x}$ defined in §1.6. Thus, to prove the lemma, it suffices to show that $H^{i}(\eta, A(N))=0$ for $i>0$ where $x$ operates as $\xi_{x}$ on $A(N)$. Let $(C(N), d)$ be the standard cochain complex for $H(, A(N))$ so that $\mathrm{C}(\mathrm{N})=\Lambda \eta^{\prime} \otimes \mathrm{A}(\mathrm{N})$ where $\eta^{\prime}$ is the dual space to $n$. If we regard $\Lambda n^{\prime}$ as the space of right invariant differential forms on N then $\mathrm{C}(\mathrm{N})$ is an $A(N)$ module of differential forms on $N$ which is stable under exterior differentiation $d_{*}$ and $d=d_{*} \mid C(N)$ as one easily sees. But since $N$ is unipotent, there exists a global coordinate system $\gamma_{i} \varepsilon A(N), i=1, \ldots, r$, on $N$ such that $S(N)$ is isomorphic to the polynomial algebra $\mathrm{C}\left[\gamma_{1}, \ldots, \gamma_{r}\right]$. Let $\left\{\mathrm{x}_{\mathrm{i}}\right\}$ be a basis for . Since $A(N)$ is an $n$-module, we may write $\xi \mathrm{x}_{\mathrm{j}}=\Sigma \mathrm{a}_{\mathrm{ij}} \frac{\partial}{\partial \gamma_{\mathrm{i}}}$ for $\mathrm{a}_{\mathrm{ij}} \varepsilon \mathrm{A}(\mathrm{N})$. Since any invertible element in $A(N)$ is necessarily a constant it follows that $n^{\prime} \subseteq \Sigma \mathrm{A}(\mathrm{N}) \mathrm{d} \gamma_{i} . \quad(\mathrm{C}(\mathrm{N}), \mathrm{d})$ is isomorphic to a Koszul complex. But this proves the lemma since one knows that a Koszul complex is acyclic.\\
Q.E.D.

Now let V be any $U$-module. Recall that $U_{\eta}(n) \subseteq U(n)$ is the kernel of $n$ as a homomorphism of\\
$U(n)$ Let


\begin{equation*}
v_{\eta}= \begin{cases}v & \left.\varepsilon V \mid\left(U_{\eta}(n)\right)^{k} v=0 \text { for some } k \varepsilon \mathbb{Z}_{+}\right\} .\end{cases} \tag{4.3.1}
\end{equation*}


If "dot" as usual denotes the $\eta$-reduced action, then note that $v \in V_{\eta}$ if and only if $P \cdot v=0$ where $P$ is some power of the augmentation ideal $n U(n)$.

Proposition 4.3.1. $V_{\eta}$ is a $\hat{U}$-submodule of $V$.

Proof. If $\eta_{0}$ is any character of $n$, let $V_{\eta_{0}}$ be defined in an analogous manner. It is convenient to think of $V_{\eta}$ as the $\eta$-weight space for the weight $\eta_{0}$. But now regarding $U$ as a $n$-module with respect to the adjoint representation, one has $U=U_{n_{0}}$ where $n_{0}$ is the trivial character. On the other hand, if $x \varepsilon n_{0} u \varepsilon U_{\text {, }}$ and $v \in v_{\eta}$ one has


\begin{equation*}
x \cdot(u v)=[x, u] v+u x \cdot v . \tag{4.3.2}
\end{equation*}


But then the usual argument on the additivity of weights shows that $U v_{n} \subseteq v_{n}$.\\
Q.E.D.

Remark 4.3. It is clear that $\mathrm{V}=\mathrm{V}_{\mathrm{\eta}}$ whenever V\\
is a Whittaker module since we have an isomorphism $\tilde{A} W h \mathrm{~V}=\mathrm{V}$.

Proposition 4.3.2. If $\mathrm{v} \varepsilon \mathrm{v}$, then


\begin{equation*}
U(n) \cdot v=U(n) v . \tag{4.3.3}
\end{equation*}


Furthermore, if $v \in v_{n}$, then $U(n) v$ is finite dimensional. If in addition $v \neq 0$, then $U(n) v$ contains a non-zero Whittaker vector.

Proof. The map $n \rightarrow U(n), x \rightarrow x-n(x) I$ clearly extends to an automorphism $\sigma: U(n) \rightarrow U(n)$. Furthermore, one clearly has $u \cdot v=\sigma(u) v$. This proves (4.3.3). Now, if $v \varepsilon V_{n}$, then $P \cdot v=0$ where $P=(U(n))^{k}$ for some $k \in \mathbb{Z}_{+}$. Thus, $U(n) v$ is finite dimensional by (4.3.3). Furthermore, if $\mathrm{v} \neq 0$, then by Engle's theorem there exists $0 \neq w \varepsilon U(n) v$ such that $x \cdot w=0$ for all $x \varepsilon n$. But then $w$ is a Whittaker vector.\\
Q.E.D.

A composition series for $V$ is a finite sequence of $U$-submodules $V_{i}, i=0, \ldots, k$, such that $V_{0}=0$, $V_{k}=V, V_{i-1}-V_{i}$ and $V_{i} / V_{i-1}$ is an irreducible U-module. If such a sequence exists, we say that $V$ has a composition series. One knows that $k$ is unique\\
and is called the length of $V$. Also, we may speak of the multiplicity of any irreducible representation of in $V$.

Now, a $U$-module $V$ is said to be $\eta$-finite if 1) $V=V_{\eta}$ and 2) $V$ has a composition series.

Proposition 4.3.3. Assume $V$ has a composition series, $V_{i}, i=0, \ldots, k$. Then $V$ is $\eta$-finite if and only if $V_{i} / V_{i-1}, i=1, \ldots, k$, are irreducible Whittaker modules.

Proof. If each $V_{i} / V_{i-l}$ is a Whittaker module, then $\mathrm{V}=\mathrm{V}_{\mathrm{n}}$ by Remark 4.3. Conversely, assume $\mathrm{V}=\mathrm{V}_{\mathrm{n}}$. But then $\left(V_{i} / V_{i-1}\right)_{\eta}=V_{i} / V_{i-1}$. Thus $V_{i} / V_{i-1}$ contains a non-zero Whittaker vector by Proposition 4.3.2. Thus $V_{i} / V_{i-1}$ is an irreducible Whittaker module since any irreducible $U$-module containing a Whittaker vector is a Whittaker module by Theorem 3.4.\\
Q.E.D.

Theorem 4.3. Assume that $V$ is an $\eta$-finite U-module. Then $H^{i}(\pi, V)=0$ for $i>0$ with respect to che $\eta$-reduced action of $\eta$ on $V$.

Proof. Let $V_{j}, j=0, \ldots, k$, be a composition series\\
for $V$ so that for $j \geq 1$ one has an exact sequence\\
(4.3.5) $\quad 0 \rightarrow v_{j-1} \rightarrow v_{j} \rightarrow v_{j} / v_{j-1} \rightarrow 0$\\
of $U$-modules. Now if $H^{i}\left(\lambda, V_{j-1}\right)=0$ for some $j \geq 1$,\\
then $H^{i}\left(\Omega, V_{j}\right)=0$ by Proposition 4.3.3, Lemma 4.3 and\\
general facts about cohomology sequences associated to an\\
exact sequence. However, $H^{i}\left(n, V_{1}\right)=0$ by Lemma 4.2.\\
This inductively proves the theorem.\\
Q.E.D.\\
4.4. Composition series and Whittaker vectors

In the case where $G=s \ell(n ; C)$, the exactness\\
of (4.4.3) was first proved by Casselman and Zuckerman\\
for the Borel subalgebra.

Theorem 4.4. Let $V$ be any $\eta$-finite $U$-module and\\
let $W h \mathrm{~V}$ denote the space of Whittaker vectors in $V$.\\
Let $V_{i}, i=0, \ldots, k$, be a composition series for $V$\\
and assume that $\operatorname{dim} W h\left(V_{i} / V_{i-1}\right)=p_{i}<\infty$ for $i=1, \ldots, k$.\\
Then WhV is finite dimensional and in fact\\
(4.4.1) $\quad \operatorname{dim} W h(V)=\sum_{i=1}^{k} \operatorname{dim} W h\left(V_{i} / V_{i-1}\right)$.

Furthermore, if $V$ is any $\eta$-finite U-module and


\begin{equation*}
0 \rightarrow v^{1} \rightarrow v \rightarrow v^{2} \rightarrow 0 \tag{4.4.2}
\end{equation*}


is an exact sequence of U-modules, then V is $\eta$-finite if and only if both $V^{1}$ and $V^{2}$ are $\eta$-finite. Furthermore, in such a case (4.4.2) induces an exact sequence

$$
\text { (4.4.3) } 0 \rightarrow W h V^{1} \rightarrow W h V \rightarrow W h V^{2} \rightarrow 0
$$

so that


\begin{equation*}
\operatorname{dim} W h V=\operatorname{dim} W h V^{1}+\operatorname{dim} W h V^{2} \tag{4.4.4}
\end{equation*}


provided that all terms are finite.

Proof. It is clear from the definitions that $V$ is $\eta$-finite if and only if $\mathrm{V}^{1}$ and $\mathrm{V}^{2}$ are $\eta$-finite in (4.4.2). But then recalling the cohomology exact sequence associated to an exact sequence one has

$$
0 \rightarrow H^{0}\left(n, v^{2}\right) \rightarrow H^{0}(n, v) \rightarrow H^{0}\left(n, v^{1}\right) \rightarrow 0
$$

with respect to the $\eta$-reduced action since $H^{\prime}\left(\eta, V^{1}\right)=0$\\
by Theorem 4.3. However, it is clear that $H^{0}\left(\pi, V^{l}\right)=W h V^{l}$. Thus, (4.4.3) and (4.4.4) follow. Finally, it follows from (4.3.5) and (4.4.3) that $\operatorname{dim} W h\left(V_{j-1}\right)+ \operatorname{dim} W h\left(V_{j} / V_{j-1}\right)=\operatorname{dim} W h\left(V_{j}\right)$. Thus, one has (4.4.1) by induction.\\
Q.E.D.\\
4.5. Submodules generated by Whittaker vectors

Let $V$ be a U-module and let WhV denote the space of Whittaker vectors in $V$. Then, it is clear that WhV is a $U(\bar{p})^{\mathrm{N}}$-module so it makes sense to talk about composition series for $U(\bar{p})^{N}$-submodules of WhV.

Lemma 4.5. Let $V$ be any $U$-module. Assume that $W \subseteq W h V$ is a $U(\bar{p})^{N}$-stable subspace. Then $W$ is $\eta$-finite if and only if $W$ has a (finite) composition series. Furthermore, in such a case


\begin{equation*}
w=w h(U w) . \tag{4.5.1}
\end{equation*}


Proof. Suppose that $W$ has a finite composition series, $W_{i}, i=0, \ldots, k$. If $k=1$, then $W$ is an irreducible $U(\bar{p})^{N}$-module and $W$ is an irreducible Whittaker module by Theorem 3.4. It is clear that (4.5.1) holds in this case. Assume $k>1$ and assume inductively\\
that the result is true for smaller values of the length of the composition series. Then, we have $W \subseteq W_{k-1}$ such that $W_{k-1}$ is $U(\bar{\rho})^{N}$-stable. But then, by induction we have $W_{k-1}$ is $\eta$-finite and $W_{k-1}=W h\left(l\left(W_{k-1}\right)\right.$. On the other hand, $W / W_{k-1}$ is an irreducible $U(\bar{\rho})^{N}-$ module and it follows from Theorem 3.4 that $U\left(W / W_{k-1}\right)=U w / d W_{k-1}$ is irreducible. Thus, $U w$ is $\eta$-finite and $W h(U W)=W$ by (4.4.3).

Conversely, suppose that $W$ doesn't have a finite composition series. Then any $U(\bar{p})^{N}$-module containing $w$ cannot have a finite composition series. In particular, since $W h(U W) \subseteq W, W h(U W)$ cannot have a finite composition series. Thus, $U W$ cannot have a finite composition series.

\section{Q.E.D.}
Proposition 4.5. Let $V$ be any $U$-module. Then $V_{\eta}$ is $\eta$-finite if and only if $W=W h(V) \subseteq V_{\eta}$ has a finite composition series. Furthermore, in such a case


\begin{equation*}
v_{n}=U w . \tag{4.5.3}
\end{equation*}


Moreover, the length of the composition series of $V_{\eta}$ is the same as the length of the composition series of $W$.

Proof. With the exception of (4.5.3), the statements in this proposition follow easily from Lemma 4.5. It remains only to show that $V_{\eta}=W$. Consider the quotient module $R=V_{\eta} /$ W. Clearly $R=R_{\eta}$. If $R_{\eta} \neq 0$, Proposition 4.3.2 asserts that $R$ contains a non-zero Whittaker vector. This means that there is a $v \varepsilon v$ such that $x \cdot v \varepsilon \quad w$ for all $x \varepsilon \wedge$. But if $\pi: \pi \rightarrow W$ is the map defined by $\pi(x)=x \cdot v$ then $\pi$ is a l-cocycle for $n$ with values in $W$. However, $H^{\prime}(\pi, W)=0$ by Theorem 4.3. Thus, there exists $u \varepsilon \quad W$ such that $\mathbf{x} \cdot \mathbf{v}=\mathbf{x} \cdot \mathbf{u}$ for all $\mathbf{x} \varepsilon \lambda$. But then $\mathbf{v}-\mathbf{u}$ is a non-zero Whittaker vector. Thus, $v-u \in W \subseteq U W$ and hence $v \varepsilon$ Uw. This is a contradiction. Thus $R=0$ and $V_{\eta}=W$.

Corollary 4.5. Let $V$ be any $U$-module. Let $V_{\eta}$ be the $U$-submodule defined by (4.3.1) and let WhV be the space of Whittaker vectors in $V$ so that WhV is a $U(\bar{p})^{N}$-module. Let $W=W h V$. then $V_{n}$ is a completely reducible $U$-module if and only if $W$ is a completely reducible $U(\bar{p})^{N}$-module. Furthermore, in such a case we have


\begin{equation*}
w=v_{n}=\sum_{i=1}^{k} l w_{i} \tag{4.5.4}
\end{equation*}


is a direct sum of irreducible $U$-modules.

Proof. This follows immediately from Theorem 4.4. Q.E.D.\\
4.6 A formula for $\operatorname{dim}$ Wh(M')

Assume that $R$ is a commutative Noetherian algebra over $C$. Let $J \subseteq R$ be an ideal. The Artin-Rees theorem implies that if M is a Noetherian module over $R$ and $\bar{M} \subseteq M$ is a submodule, then for any $k \in \mathbb{Z}_{+}$, there exists $i \varepsilon \mathbb{Z}_{+}$such that


\begin{equation*}
\mathrm{J}^{\mathrm{i}} \mathrm{M} \cap \overline{\mathrm{M}} \subseteq \mathrm{~J}^{\mathrm{k}} \overline{\mathrm{M}} . \tag{4.6.1}
\end{equation*}


$N$. Wallach has shown that this result holds if $R$ is replaced by the enveloping algebra of a nilpotent Lie algebra.

Proposition 4.6. Let $M$ be any finitely generated $U(n)$-module. Put $J=\left(U_{n}(n)\right)$. Then for any submodule $\overline{\mathrm{M}} \subseteq \mathrm{M}$ and for any $\mathrm{k} \varepsilon \mathrm{Z}_{+}$there exists $\mathrm{i} \varepsilon \mathrm{Z}_{+}$such that one has the inclusion (4.6.1).

Proof. See Kostant [9] or Wallach [17]. Q.E.D.

Now assume that


\begin{equation*}
0 \leftarrow \overline{\overline{\mathrm{M}}} \leftarrow \mathrm{M} \leftarrow \overline{\mathrm{M}} \leftarrow 0 \tag{4.6.2}
\end{equation*}


is an exact sequence of $U$-modules. Let $V=M^{\prime}$, the full dual space to $M$, regarded contragradiently as a U-module. Let $\overline{\bar{V}} \subseteq V$ be the orthocomplement of $\bar{M}$ in V and put $\overline{\mathrm{V}}=\mathrm{V} / \overline{\overline{\mathrm{V}}}$ so that one has an exact sequence


\begin{equation*}
0 \rightarrow \overline{\bar{V}} \rightarrow V \rightarrow \bar{V} \rightarrow 0 \tag{4.6.3}
\end{equation*}


of $U$-modules where we can regard $\overline{\overline{\mathrm{V}}}=(\overline{\overline{\mathrm{M}}})^{\prime}, \overline{\mathrm{V}}=(\overline{\mathrm{M}})^{\prime}$.\\
The argument in the following lemma is due to $N$. Wallach in the case when $U_{\eta}(n)$ is replaced by the augmentation ideal $\bigcup(n)$.

Lemma 4.6. Assume that the $U$-module $M$ is finitely generated over $U(N)$. Using the notation of (4.3.1), the exact sequence (4.6.3) induces an exact sequence


\begin{equation*}
0 \rightarrow \overline{\bar{v}}_{\eta} \rightarrow v_{\eta} \rightarrow \bar{v}_{\eta} \rightarrow 0 \tag{4.6.4}
\end{equation*}


of $n$-finite U-modules.

Proof. Let $J \subseteq U(N)$ be the ideal defined by $J=\left(U_{n}(n)\right)$ where we recall that $U_{n}(\pi)=\operatorname{ker} \eta$.

Then, since for $u \varepsilon U(n), m \varepsilon M, v \varepsilon V$ one has <um, v> = <m, uv>, it follows that


\begin{equation*}
(\mathrm{JM})^{\circ}=\mathrm{WhV} \tag{4.6.5}
\end{equation*}


where the superscript o denotes orthocomplement in V and that furthermore


\begin{equation*}
\sum_{k=0}^{\infty}\left(J^{k} M\right)^{\circ}=V_{\eta} . \tag{4.6.6}
\end{equation*}


But now, since $J$ has codimension 1 in $U(M)$ and $M$ is $U(n)$-Noetherian (i.e. finitely generated) it is obvious that WhV is finite dimensional. But then, WhV clearly has a finite composition series. By Theorem 4.4, $V_{\eta}$ is $\eta$-finite. The same argument applies to $\bar{V}_{\eta}$ and $\overline{\overline{\mathrm{V}}}_{\eta}$ since $\overline{\mathrm{M}}$ and $\overline{\overline{\mathrm{M}}}$ are also Noetherian (since $U(n)$ is a Noetherian algebra.)

But now the exact sequence (4.6.3) obviously induces the exact sequence $0 \rightarrow \overline{\overline{\mathrm{~V}}}_{\eta} \rightarrow \mathrm{V}_{\eta} \rightarrow \overline{\mathrm{V}}_{\eta}$. The only problem is to show that the map $V_{\eta} \rightarrow \bar{V}_{\eta}$ is surjective. Let $\overline{\mathrm{v}} \varepsilon \overline{\mathrm{V}}_{\mathrm{n}}$. Then there exists $\mathrm{k} \varepsilon \mathrm{Z}_{+}$such that $\overline{\mathrm{v}}$ vanishes on $\mathrm{J}^{\mathrm{k}} \overline{\mathrm{M}}$. However, by (4.6.1) there exists $\mathrm{i} \varepsilon \mathrm{Z}_{+}$such that $\bar{v}$ vanishes on $\mathrm{J}^{i} \mathrm{M} \quad \overline{\mathrm{M}}$. But then it is clear\\
that $\overline{\mathrm{V}}$ extends to a linear functional $\mathrm{V} \varepsilon \mathrm{V}$ which vanishes on $\mathrm{J}^{\mathrm{i}} \mathrm{M}$. But then $\mathrm{v} \varepsilon \mathrm{V}_{\mathrm{n}}$ and v induces $\overline{\mathrm{v}}$ on $\bar{M}$. This proves $V_{\eta} \rightarrow \bar{V}_{\eta}$ is surjective.

\section{Q.E.D.}
In the case of a Borel subalgebra, $U(\bar{\rho})^{N}$ is abelian and Lemma 4.6 is quite sufficient for the Whittaker theory in this case. However, in general, $U(\bar{p})^{\mathrm{N}}$ has irreducible infinite-dimensional representations and one would like to replace the assumption that $M$ is finitely generated over $U(n)$ with some condition that utilizes the fact that $V_{\eta}$ is $\eta$-finite if and only if WhV has a finite composition series.

Now, assume that $M$ is a $U$-module and one has a reverse finite filtration


\begin{equation*}
M=M_{0} \supseteq \cdots \geq M_{i-1} \geq M_{i} \supseteq \cdots \geq M_{k}=0 \tag{4.6.7}
\end{equation*}


of $U$-submodules. Let $V$ be the full dual to $M^{\prime}$ and let $V_{i} \subseteq V$ be the orthocomplement to $M_{i}$ so that


\begin{equation*}
0=\mathrm{v}_{0} \subseteq \cdots \subseteq \mathrm{v}_{\mathrm{i}-1} \subseteq \mathrm{v}_{\mathrm{i}} \subseteq \cdots \subseteq \mathrm{v}_{\mathrm{k}}=\mathrm{v} \tag{4.6.8}
\end{equation*}


is a filtration of $V$ by $U$-submodules. Clearly


\begin{equation*}
\left(M_{i-1} / M_{i}\right)^{\prime}=V_{i} / V_{i-1} . \tag{4.6.9}
\end{equation*}


Theorem 4.6. Let the notation be as above. Assume now that the $U$-module $M$ is finitely generated over $U(n)$. Then for any $i=1, \ldots, k$, the quotient map defines an exact sequence


\begin{equation*}
0 \rightarrow\left(V_{i-1}\right)_{\eta} \rightarrow\left(V_{i}\right)_{\eta} \rightarrow\left(V_{i} / V_{i-1}\right)_{\eta} \rightarrow 0 \tag{4.6.10}
\end{equation*}


of $\eta$-finite modules. Moreover this induces an exact sequence

$$
\text { (4.6.11) } 0 \rightarrow W h\left(V_{i-1}\right) \rightarrow W h\left(V_{i}\right) \rightarrow W h\left(V_{i} / V_{i-1}\right) \rightarrow 0
$$

and one has the formula


\begin{equation*}
\operatorname{dim} W h(V)=\sum_{i=1}^{k} \operatorname{dim} W h\left(V_{i} / V_{i-1}\right) . \tag{4.6.12}
\end{equation*}


Proof. Since $M$ is finitely generated, it follows that $M / M_{i}$ is finitely generated over $U(x)$ for all $i$. But then (4.6.10) is an exact sequence of $\eta$-finite U-modules by (4.6.9) and (4.6.4). But then by (4.4.3) and (4.4.4) one has (4.6.11) and the relation\\
$\operatorname{dim} W h\left(V_{i}\right)=\operatorname{dim} W h\left(V_{i-1}\right)+\operatorname{dim} W h\left(V_{i} / V_{i-1}\right)$. The\\
formula (4.6.12) follows immediately by induction.\\
Q.E.D.

\section{Chapter 5}
\section{Whittaker vectors and Verma modules}
The material in this chapter has been treated by Kostant [9, Chapter 3] when $\rho$ is a Borel subalgebra.\\
5.1. The completions of generalized Verma modules.

In this chapter, $O$ will denote a complex, semisimple Lie algebra and $p=m \oplus n$ will denote the decomposition of a (standard) parabolic subalgebra of $y$ into reductive and nilpotent components. The element $x_{0} \varepsilon a_{0} \subseteq m$ and the sufspaces $d_{i}$ (eigenspaces for ad $x_{0}$ in (7.) will be as in §1.1.

Let H be an irreducible, finite dimensional left $U(m)$-module. We let $H$ be considered as a left $U(p)$-module by letting $n$ act trivially. The generalized Verma module $\mathrm{V}=\mathrm{V}_{\mathrm{H}}$ is then defined by


\begin{equation*}
v=U_{\otimes_{U(p)}}{ }^{H} . \tag{5.1.1}
\end{equation*}


Clearly, $V$ is a $U=U(\sigma)$-module under left multiplication.\\
Recall that $m=\mathrm{d}_{0}=\sigma \gamma^{\mathrm{x}_{0}}$. Thus, $\mathrm{x}_{0} \mathrm{mv}=\mathrm{mx}_{0} \mathrm{v}$ for all m ɛM, v ɛ H. Since H is irreducible, Schur's lemma\\
implies that there is a constant $\lambda \varepsilon \mathbb{C}$ such that $\mathrm{x}_{0} \mathrm{v}=\lambda \mathrm{v}$ for all $\mathrm{v} \varepsilon \mathrm{H}$.

Recall that ad $\mathrm{x}_{0}$ restricted to $\lambda$ (resp. $\bar{n}$ ) has strictly positive (resp. negative) integral eigenvalues so that we may define gradations of $U(n)$ and $U(\bar{n})$ such that for $i \varepsilon \mathrm{Z}_{+}$


\begin{align*}
& U(n)_{i}=\left\{u \varepsilon U(n) \mid \operatorname{ad} x_{0}(u)=i u\right\}  \tag{5.1.2}\\
& U(\bar{n})_{i}=\left\{u \varepsilon U(\bar{n}) \mid \operatorname{ad} x_{0}(u)=-i u\right\} .
\end{align*}


Then if $\mathrm{v} \varepsilon \mathrm{V}$ is such that $\mathrm{x}_{0} \mathrm{v}=\mu \mathrm{v}, \mu \varepsilon \mathrm{C}$, then $x_{0} u v=(\mu+i) v$ depending on whether $u \varepsilon U_{(n)_{i} \text { or }}$ u $\varepsilon \bigcup(\bar{n})_{i}$. Thus, if we define


\begin{equation*}
V(\mu)=\left\{v \varepsilon V \mid x_{0} v=\mu v, \mu \varepsilon C\right\} \tag{5.1.3}
\end{equation*}


then it is clear that


\begin{equation*}
V=\underset{i \varepsilon Z_{+}}{\oplus} V(\lambda-i) \tag{5.1.4}
\end{equation*}


is a gradation of $V$.\\
Now, $U(n)$ is an $n$-module under left multiplication. If $U(n)$ ' is the fuall dual to $U(n)$, then $U(n)$ ' inherits\\
an $\lambda$-module structure by contragradience. Thus, if $\mathbf{u} \varepsilon U(n), \beta \varepsilon U(n)^{\prime}$ and $\mathbf{x} \varepsilon n$, then $\langle\mathbf{x} \beta, \mathbf{u}\rangle=-\langle\beta, \mathbf{x u}\rangle$. As usual, $U(n)^{*}$ will denote the $n$-submodule of $U(n)$ ' consisting of those elements which vanish on a power of the augmentation ideal $n U(\pi)$.

Now, let $\tilde{V}=\underset{i>0}{\oplus} V(\lambda-i)$ so that $V=\tilde{V} \oplus H$ is a direct sum. Let $\ell: V \rightarrow H$ be the projection onto the second summand under this decomposition. For any $v \in v$, we let $\beta_{v} \varepsilon U(n)^{\prime} \otimes H$ (linear functionals on $U(n)$ with values in H$)$ be defined by


\begin{equation*}
\beta_{v}(u)=\ell(\check{u} v) \tag{5.1.5}
\end{equation*}


where the linear isomorphism $U(n) \rightarrow U(n), u \rightarrow \stackrel{v}{u}$ was defined in §4.1.

Proposition 5.1. The correspondence $\mathrm{v} \rightarrow \beta_{\mathrm{v}}$ defines a map


\begin{equation*}
\mathrm{V} \rightarrow U(n)^{*} \otimes \mathrm{H} \tag{5.1.6}
\end{equation*}


of $n$-modules. Furthermore, (5.1.6) is an isomorphism of $n$-modules if and only if $V$ is irreducible.

Proof. If $v \varepsilon v$, it follows easily from (5.1.4) that $\beta_{v}$ vanishes on a power of $\varkappa U(n)$ and hence one has a map (5.1.6). Now, if $\mathbf{x} \varepsilon \eta, \mathbf{v} \varepsilon \mathbf{v}$ and $\mathbf{u} \varepsilon U(n)$, $\beta_{x v}(u)=\ell(u x v)=-\ell\left((x u) v=-\beta_{v}(x u)=(x \beta)(u)\right.$. Thus, $\beta_{\mathbf{x V}}=\mathbf{x} \beta_{\mathbf{V}}$ so that (5.1.6) is a map of $n$-modules.

By contragradience, (5.1.2) gives a gradation of $U(n)^{*}$ so that $U(n)^{*}=\underset{i \varepsilon Z_{+}}{\oplus} U(n)_{-i}^{*}$. Now, if $\mathrm{v} \varepsilon \mathrm{v}(\lambda-i)$ and $u \in U(n)_{j}$, then ùv $\varepsilon V(\lambda-i+j)$ so that $\ell(\check{u} v)=0$ unless $i=j$. Thus, (5.1.6) restricts to a map $V(\lambda-i) \rightarrow U(n)_{-i}^{*} \otimes H$. Thus, (5.1.6) is an isomorphism of $\lambda$-modules if and only if $V(\lambda-i) \rightarrow U(\lambda)_{-i}^{*} \otimes H$ is a linear isomorphism for all $\mathrm{i} \varepsilon \mathrm{Z}_{+}$.

Now, if $\mathrm{V}_{1} \subseteq \mathrm{~V}$ is a proper $n$-submodule, then $\mathrm{V}_{1}$ is spanned by ad $x_{0}$-eigenvectors. However, since $H$ is $U(m)$-irreducible, we must have that $(\mathrm{H}=\mathrm{V}(\lambda)) \cap \mathrm{V}_{1}=\{0\}$. Thus, $\beta_{\mathrm{V}}=0$ for all $\mathrm{V} \varepsilon \mathrm{V}_{1}$ so that (5.1.6) is not an isomorphism.

Conversely, suppose that $V$ is irreducible. If $0 \neq \mathrm{v} \varepsilon \mathrm{V}(\lambda-\mathrm{i})$ and $o \neq \mathrm{v}_{0} \varepsilon \mathrm{H}$, there exists $\mathrm{u} \varepsilon$ U such that $u v=v_{0}$. By writing $u$ in a Birkhoff-Witt basis with respect to $U=U(\bar{n}) \otimes U(m) \otimes U(n)$, a simple consideration of $\mathrm{x}_{0}$-weights shows that there exists an element $u^{\prime} \varepsilon U(n)_{i}$ such that $0 \neq v_{1}=u^{\prime} v \varepsilon H=V(\lambda)$.

Thus, (5.1.6) is injective. However, (5.1.6) is also onto since $\operatorname{dim} V(\lambda-i)=\left(\operatorname{dim} \bigcup(\bar{n})_{i}\right)\left(\operatorname{dim} V_{\lambda}\right)=$ ( $\operatorname{dim} U\left(N_{i}^{*}\right)(\operatorname{dim} H)$.\\
Q.E.D.

It would be very convenient if Proposition 5.1 were to hold under more general conditions. In particular, we have

Conjecture 5.1. Let $H$ be an irreducible, infinitedimensional representation of $U_{(m)}$. Let $V=U_{\Theta_{U(p)}}{ }^{H}$ as in (5.1.1). Let $\mathrm{V} \rightarrow U(n)^{*} \otimes \mathrm{H}$ be defined as in (5.1.6). Then this map is surjective if $V$ is irreducible.

If Conjecture 5.1 holds, then Proposition 5.1 holds since the remainder of the proof carries through directly. We will assume that $H$ is finite dimensional in the remainder of this chapter, but we will phrase the proofs so that they hold if $H$ is infinite dimensional and Conjecture 5.1 holds.

Recalling the direct sum decomposition (5.1.4) of $V$, we now define the completion $\bar{V}$ of $V$ to be the direct product\\
(5.1.7)

$$
\bar{V}=\prod_{i \varepsilon Z_{+}} V(\lambda-i) .
$$

Thus, an element $v \varepsilon \bar{V}$ is a formal infinite sum


\begin{equation*}
v=\sum_{i \varepsilon Z_{+}} v_{i} \tag{5.1.8}
\end{equation*}


where $v_{i} \varepsilon V(\lambda-i)$. We now observe that the U-module structure on $V$ extends to $\bar{V}$ so that if $u \in U$ is such that ad $x_{0}(u)=j u$, then $u v=y$ where $y_{i}=u v_{i-j}$. As a $U$-module, we refer to $\bar{V}$ as the completion module of V.

By considering ad $\mathrm{x}_{0}$-weights, it is easy to see that the generalized Verma module $V$ cannot contain a Whittaker vector. However, one has

Theorem 5.1. Let $\mathrm{V}=\mathrm{V}_{\mathrm{H}}$ be an irreducible generalized Verma module and let $\bar{V}$ be its completion. Suppose that there is an admissible nilpotent for $\rho$ in $d_{-1}$ and let $\eta: U(\eta) \rightarrow \mathbb{C}$ be the corresponding admissible homomorphism. See §2.2. Then, the space of Whittaker vectors in $\overline{\mathrm{V}}$, Wh $(\bar{V})$, is isomorphic to $H$ as a $U(\bar{\rho})^{N}$-module where $U(\bar{p})^{\mathrm{N}}$ acts on H via the generalized Harish-Chandra homomorphism. (See Corollary 2.3.2 and Theorem 3.2.)

Proof. Recall that the full dual $U(n)^{\prime}$ of $U(n)$ is an $n$-module. On the other hand, it is clear that $U(n)$ ' is the direct product of the spaces $U(n)_{-i}^{*}$ over\\
all $i \varepsilon Z_{+}$. It is also clear that the map (5.1.6) extends in an obvious way to a map


\begin{equation*}
\overline{\mathrm{V}} \rightarrow U(n)^{\prime} \otimes \mathrm{H} \tag{5.1.9}
\end{equation*}


and that by Proposition 5.1, (5.1.9) is a map of $\pi$-modules. But it is also clear that (5.1.9) is an isomorphism of $\eta$-modules in case $V$ is irreducible. Recalling the definition of $U_{\eta}(\eta)$ in §2.2, it is clear that Wh $\overline{\mathrm{V}}$ is just the preimage of the orthocomplement of $U_{\eta}(\eta)$ in $U(n)$ ' $H$. However, since $U_{\eta}(n)$ has codimension 1 in $\hat{U}(n)$ one has a map


\begin{equation*}
\mathrm{H} \rightarrow \mathrm{~Wh}(\overline{\mathrm{~V}}) \tag{5.1.10}
\end{equation*}


and this map is a linear isomorphism.\\
Now, let $w \varepsilon W h(\bar{V})$ and write $w=\sum_{i \varepsilon Z_{+}} w_{i}$ where $w_{i} \varepsilon V(\lambda-i)$. If $w \neq 0$, let $I$ be the least integer such that $w_{I} \neq 0$. By Proposition 5.1, if $I \neq 0$, then $\{0\} \neq \lambda \mathrm{w}_{I} \subseteq \mathrm{~V}(\lambda-\mathrm{I}+1)$. However, if w is a Whittaker vector and $x \in \lambda$, then $x w=\eta(x) w$ and it follows that $x w_{I}=0$. Thus, $I=0$. Finally, we note that if $u \varepsilon U(\bar{\rho})^{N}$ and if $u_{1}$ is the image of $u$ in $U(m)$\\
under the generalized Harish-Chandra homomorphism, then the component of $u w$ in $V(\lambda)$ is precisely $u_{1} w_{0}$. Q.E.D.\\
5.2. The annihilators of Whittaker modules.

Now, let $\mathrm{V}=\mathrm{V}_{\mathrm{H}}$ be a Verma module and let $\mathrm{V}_{1}$ be an irreducible $U$-submodule of $V$. We may now calculate the annihilator of any Whittaker submodule of $\overline{\mathrm{V}}_{1} \subseteq \overline{\mathrm{~V}}:$

Theorem 5.2. Let the notation be as above and let $w$ be a Whittaker submodule of $\bar{v}_{1}$. Let $U_{w}$ and $U_{v_{1}}$ denote the annihilators of $W$ and $V_{1}$ in $U$. Then


\begin{equation*}
u_{w}=u_{v_{1}} \tag{5.2.1}
\end{equation*}


Proof. Clearly, $U_{W} \subseteq U_{V_{1}}$ since $w \subseteq \bar{V}_{1}$. Now, since the annihilator of a module is stable under the adjoint action of $O$ and hence $\mathrm{x}_{0}$, it is spanned by ad $x_{0}$-eigenvectors. Thus, using the obvious notation, it is enough to show that $U_{V_{1}}{ }^{(i)} \supseteq U_{W^{(i)}}$ for all $i \varepsilon Z_{+}$. So, let $u \varepsilon U_{W}(i)$ and suppose that $u \notin U_{v_{1}}$ (i). Then there exists $z \varepsilon V_{1}$ such that $u z \neq 0$. It is clearly sufficient to assume that $z \varepsilon V_{1}(\lambda-j)$ for some\\
$j \varepsilon Z_{+}$. But now, if $v \varepsilon W$, we can write $v$ as an infinite "sum", $\mathrm{v}=\sum_{\mathrm{k} \varepsilon \mathrm{Z}_{+}} \mathrm{v}_{\mathrm{k}}$ where $\mathrm{v}_{\mathrm{k}} \varepsilon \mathrm{v}_{1}(\lambda-\mathrm{k})$. Furthermore, choosing $v \neq 0$, there exists $r \varepsilon \mathrm{Z}_{+}$ such that $v_{r} \neq 0$. By the irreducibility of $v_{1}$, there must exist $y \in \chi_{(j-r)}$ such that $u v_{r}=z$. Thus, yv $\varepsilon V_{1}$ and if we write $y v=\sum(y v)_{s^{\prime}}$ then $(y v)_{j}=z$. But then also $u y v=\sum(u y v)_{s}$ where $(u y v)_{j+i}=u z$. However, $u z \neq 0$ by assumption but $u y v=0$ since $u \varepsilon U_{w}$. This is a contradiction proving the theorem. Q.E.D.

Note that if Conjecture 5.1 holds, then the above theorem holds if we let $V_{1}$ be an irreducible submodule of a generalized Verma module $V=V_{H}$ where $H$ is an infinite-dimensional irreducible $U(m)$-module.

\subsection{Comments}
Let $\Delta\left(n, m_{A}\right)$ denote the set of roots for the pair $\left(n, m_{\mathrm{A}}\right)$. We select a fundamental system of roots for $\Delta\left(n, m_{\mathrm{A}}\right), \pi_{\mathrm{A}}$, and note that we can identify $\pi_{\mathrm{A}}$ with $\left\{\alpha_{i} \varepsilon \pi \mid \alpha_{i} \notin \Sigma\right\}$. See §1.1. Thus, $\pi_{A}$ is a basis for the dual of $m_{\mathrm{A}}$. We define a dual basis $\left\{\mathrm{a}_{\mathrm{i}} \varepsilon m_{\mathrm{A}} \mid \alpha_{\mathrm{i}} \varepsilon \Sigma\right\}$ so that $\alpha_{i}\left(a_{j}\right)=\delta_{i j}$ for all $\alpha_{i} \varepsilon \pi$.

We have the following proposition.

Proposition 5.3. Let $H$ be an irreducible finitedimensional $U(m)$-module as above and let v be the associated character on $U\left(m_{A}\right)$. Let $\rho=\frac{1}{2} \sum_{\phi \varepsilon \Delta}{ }^{\phi}$. Then, if $(v+\rho)\left(a_{i}\right)$ is not a positive integer for each $i$, then the induced module $\mathrm{v}_{\mathrm{H}}=U \otimes_{U(\rho)^{\mathrm{H}}}$ is irreducible.

Proof. See Jantzen [5, Lemma 2].\\
Q.E.D.

Thus, Theorem 5.1 gives us many Whittaker modules.

There remain many open questions concerning the relation between Whittaker modules and completions of Verma modules. We have already mentioned the possibility of considering modules induced from infinite-dimensional $\|(m)$-modules. See Conjecture 5.1.

On the other hand, given a (finite-dimensional) U ( $m$ )-module $H$ such that $V_{H}$ is irreducible one can ask whether the Whittaker vectors given by Theorem 5.1 generate an irreducible Whittaker submodule of $\overline{\mathrm{V}}_{\mathrm{H}}$. One easily sees that this is equivalent to asking whether ${ }^{\left.\mathrm{H}\right|_{U(\bar{p})} \mathrm{N}}$ (via the generalized Harish-Chandra homomorphism) is irreducible. However, except for very simple cases (see Chapter 8) the image of the restriction of $U(\bar{p})^{\mathrm{N}}$ to $U(m)$ is quite complicated.

Further, it would be nice to know whether every irreducible Whittaker module, W, (with $\operatorname{dim} \mathrm{Wh}(\mathrm{W})<\infty$ ?) occurs as a submodule of the completion of a Verma module. One would also like to know which of these Whittaker modules are isomorphic. Recall [9] that if $p$ is a Borel subalgebra then the image of $U(\bar{p})^{\mathrm{N}}$ in $U(m)$ is precisely the (translated) Weyl group invariants in the enveloping algebra of a Cartan subalgebra. It would be quite interesting if some analogous finite group played a role in the generalized Whittaker theory.

Lastly, given a (not necessarily irreducible) generalized Whittaker module, $V_{H}$, it is probably true that $\operatorname{dim} W h\left(V_{H}\right)=\operatorname{dim}(H)$. This result is true for Borel subalgebras. Since every generalized Verma module contains an irreducible generalized Verma module ( $V_{H}$ has a finite composition series and any irreducible submodule of $V_{H}$ is a generalized Verma module) it would be nice to know whether all of the Whittaker vectors in $\overline{\mathrm{V}}_{\mathrm{H}}$ actually lie in the completion of the irreducible submodule (submodules?). This would be true if the irreducible submodule took the form $\mathrm{V}_{\mathrm{H}_{1}}$ with $\operatorname{dim}\left(\mathrm{H}_{1}\right)=\operatorname{dim}(\mathrm{H})$.

Chapter 6

\section{Principal series and Whittaker vectors}
\subsection{Preliminaries}
Let $G_{\mathbb{C}}$ be a connected complex Lie group with Lie algebra $\sigma \gamma$. Let $G$ be a Lie subgroup of $G_{C}$ such that its Lie algebra, $\mathcal{J}_{0}$, is a real form of $\mathcal{\gamma}$. For our purposes, we may as well assume that $\mathbf{G}_{\mathbb{C}}$ is simply connected and that $G$ is connected although the extensions to moregeneral cases is direct. Let $G=K A N$ (resp. $\left.\gamma_{0}=k_{0} \oplus \sigma_{0} \oplus \pi_{0}\right)$ be an Iwasawa decomposition of $G$ (resp. $G_{0}$ ). If $M$ is the centralizer of $A$ in $K$ with Lie algebra $m_{0}$, then $\mathrm{P}=$ MAN (resp. $\rho_{0}=m_{0} \oplus \Omega_{0} \oplus \Omega_{0}$ ) is a mimimal parabolic subgroup (resp. subalgebra) of $G$ (resp. $\gamma_{0}$ ). Of course, $K$ is the normalizer of $k_{0}$ in $G, A=\exp \left(\sigma_{0}\right)$ and $N=\exp \left(\pi_{0}\right)$.

Now, let $T$ denote a Cartan subgroup of $M$ with Lie algebra $t_{0}$. Then, $H=T A$ is a maximally split Cartan subgroup of $G$ with Lie algebra $h_{0}=\Omega_{0}+t_{0}$. Note that H is abelian.

Note that the irreducible finite-dimensional representations of $\mathrm{P}, \mathrm{P}^{\wedge}$, may be characterized by $\hat{P}^{\wedge}=\hat{M}^{\wedge} \times \tau^{\prime}$ where $\hat{M}$ is the collection of\\
irreducible finite-dimensional representations of M and, given $v \varepsilon \sigma^{\prime}$, we let $a \rightarrow a^{v}$ be the character (one-dimensional representation) of A given by $a^{v}=\exp (v(x))$ if $a=\exp (x), x \varepsilon \iota_{0}$. Let $\Gamma$ denote the collection of irreducible finite-dimensional representations of $G_{\text {. }}$ Then, as is well known, any $\gamma \varepsilon \mathrm{M}^{\wedge}$ is a direct summand of $\tilde{\gamma} \mid \mathrm{G}$ for some $\gamma \varepsilon \Gamma$.

Let $H_{\mathbb{C}}$ denote the Cartan subgroup of $G_{\mathbb{C}}$ having Lie algebra $h$. Let $B_{\mathbb{C}}=H_{\mathbb{C}} \tilde{N}_{\mathbb{C}}$ be a Borel subgroup of $G_{\mathbb{C}}$ such that $N \subseteq \tilde{N}_{\mathbb{C}}$. Let $\Delta(\tilde{r}, h)$ define a system of positive roots for $h$ in $0 \gamma$. Then, on restriction to $\Omega, \Delta(\tilde{n}, h)$ maps onto $\{0\} \cup \Delta\left(n_{\circ} \Omega\right)$.

Now, we say that a weight (one-dimensional representation) of $\mathrm{H}_{\mathbb{C}}$ is dominant if its differential is a linear combination of elements of $\Delta(\tilde{h}, h)$. Every irreducible finite-dimensional representation $F$ of $G$ has a unique highest weight $\mu \varepsilon \hat{H}_{\mathbb{C}}$ which is dominant, occurs with multiplicity one and characterizes $F$. We write $F=F(\mu)$. Conversely, every dominant weight corresponds to a unique representation as above. The set of weights of $F$ will be written $\Delta(F) \subseteq \hat{H_{\mathbb{C}}}$.

Let $C \subseteq \sigma^{\prime}$ be the open complex Weyl chamber defined by\\
(6.1.1) $C=\left\{v \varepsilon \sigma^{\prime} \mid \operatorname{Re}\langle v, \phi\rangle>0\right.$ for all $\left.\phi \varepsilon \Delta_{+}\right\}$\\
where < , > is the usual extension of the Killing form to $\Omega^{\prime}$. (Since the Killing form is nondegenerate on $\sigma_{0}$, take a dual basis,...) Further, let $\tilde{\rho}=\frac{1}{2} \sum_{\varepsilon \phi \Delta(\tilde{\pi}, \sigma)}$. Let $\rho=\tilde{\rho} \mid \sigma$ and define\\
(6.1.2) $\quad C_{\rho}=\rho+C$.

Note that for any $(\sigma, v) \varepsilon M^{\wedge} \times \sigma^{\prime}=P^{\wedge}$, there is a $\gamma \varepsilon \Gamma$ with an irreducible M-subrepresentation $\gamma \geqslant \gamma_{1} \varepsilon \hat{\mathrm{MA}}$ such that $(\sigma, \mathrm{v})$ is equivalent to $\gamma_{1} \otimes v_{1}$ where $v_{1}$ is trivial on $M$ and $v_{1}$ restricted to A corresponds to an element of C .

Now let $(\sigma, v) \varepsilon M^{\wedge} \times \sigma^{\nu}$. Consider the space $X_{\sigma, v}^{\infty}$ of all $C^{\infty}$ complex valued functions $f$ on $G$ with values in a representation space of $\sigma$ such that


\begin{equation*}
f\left(g p^{-1}\right)=a^{v_{\sigma}(m) f(q)} \tag{6.1.3}
\end{equation*}


for $p=\operatorname{man}, m \varepsilon M, a \varepsilon A, n \varepsilon N, g \varepsilon G$. Then $X_{\sigma, V}^{\infty}$ becomes a G-module with respect to the induced representation, $\operatorname{Ind}_{p}^{G}(\sigma \otimes v \otimes 1)$ (a smooth principal series\\
representation), where $g \varepsilon G$ acts on $f \varepsilon X_{\sigma, V}^{\infty}$ so that $(g f)(h)=f\left(g^{-1} h\right)$ for all $h \varepsilon G$. Moreover, $X_{\sigma, V}^{\infty}$ is also a $U=U(\sigma)$-module where if $x \in \sigma, p \varepsilon P$, and $f \varepsilon X_{\sigma, V}^{\infty}$, then


\begin{equation*}
(x f)(h)=\left.\frac{d}{d t} f((\exp -t x) h)\right|_{t=0} . \tag{6.1.4}
\end{equation*}


For any $(\sigma, v) \varepsilon \hat{M} \times \Omega^{\prime}$, let $X_{\sigma, v}$ denote the space of K-finite ( $k$-finite) elements in $X_{\sigma, V}^{\infty}$.\\
6.2. The $U_{(n)}$-freeness of $X_{0, v}$ for $v \varepsilon C_{\rho}$.

The material in this section is due to Kostant [9, Chapter 5] and [8].

Let $W$ be the "little" Weyl group $W(0, \sigma)$. Regarding $W$ as operating on $\Omega^{\prime}$, one also introduces the translated Weyl group $W_{T}$ of all maps $\sigma^{\prime} \rightarrow \sigma^{\prime}$ of the form $v \rightarrow \tau(v+\tilde{\rho})-\tilde{\rho}$ for $\tau \varepsilon W$. Both $W$ and $W_{T}$ then operate as groups of automorphisms of $U(\sigma)$ since U (O) can be identified with the algebra of polynomial functions on 0 . Their respective algebras of invariants are denoted by $U(\pi)^{W}$ and $U(\pi)^{W}$.

By the Birkhoff-Witt theory one has


\begin{equation*}
U=U(k) \otimes U(0) \otimes U(\pi) \tag{6.2.1}
\end{equation*}


and hence also a direct sum $U=U(\Omega) \oplus(k U+U \Omega)$. For any $u \varepsilon U$, let $p_{u} \varepsilon U(\sigma)$ be the unique element such that $u-p_{u} \varepsilon k U+U_{n}$. Now let $U^{k}$ be the centralizer of $k$ in $U$. It is then a result of Harish-Chandra that the correspondence $\mathrm{u} \rightarrow \mathrm{p}_{\mathrm{u}}$ defines an algebra homomorphism $U^{\mathrm{k}} \rightarrow U(\sigma){ }^{\mathrm{W}_{\mathrm{T}}}$. In fact, the result asserts that $I=k U+U^{k}$ is an ideal in $U^{k}$ and the correspondence $u \rightarrow p_{u}$ defines an algebra exact sequence


\begin{equation*}
0 \rightarrow I \rightarrow U^{k} \rightarrow \mathcal{U}(\Omega)^{W_{T}} \rightarrow 0 . \tag{6.2.2}
\end{equation*}


See, e.g., [12], Proposition 3.1. One then notes that any $\quad v \varepsilon \sigma^{\prime}$ defines a maximal ideal $\bigcup_{v}^{k}$ in $U^{k}$ where

$$
\mathcal{U}_{v}^{k}=\left\{u \in \mathcal{U}^{k} \mid p_{u}(\kappa v)=0\right\}
$$

Here $\kappa$ is the element in $W$ taking all positive roots to negative roots.

Let $0 \varepsilon \mathrm{M}$ denote the trivial representation so that the modules $X_{0, V}$ belong to the spherical principal series and were studied in [8]. A distinguished element\\
in $X_{0, v}$ is the K-invariant $l_{v} \varepsilon X_{0, v}$ where if $g=k a n$ is the Iwasawa decomposition of $g \varepsilon G$, then


\begin{equation*}
l_{v}(g)=a^{v} . \tag{6.2.3}
\end{equation*}


The following result was proved in [8].\\[0pt]
Theorem 6.2.1. [Kostant] Let $c_{\rho} \subseteq 0 \tau^{\prime}$ be defined by (6.1.2). Then for any $v \varepsilon C_{\rho}$, the $U$-module $x_{0, v}$ is cyclically generated by $1_{\mathrm{v}}$. That is

$$
x_{0, v}=U l_{v} .
$$

Furthermore, if $U_{v}$ denotes the annihilator of $I_{v}$, then


\begin{equation*}
U_{\mathrm{v}}=U_{\mathrm{v}}^{\mathrm{k}}+U_{\mathrm{k}} . \tag{6.2.4}
\end{equation*}


The following lemma is Lemma 5.2 of [ 9 ].

Lemma 6.2. For any $v \varepsilon G^{\prime}$, one has the direct sum

$$
\left(U U_{v}^{k}+U k\right) \oplus(U(n) \otimes D)
$$

where $\mathrm{D} \subseteq \Omega$ is a graded subspace such that\\
$\operatorname{dim} D=\operatorname{card}(W)$.

Together with Theorem 6.2.1, this implies

Theorem 6.2.2. Let the translated Weyl chamber $c_{\rho} \subseteq \Omega^{\prime}$ be defined by (6.1.2). Let $v \varepsilon C_{n}$ and let $X_{0, v}$ be the $K$-finite spherical principal series $U$-module defined in this section. Then $X_{0, V}$ is a free U( $n$ )-module with card ( $W$ ) generators where $W$ is the Weyl group $W(G, C l)$.

Proof. Let $U(n) \otimes D \rightarrow x_{0, v}$ be the map defined by $\mathrm{u} \rightarrow \mathrm{ul}_{\mathrm{v}}$. Then this map is a bijection by Lemma 6.2 and Theorem 6.2.1. This proves the theorem since $\operatorname{dim} D=\operatorname{card} W$.\\
Q.E.D.

Corollary 6.2. If $\mathrm{v} \varepsilon \mathrm{C}_{\rho}$ and $\mathrm{\eta}$ is an admissible homomorphic of $n$ into $\mathbb{C}$, then\\
(6.2.5) $\operatorname{dim} W h\left(X_{0, v}^{\prime}\right)=\operatorname{card}(W)$\\
where $w$ is the Weyl group $w\left(\sigma_{y}, \Omega\right)$ and $W h\left(X_{0, v}^{\prime}\right)$ is the space of Whittaker vectors in the (algebraic) dual $U$-module of $x_{0, v}$.

Proof. Since $W h\left(X_{\sigma, v}^{\prime}\right)$ is the orthocomplement of $U_{\eta}(n) X_{\sigma, v}$, the lemma follows immediately from

Theorem 6.2.2.\\
Q.E.D.

The remainder of this chapter will be devoted to calculating $\operatorname{dim} W h\left(X_{\sigma, v}^{\prime}\right)$ for $(\sigma, v) \varepsilon M^{\wedge} \times \Omega^{\prime}$.\\
6.3. Coherent continuation of characters

The material in this section is due to Vogan [15] and Speh-Vogan [12].

The following lemma is well known.

Lemma 6.3.1. Let $(\sigma, \mathrm{v}) \varepsilon \mathrm{M}^{\wedge} \times \sigma{ }^{\prime}$. If F is a finite dimensional representation of $G$, choose a family $0=\mathrm{F}_{0} \subseteq \ldots \subseteq \mathrm{~F}_{\mathrm{i}-1} \subseteq \mathrm{~F}_{\mathrm{i}} \subseteq \ldots \subseteq \mathrm{F}_{\mathrm{k}}=\mathrm{F}$ of P-invariant subspaces of $F$ such that $V_{i}=F_{i} / F_{i-1}$ is an irreducible P-module. Then, $X_{\sigma, V} \otimes F$ has a family $0=\mathrm{x}_{0} \subseteq \mathrm{x}_{1} \subseteq \cdots \subseteq \mathrm{x}_{\mathrm{k}}=\mathrm{x}_{\sigma, \mathrm{v}} \otimes \mathrm{F}$ of G-invariant subspaces, such that


\begin{equation*}
x_{i} / x_{i-1} \cong x_{(\sigma, v) \otimes V_{i}} \tag{6.3.1}
\end{equation*}


where $X_{(\sigma, V) \otimes V_{i}}$ is the $K$-finite induced module from the tensor product representation $(\sigma, \mathrm{v}) \otimes \mathrm{V}_{\mathrm{i}}$ of P .

Proof. For formal reasons, $X_{\sigma, V} \otimes F \cong \operatorname{Ind}_{p}^{G}\left(\left.(\sigma \otimes v \otimes 1) \otimes_{p} F\right|_{p}\right)$. The result follows by the exactness of Ind.\\
Q.E.D.

Of course, Lemma 6.3 is true in a much wider context.

Combining Lemma 6.3 with $(4.6 .12)$ we see that in the above notation, $\operatorname{dim} W h\left(X_{\sigma, v}^{\prime} \otimes F^{\prime}\right)=\sum_{i=1}^{k} \operatorname{dim} W h\left(X_{(\sigma, v) \otimes V_{i}}^{\prime}\right)$.

We now recall the basic facts of the character theory of G. (See Harish-Chandra [18].) Let $\pi$ be an admissible representation of $G$ on a Hilbert space, with a finite composition series and such that $\left.\pi\right|_{K}$ is unitary. If $f \varepsilon C_{C}^{\infty}(G)$, we define $\pi(f)=\int_{G} f(g) \pi(g) d g$. Then, $\pi(f)$ is an operator of trace class and $f \rightarrow \operatorname{tr} \pi(f)$ defines a distribution on $G$. This distribution is called the character of $\pi$ and is written $\theta(\pi)$. If $X$ is the Harish-Chandra module of K-finite vectors for $\pi$, we may define $\Theta(X)=\Theta(\pi)$. Every irreducible Harish-Chandra module can be realized in this way ([18]) and thus $\Theta(\cdot)$ is well defined. If $X$ has a finite composition series, suppose that $X$ has the irreducible composition factors $X_{1}, \ldots, X_{r}$ (listed with multiplicity). Then\\
$\theta(X)=\sum_{i=1}^{r} \Theta\left(X_{i}\right)$. We take this as a definition of $\Theta(X)$ whenever $X$ has a finite composition series. If $X_{1}, \ldots, X_{r}$ are inequivalent irreducible Harish-Chandra modules, then $\theta\left(X_{1}\right), \ldots, \theta\left(X_{r}\right)$ are linearly independent.

By a virtual representation, we will mean a formal finite combination of irreducible representations with integral coefficients. By the above, to any virtual representation we may assign a distribution character which vanishes if and only if the virtual representation does. By a character, we will always mean the character of a virtual representation.

Now, if $F$ is a finite-dimensional representation of $G$ and $X$ is a Harish-Chandra module, then $X \otimes F$ is a Harish-Chandra module. Moreover, $X \otimes F$ has a finite composition series if $X$ does. In this case, we have that $\theta(X \otimes F)=\theta(X) \theta(F)$. To see this, it is enough to assume that $X$ is irreducible. But then, the above formula follows directly from the definition. (Notice that $\Theta(F)$ is a smooth function on $G$ so that $\Theta(X) \Theta(F)$ is defined.) It follows that multiplication by $\Theta(F)$ defines a homomorphism from the group of characters to itself.

Now, if $\mu \varepsilon \widehat{H_{\mathbb{C}}}$ is the weight of a finite-dimensional representation, a map $S \mu$ taking characters into characters is defined. If a character $\Theta$ has an infinitesimal character represented by a nonsingular dominant weight $\gamma \varepsilon h^{*}$, we may then write $\theta=\theta(\gamma)$ and $S \mu \cdot \theta=\theta(\gamma+\mu)$.

We will use the following result

Lemma 6.3.2. If $F$ is a finite-dimensional representation of $G$, then\\
(6.3.2)

$$
\theta \cdot \theta(F)=\sum_{\mu \varepsilon \Delta(F)} S \mu \cdot \theta
$$

(as functions on the regular elements of G). Furthermore, if $\theta=\theta(\gamma)$ with $\gamma$ dominant and nonsingular, then\\
(6.3.3) $\quad \theta(\gamma) \cdot \theta(F)=\sum_{\mu \varepsilon \Delta(F)} \theta(\gamma+\mu)$.

Proof. See Lemma 5.3 in [15].\\
Q.E.D.\\
6.4. The theorem $\operatorname{dim} W h\left(X_{\sigma, v}^{\prime}\right)=|W| \operatorname{dim}(v)$.

Now, fix a representation $(\sigma, v) \in \mathrm{P}^{\wedge}$ and consider it as a representation of $m$. Then, the weights for the action of $h$ in ( $\sigma, v$ ) take the form $\lambda_{0}+\mu$ for $\lambda_{0} \varepsilon C p$ and $\mu$ is the weight of a finite-dimensional y-module. Let $\Lambda=\left\{\lambda_{0}+\mu \mid \mu\right.$ is the weight of a finite-dimensional $U$-module $\}$ and let $\left\{\theta\left(\lambda_{0}+\tilde{\mu}\right) \mid \mu\right.$ as above $\}$ be coherent family such that $\theta\left(\lambda_{0}+\tilde{\mu}_{0}\right)=\theta\left(X_{\sigma, v}\right)$. (See [12].)

Theorem 6.4. Let $(\sigma, v) \varepsilon P^{\wedge}$ be a representation of $P$ and let $X_{\sigma, V}$ be the induced Harish-Chandra module. Let Wh ( $X_{\sigma, v}^{\prime}$ ) denote the space of Whittaker vectors in the dual of $X_{\sigma, v}$. Then,\\
(6.4.1) $\quad \operatorname{dim} W h\left(X_{\sigma, v}^{\prime}\right)=\operatorname{dim}(\sigma)|W|$\\
where $\operatorname{dim}(\sigma)$ is the dimension of a representation space of $\sigma$. Here $|w|=\operatorname{card} w(0, \Omega)$.

Proof. For any $(\sigma, v) \varepsilon \hat{P}$, we define


\begin{equation*}
f(\sigma, v)=\operatorname{dim} w h\left(X_{\sigma, v}^{\prime}\right) . \tag{6.4.2}
\end{equation*}


We claim that $f(\sigma, v)<\infty$ for all $(\sigma, v) \varepsilon \hat{P}$. However, we note that $X_{\sigma, v}$ may be realized as one of the\\
subquotients of a tensor product $X_{0, v_{1}} \otimes F$ with $v_{1} \varepsilon C p$ by Lemma 6.3.1 and the comments following (6.1.2). By (4.6.12) and (4.2.2),\\
$\infty>\operatorname{dim} W h\left(\left(X_{0, V_{1}} \otimes F\right)^{\prime}\right)=|W|(\operatorname{dim} F) \geq \operatorname{dim} W h\left(X_{\sigma, V}^{\prime}\right)$\\
where $|W|$ is defined above. This proves the claim.\\
Now, by (4.6.12) we have that


\begin{equation*}
f(\sigma, v)=\operatorname{dim} W i n\left(X_{\sigma, v}^{\prime}\right)=\Sigma \operatorname{dim} W h\left(Y^{\prime}\right) \tag{6.4.3}
\end{equation*}


where $Y$ is an irreducible subquotient of a composition series of $X_{\sigma, V}$ (listed with multiplicites). Thus, if $\theta$ is an integral combination of irreducible characters, we can define $f(\theta) \varepsilon Z$.

Now, let $\Lambda$ be as above. We may then define $f(\theta(\lambda))=f(\lambda)$ for any $\dot{\lambda} \varepsilon \Lambda$. Recall that $\theta(\lambda) \otimes F=\sum_{\mu \varepsilon \Delta(F)} \theta(\lambda+\mu)$ for any finite-dimensional representation, $F$. Thus, $f(\theta(\lambda) \otimes F)=\Sigma f(\theta(\lambda+\mu))$. $\mu \varepsilon \Delta(\mathrm{F})$

On the other hand, (4.6.12) says that $f(\theta(\lambda) \otimes F)=(\operatorname{dim} F) f(\theta(\lambda))$. Thus,


\begin{equation*}
f(\lambda)=\frac{1}{\operatorname{dim} F} \sum_{\mu \varepsilon \Delta(F)} f(\lambda+\mu) \tag{6.4.4}
\end{equation*}


By Lemma 4.3 in [ 5 ], this implies that $f(\lambda)=f(\theta(\lambda))$ is a harmonic polynomial in $\lambda$ on the lattice $\Lambda$.

Now, given $\lambda \varepsilon \Lambda$, then $\lambda$ corresponds to $(\underline{+})$ a representation of MA. Further, $f$ is skew under the action of the Weyl group of $(m, \sigma), W_{M}$. This implies that $f$ is divisible by the dimension function $\operatorname{dim}(\sigma)$. Thus,


\begin{equation*}
f(\sigma, v)=\operatorname{dim}(\sigma) g(\sigma, v) \tag{6.4.5}
\end{equation*}


for some $\mathrm{W}_{\mathrm{M}}$-invariant polynomial $\mathrm{g}(\sigma, \mathrm{v})$. (See Proposition 4,9 in [1.5] for a proof of a similar statement.)

Recall that $f(\sigma, v)$ is a non-negative integer. Thus, $g(\sigma, v)>0$. From this, it easily follows that we may find $(\bar{\sigma}, \bar{v}) \varepsilon P^{\wedge}$ such that $g(\bar{\sigma}, \bar{v})$ is minimal. Let $p=g(\bar{\sigma}, \bar{v})$. Let $F$ be a finite dimensional representation of $G$ such that $\left.(\bar{\sigma}, \bar{v}) \otimes F\right|_{\text {MA }}$ contains a representation $\left(0, \mathrm{v}_{1}\right) \varepsilon \mathrm{M}^{\prime} \times \mathrm{C}_{\rho}$. Then, by (4.6.12) we have


\begin{equation*}
p(\operatorname{dim} F)=\Sigma f\left(\sigma_{i}, v_{i}\right) \tag{6.4.6}
\end{equation*}


where the $\left(\sigma_{i}, v_{i}\right)$ run over the collection of MA-subrepresentations of $(\bar{\sigma}, \bar{v}) \otimes F$. However,\\
$\Sigma f\left(\sigma_{i}, v_{i}\right) \geq \Sigma \rho\left(\operatorname{dim} \sigma_{i}\right)$. Since $\rho$ is minimal, this implies that $f\left(\sigma_{i}, v_{i}\right)=\left(\operatorname{dim} \sigma_{i}\right)$ for all ( $\sigma_{i}, v_{i}$ ) above. However, $f\left(0, v_{1}\right)=|W|$ by (6.2.5). Thus, $\rho=|W|$. Conversely, given $(\sigma, v) \varepsilon \hat{P}$, find $\left(0, v_{1}\right) \varepsilon \hat{P}$ such that $v_{1} \varepsilon C \rho$ and there is a finite-dimensional representation $F$ of $G$ such that ( $\sigma, v$ ) occurs in $\left.\left(0, v_{1}\right) \otimes F\right|_{M A}$. Again by (4.6.12) we have that $(\operatorname{dim} F)|w|=\Sigma f\left(\sigma_{i}, v_{i}\right) \geq|w| \Sigma \operatorname{dim}\left(\sigma_{i}\right)$ where the sum runs over the MA-subrepresentations of $\left(0, v_{1}\right) \otimes F$. But then, we have that $\operatorname{dim}(\sigma)|w|=f(\sigma, v)$ for any $(\sigma, v) \in P^{\wedge}$.\\
Q.E.D.

A great deal of the above proof was suggested to me by conversations with Dave Vogan. It would be very interesting to see if there is any connection between Theorem 6.4 and Conjecture 4.10 of [15]. More generally, it would be nice to establish a connection between the generalized Whittaker theory and other results in [15]. For background in the techniques used in the above proof, see [12] and [15].

Classification of admissible nilpotent elements\\
7.1. Some generalities

In this chapter we will use some of the original notation of §1.1. Let $\mathcal{J}_{0}=k_{0} \oplus \Omega_{0} \oplus n_{0}$ be an Iwasawa decomposition of a real, semi-simple Lie algebra $\gamma_{0}$. Let $\rho_{0}=m_{0} \oplus \Omega_{0} \oplus n_{0}$ be the associated minimal parabolic. Let $\Delta_{+}=\Delta\left(n_{0}, \Omega_{0}\right)$ denote the set of roots for the action of $\Omega_{0}$ on $\Omega_{0}$ and let $\pi$ denote a fundamental system of roots in $\Delta_{+}$. For any $\Sigma \subseteq \pi$, we let $f_{0}(\Sigma)=m_{0}(\Sigma) \oplus_{0}(\Sigma)$ denote the decomposition of the parabolic subalgebra $\rho_{0}(\Sigma)$ into its reductive and nilpotent components. As usual, dropping the subscript "0" denotes complexification.

Let $G_{0}$ be any real Lie group with Lie algebra $\mathcal{J}_{0}$. Since, at this point, we are only interested in the adjoint action of $G_{0}$ on $\mathcal{J}_{0}$, we may consider $G_{0}$ to be a subgroup of $G$, the adjoint group of of. As usual, $P_{0}(\Sigma)$ will denote the normalizer of $\rho_{0}(\Sigma)$ in $G_{0}$ and we have a semi-direct product $P_{0}(\Sigma)=M_{0}(\Sigma) N_{0}(\Sigma)$ where $\mathrm{M}_{0}(\Sigma)$ is reductive and $\mathrm{N}_{0}(\Sigma)=\exp \left(\eta_{0}(\Sigma)\right)$.

Rather than consider admissible nilpotent elements directly, we define a slightly more general notion. We say that an element $f \varepsilon \forall \gamma_{0}$ is a real admissible\\
nilpotent element for $\rho_{0}(\Sigma)$ if

$$
\text { 1) } \pm \varepsilon d_{-1} \cap O \gamma_{0} \text {, and }
$$

(7.1.1)\\
2) $\left(\eta_{0}(\Sigma)\right)^{f}=\{0\}$.

Note that since $\left[i \Omega_{0}(\Sigma), f\right] \subseteq i \Omega_{0}(\Sigma), i=\sqrt{-1}$, we have that $f$ is actually an admissible nilpotent element for the complex parabolic subalgebra $p(\Sigma)$ in the sense of §1.2. Further, note that an admissible nilpotent for $p(\Sigma)$ is a real admissible nilpotent for $p(\Sigma)$ viewed as a real Lie algebra. Thus, to classify admissible nilpotent elements, it suffices to classify real admissible nilpotent elements.

Concerning the uniqueness of conjugacy classes of admissible nilpotents, we have the following theorem:

Theorem 7.1. If $f$ is an admissible nilpotent element for $p(\Sigma)$ (in the sense of §1.2) then $f$ is a Richardson element of $\bar{\Omega}(\Sigma)$. That is $f \varepsilon \bar{\Omega}(\Sigma)$ and $\operatorname{dim}(G f)=2 \operatorname{dim}(G / P(\Sigma))$. Moreover, if $f^{\prime}$ is another admissible nilpotent for $p(\Sigma)$, then $f^{\prime}$ is conjugate to $f$ under $M(\Sigma)$.

Proof. If $g \varepsilon M(\Sigma), f, x \varepsilon O \gamma$, then $[\operatorname{Ad}(g) f, x]=\operatorname{Ad}(g)\left[f, \operatorname{Ad}\left(g^{-1}\right) x\right]$. Since $\operatorname{Ad}(g) \wedge(\Sigma)=\eta(\Sigma)$\\
for any $\mathrm{g} \varepsilon \mathrm{M}(\Sigma)$, it follows that any $\mathrm{M}(\Sigma)$-conjugate of an admissible nilpotent element is also admissible.

By (Corollary l. 7.3) we have that if $f$ is an admissible nilpotent element for $\rho(\Sigma)$, then $\operatorname{dim}\left(\operatorname{of}^{f}\right)=\operatorname{dim} m(\Sigma)$. Thus, $\operatorname{dim}(G f)=\operatorname{dim} G-\operatorname{dim}\left(G^{f}\right)=\operatorname{dim} y-\operatorname{dim}\left(y^{f}\right)= 2 \operatorname{dim}(\pi(\Sigma))=2 \operatorname{dim}(G / P(\Sigma))$. Hence, $f$ is a Richardson element. Now, by Lemma 1.1 in [4], we have that the $\overline{\mathrm{P}}(\Sigma)$ orbit of f is dense in $\bar{\Omega}(\Sigma)$. This proves the theorem since if $x \in d_{-1}$, then $P^{-}(\Sigma) x \cap d_{-1}=M(\Sigma) x$.\\
Q.E.D.

Thus, admissible nilpotent elements are precisely those Richardson elements lying in $d_{-1}$. By a theorem of Richardson ([11], p. 23 ) every parabolic subalgebra of of has a Richardson element. On the other hand, as we will see below, it is not true that every parabolic subalgebra has admissible nilpotent elements.

Corollary 7.1. Let $f$ be an admissible nilpotent element for $\rho(\Sigma)$. Then the set of real admissible nilpotent elements for $p_{0}(\Sigma)$ is given by


\begin{equation*}
\mathrm{M}(\Sigma) \mathrm{f} \cap g_{0} . \tag{6.1.2}
\end{equation*}


Thus, if $\rho_{0}(\Sigma)$ has a real admissible nilpotent element, then the set of all such elements is dense in $d_{-1} \cap \mathcal{J}_{0}$. Moreover, there are only a finite number of $M(\Sigma)$-orbits of real admissible nilpotent elements in $d_{-1} \cap \delta_{0}$.

Proof. The first two statements follow from Theorem 7.1. The last statement follows from the fact that there are only a finite number of $\mathrm{M}_{0}(\Sigma)$-orbits in $\bar{n}(\Sigma)$ 。\\
Q.E.D.

In passing, we mention that if $f$ is a real admissible nilpotent for $\rho_{0}(\Sigma)$, then $\bar{\rho}(\Sigma)$ is a real admissible polarization of $f$ in the notation of Ozeki-Wakimoto [10] or Hesselink [4].

We now proceed to the classification problem. We will not try to distinguish between the various $M_{0}(\Sigma)$-orbits of real admissible nilpotent elements. Instead, we will content ourselves with showing which parabolics have real admissible nilpotents and to provide examples for the classical algebras and the various real forms of $\mathrm{G}_{2}$ and $\mathrm{F}_{4}$. For the real forms of $\mathrm{E}_{6}, \mathrm{E}_{7}$, and $\mathrm{E}_{8}$ we will study only certain parabolics since, although direct, the calculations are rather tedious for these algebras.

Clearly, in what follows we may restrict ourselves to\\
simple Lie algebras.\\
7.2. Conditions for the non-existence of real admissible nilpotents for the classical algebras.

We note that the classical Lie algebras may be grouped into three classes: $\left\{s \ell(n ; R), s \ell(n ; C), s u^{*}(2 n)\right\}$, $\{s o(n ; C), s o(p, q), s u(p, q), s p(p, q)\}$ and $\left\{s p(n ; R), s p(n ; C), s_{0}^{*}(2 n)\right\}$. Rather than study $s \ell(n ; R), s \ell(n ; C), s u *(2 n)$, and $s u(p, q)$ directly, it is equivalent to study the algebras $g \ell(n ; K), K=R, C$, or $H$ the quaternions, and $u(p, q)$ since these differ from the above algebras only by an abelian summand (resp. $R, C$, $R$, and $u(1))$. For convenience, we now take specific realizations of these algebras.

Let $\left\{e_{i j} \mid 1 \leq i, j \leq n\right\}$ be the usual basis for $g \ell(n ; K)$ over $K$. Let $\{\alpha\}$ be a basis for $K$ over R and set\\
(7.2.1)

$$
\begin{aligned}
x_{v_{i}-v_{j}}^{\alpha} & =\alpha e_{i j} \quad i \neq j \\
H_{i} & =e_{i i} \\
T_{i}^{\alpha} \quad & =\alpha e_{i i} \quad(\alpha \text { is purely imaginary })
\end{aligned}
$$

Then, (7.2.1) is a real basis for $g \ell(n ; K)$ and we may .\\
take $a_{0}=\sum_{i=1}^{n} R H_{i}, m_{0}=\sum_{i=1}^{n} R T_{i}^{\alpha} \quad$, and α purely imaginary\\
$n_{0}=\sum_{i<j} R x_{v_{i}-v_{j}}^{\alpha}$. Let $\left\{v_{k}\right\}, r=1, \ldots, n$, be the dual $\alpha$ basis to $\Omega_{0}$ defined by $v_{k}\left(H_{i}\right)=\delta_{i k}$. Then $H_{v_{i}-v_{k}}^{\alpha}$ has the $\Omega_{0}$-weight $\left(v_{i}-v_{k}\right)$ and we may take $\Delta_{+}=\Delta\left(\eta_{0}, \Omega_{0}\right)=\left\{v_{i}-v_{j} \mid i<j\right\}$ and $\Pi=\left\{v_{i}-v_{i+1} \mid i=1, \ldots, n-1\right\}$. Thus, $\Delta_{+}$is a root system of type $A_{n-1}$.

Next, let $\mathrm{n}=\mathrm{p}+\mathrm{q}$ and define the following elements in $\mathrm{g} \ell(\mathrm{n} ; \mathrm{K})$ :\\
(7.2.2)

$$
\begin{aligned}
& x_{v_{i}-v_{j}}^{\alpha}=\alpha e_{i j}-\bar{\alpha} e_{n+1-j, n+1-i} \quad(i<j) \\
& x_{v_{i}+v_{j}}^{\alpha}=\alpha e_{i, n+1-j}-\overline{\alpha e} e_{j, n+1-i} \\
& x_{-\left(v_{i}+v_{j}\right)}^{\alpha}=\alpha e_{n+1-i, j}-\overline{\alpha e} e_{n \_1-j, i} \\
& x_{v_{i}}^{\alpha, \rho}=\alpha e_{i, q+\rho}-\overline{\alpha e}{ }_{q+\rho, n+1-i} \\
& x_{-v_{i}}^{\alpha, \rho}=\alpha e_{n+1-i, q+\rho}-\overline{\alpha e}{ }_{q+\rho, i} \\
& T_{i}^{\alpha}=\alpha\left(e_{i i}+e_{n+1-i, n+1-i}\right) \quad \text { ( } \alpha \text { purely imaginary) } \\
& M_{\rho \sigma}^{\alpha}=\alpha e_{q+\rho, q+\sigma}-\overline{\alpha e} e_{q+\sigma, q+\rho} \text { and } \\
& H_{i}=e_{i i}-e_{n+1-i, n+1-i}
\end{aligned}
$$

where $1 \leq i, j \leq q, 1 \leq \rho, \sigma \leq p-q$, and $\{\alpha\}$ is a basis for $K$ over $R$. Let $g(p, q ; K)$ be the Lie algebra over $R$ spanned by this basis. Then,\\
(7.2.3)

$$
\begin{aligned}
& g(p, q ; R)=s o(p, q) \\
& g(p, q ; C)=u(p, q) \text {, and } \\
& g(p, q ; H)=s p(p, q)
\end{aligned}
$$

If we let (7.2.2) be a complex basis and we take $\alpha=1$ throughout, then (7.2.2) is a basis for so $(2 q ; C)$ if\\
$p=q$ and a basis for so $(2 q+1 ; C)$ if $p=q+1$. As usual, we take $\sigma_{0}$ to be spanned by the $H_{i}, m_{0}$ to be spanned by the $T_{i}^{\alpha}$ and $M_{\rho \sigma}^{\alpha}$ and $\eta_{0}$ to be spanned by the $x_{v_{i}-v_{j}}^{\alpha}(i<j), x_{v_{i}+v_{j}}^{i<}$ and $x_{v_{i}}^{\alpha, \rho}$. Again, if $\left\{v_{k}\right\}$ is the dual basis to $\Omega_{0}$ given by $v_{k}\left(H_{i}\right)=\delta_{i k}$, then the various X's have the weights indicated by their subscripts. From this it follows that $\mathbb{I}=\left\{v_{i}-v_{i+1}, \beta \mid i=1, \ldots, q-1\right\}$ where $\beta=v_{q}$ if we are are considering $q(p, q ; K), p>q, p \neq q$ or $\operatorname{so}(2 q+1 ; C)$, $\beta=2 v_{q}$ if we are considering $g(g, q ; K), K=C$ or $H$, and $\beta=v_{q-1}+v_{q}$ if we are considering $g(q, q ; R)$ or so $(2 q ; C)$. The corresponding root systems, $\Delta_{+}$, are of types $B_{q}, C_{q}$, and $D_{q}$ respectively.

Finally, we define the following elements in $\mathrm{g} \ell(2 \mathrm{n}+1 ; \mathrm{K})$ :\\
(7.2.4)

$$
\begin{aligned}
& x_{v_{i}-v_{j}}^{\alpha}=\alpha e_{i j}-\bar{\alpha} e_{n+j, n+i} \\
& x_{v_{i}+v_{j}}^{\alpha}=\alpha e_{i, n+j}+\overline{\alpha e} e_{j, n+i} \\
& x_{-\left(v_{i}+v_{j}\right)}^{\alpha}=\alpha e_{n+j, i}+\bar{\alpha} e_{n+i, j} \\
& x_{v_{i}}^{\alpha}=\alpha e_{i, 0}+\overline{\alpha e} e_{0, n+i} \\
& x_{-v_{i}}^{\alpha}=\alpha e_{0, i}+\overline{\alpha e} e_{n+i, 0} \\
& H_{i}=e_{i i}-e_{n+i, n+i} \\
& T_{i}^{\alpha}=\alpha\left(e_{i i}+e_{n+i, n+i}\right) \quad \text { ( } \alpha \text { is purely imaginary) } \\
& T_{0}^{\alpha}=\alpha e_{00} \quad \text { ( } \alpha \text { is purely imaginary) }
\end{aligned}
$$

$I \leq i, j \leq n,\{\alpha\}$ is a basis for $K$ over $R$. For $K=H$, this defines a real basis for $s^{*}(4 n+2)$. If we drop the elements $T_{0}^{\alpha}, X_{v_{i}}^{\alpha}$ and $X_{-v_{i}}^{\alpha}$, this defines a real basis for $\operatorname{so} *(4 n)$. If we further assume that $\alpha=1$, this defines a real basis for $s p(n ; R)$ and a complex basis for $\operatorname{sp}(n ; C)$. Take the $\left\{v_{i} \mid i=1, \ldots, n\right\}$ as usual and note that the X's have the weights indicated by their subscripts. Note that $\sigma_{0}$ is spanned by the $H_{i}, m_{0}$ is spanned by the $T_{i}^{\alpha}$ and $n_{0}$ is spanned by the $X_{v_{i}-v_{j}}^{\alpha}(i<j)$,\\
$x_{v_{i}+v_{j}}^{\alpha}$ and $x_{v_{i}}^{\alpha}$. Note that $\Pi=\left\{v_{i}-v_{i+1}, v_{n} \mid i=1, \ldots, n-1\right\}$ for $s^{*}(4 n+2)$ and $\Pi=\left\{v_{i}-v_{i+1}, 2 v_{n} \mid i=1, \ldots, n-1\right\}$ for. the other algebras. Thus, $\Delta_{+}$is a root system of type $B_{n}$ for so*(4n+2) and of type $C_{n}$ for the other algebras.

Note that there has been a degree of sloppiness in the defining of $\Omega_{0}$ above for the algebras $s o(n ; C)$ and $s p(n ; C)$ in that we should take $\sum_{i=1}^{n} \sqrt{-1} \mathrm{RH}_{i} \subseteq m_{0}$ rather than in $\sigma_{0}$ as the above definitions indicate. However, this distinction has no bearing on what follows.

If $\Delta_{+}$is of type $B_{n}, C_{n}$, or $D_{n}$ we define $\beta=v_{n}, 2 v_{n}$, or $v_{n-1}+v_{n}$ respectively. If $\Delta_{+}$is of type $A_{n-1}$, we ignore $\beta$. Then, for any of these root systems we have $\Pi=\left\{v_{i}-v_{i+1}, \beta \mid i=1, \ldots, n-1\right\}$. Now, let $0=I_{0}<1 \leq I_{1}<I_{2}<\ldots<I_{R}<n=I_{R+1}$ be integers. Then, any $\Sigma \subseteq \Pi$ is given by $\Sigma=\left\{v_{I_{p}}-v_{I_{p}+1} \mid p=1, \ldots, R\right\}$ or $\Sigma=\left\{v_{I_{p}}-v_{I_{p}+1}, \beta \mid p=1, \ldots, R\right\}$ for some $\left(I_{p}\right)$ as above. We let $\ell_{p}=I_{p}-I_{p-1}$.

Let $\Sigma=\left\{\mathrm{v}_{\mathrm{I}_{\mathrm{p}}}-\mathrm{v}_{\mathrm{I}_{\mathrm{p}}+1}\right\}$ as above and let $\Sigma^{\prime}=\Sigma U\{\beta\}$ if $\Delta\left(\pi_{0}, \pi_{0}\right)$ is of type $B_{n}, C_{n}$, or $D_{n}$. Let\\
$d_{-1, p} p=1, \ldots, r$, be the space spanned by the elements $\mathrm{x}_{\mathrm{v}_{I_{p}+\mathrm{i}}^{\alpha}}^{\alpha}-\mathrm{v}_{I_{p-1}+j}, i \leq \ell_{p+1}, j \leq \ell_{p}$, where if $\Delta\left(\eta_{0}, \sigma_{0}\right)$ is of type $B_{n}, C_{n}$, or $D_{n}$, we adjoin the elements $x_{-v_{I_{r-1}}^{\alpha}+i}$ and $\left.x_{-\left(v_{1_{r-1}+i}^{\alpha}\right.}+v_{I_{r}+j}\right)$ when defined and $\beta \notin \Sigma$. If $\beta \varepsilon \Sigma$, we let $d_{-1, r+1}$ be the sapce spanned by the $\left.x_{-\left(v_{I_{r}+i}^{\alpha+v_{I_{r}}+j}\right.}\right)$ if $\beta=2 v_{q}$ or $v_{q-1}+v_{q}$ and $\mathrm{X}_{-\mathrm{v}_{\mathrm{I}_{r}+\mathrm{i}}}^{\alpha}$ if $\beta=\mathrm{v}_{\mathrm{q}}$. Then, we have a direct sum $d_{-1}=\stackrel{r, r+1}{p=1} d_{-1, p}$. Thus, given any $f \varepsilon d_{-1}$, we have $f=\sum_{p=1}^{r, r+1} f_{p}$ where $f_{p} \varepsilon d_{-1, p}$. Note that $f_{p}(p<r$ or $r+1)$ may be regarded as a $\ell_{p} \times \ell_{p+1}$ matrix and has a rank less than or equal to the minimum of $\ell_{p}$ and $\ell_{p+1}$.

Lemma 7.2. Let $\Delta^{+}$be of type $B_{n}$ or $C_{n}$. Then if $d_{-l}$ contains an admissible nilpotent element, we have (in the above notation)\\
(7.2.4)

$$
\ell_{1} \leq \ell_{2} \leq \cdots \leq \ell_{r} .
$$

If $\Delta^{+}$is of type $D_{n}$ and $d_{-1}$ contains an\\
admissible nilpotent element, then there is an $R \varepsilon \mathrm{Z}_{+}$ such that\\
(7.2.5)

$$
\ell_{1} \leq \cdots \leq \ell_{R}=\ell_{R+1}+1=\ldots=\ell_{r+1}+1 .
$$

(Note that we may have $\mathrm{R}=\mathrm{r}+\mathrm{l}$.)\\
If $\Delta^{+}$is of type $A_{n}$ and $d_{-1}$ contains an admissible nilpotent element, then there is an $R \varepsilon Z_{+}$ such that


\begin{equation*}
\ell_{1} \leq \cdots \leq \ell_{R} \geq \ell_{R+1} \geq \cdots \geq \ell_{r+1} . \tag{7.2.6}
\end{equation*}


Proof. If $\Delta^{+}$is of type $B_{n}$ or $C_{n}$ suppose that (7.2.4) doesn't hold. Then there is a $p$ such that $\ell_{p}>\ell_{p+1}$. If $f \varepsilon d_{-1}$, then $f_{p}$ has rank at most $\ell_{p+1}$. Now, given $x=x_{+v_{I_{p-1}}^{\alpha}}^{\alpha}$. or $x_{+2 v_{I_{p-1}+i}^{\alpha}}^{\alpha} \varepsilon n_{o^{\prime}}$\\
we have that $[f, x]=\left[f_{p}, x\right]$ as one easily sees. However, these span an $\ell_{p}$ dimensional space whereas $f_{p}$ has rank at most $\ell_{p+1}<\ell_{p}$. Thus, there is an $x$ in the span of these elements so that $\left[f_{p}, x\right]=[f, x]=0$. If follows that there are no real admissible nilpotents in $d_{-1}$ unless (7.2.4) holds.

Similar arguments show that (7.2.5) and (7.2.6) must hold by considering spaces spanned by terms of the form $\mathrm{X}_{\mathrm{v}_{\mathrm{p}}+\mathrm{i}+\mathrm{v}_{\mathrm{I}_{\mathrm{q}}+\mathrm{j}}^{\alpha}}$ with $\mathrm{i} \leq \ell_{\mathrm{p}+1}, \mathrm{j} \leq \ell_{\mathrm{q}+1}$.\\
Q.E.D.\\
7.3. Conditions for the existence of real admissible nilpotents for the classical algebras.

In this section, we will show that the necessary conditions for the existence of real admissible nilpotent elements given by Lemma 7.2 are actually sufficient.

It is interesting to note that for any nilpotent element in $\mathrm{g} \ell(\mathrm{n} ; \mathrm{C})$, there is a parabolic subalgebra for which it is admissible. This can be easily seen from the Jordan form of any nilpotent and the forms of the elements given by (7.3.2).

Proposition 7.3.1. Let $\gamma_{0}=g \ell(n ; K), K=R, C$, or $H$. Let $\Sigma=\left\{v_{I_{p}}-v_{I_{p}+1} \mid 1 \leq I_{p}<I_{p+1}<n, p=1, \ldots, r\right\} \leq \mathbb{I}$.

Then $p_{0}(\Sigma)$ has a real admissible nilpotent element if and only if there is an $\mathrm{R} \varepsilon \mathrm{Z}_{+}$such that\\
(7.3.1)

$$
\ell_{1} \leq \ell_{2} \leq \cdots \leq \ell_{R} \geq \ell_{R+1} \geq \cdots \geq \ell_{r+1} .
$$

where $\ell_{p}=I_{p}-I_{p-1}$ and $I_{0}=0, I_{r+1}=n$.

Proof. If $\Sigma$ is not of the above form, then $\rho_{0}(\Sigma)$ has no real admissible nilpotents by Lemma 7.2.

Conversely, suppose that $\Sigma$ is of the above form. Consider the element $f \varepsilon d_{-1}, f=\sum_{p=1}^{r} f_{p}$, where


\begin{equation*}
f_{p}=\sum_{\sigma=1}^{\min \left(\ell_{p}, \ell_{p+1}\right)} e_{I_{p}+\sigma, I_{p-1}+\sigma .} \tag{7.3.2}
\end{equation*}


We claim that $f$ is a real admissible nilpotent. Suppose there is an $x \in \eta$ such that $[x, f]=0$. Since $f$ is ad $x_{0}$-stable, we may suppose that $x \varepsilon d_{s}$ for some $s>0$. Further, the form of $f$ shows that we need only consider $x$ of the form $x=\sum a_{p} e_{I_{p}+i, I_{p+s}+j}$ such that $a_{p} \varepsilon K$ and $i \leq \ell_{p+1}, j \leq \ell_{p+s+1}$ for all $p$ such that $a_{p} \neq 0$. Assume that $i \geq j$. (A similar proof works for $i \leq j$. ) Let $p_{0}$ be minimal such that $a_{p_{0}} \neq 0$. Then, $[f, x]=-a_{p_{0}}^{\prime} e_{I_{p_{0}}}+i, I_{p_{0}+s-1}+j+$ (terms in $d_{s-1}$ of the form $b_{p} e_{I_{p}+i, I_{p+s-1}+j}$ with $p>p_{0}$.) Thus, if $[f, x]=0$, then we must have $a_{p_{0}}^{\prime}=0$. Now, if $p_{o}+s>R$, then $a_{p_{0}}^{\prime}$ cannot be zero since then $\ell_{p_{0}+s} \geq \ell_{p_{0}+s+1} \geq j$.

On the other hand if $\mathrm{p}_{0}+\mathrm{s} \leq \mathrm{R}$, then $\ell_{\mathrm{p}_{0}+\mathrm{s}} \geq \ell_{\mathrm{p}_{0}+1} \geq \mathrm{i} \geq \mathrm{j}$ by assumption. Thus, $a_{p_{0}}^{\prime}$ cannot be zero. This is a contradiction. Hence, $\mathcal{n}_{0}^{f}=\{0\}$.\\
Q.E.D.

Proposition 7.3.2. Let $\mathcal{J}_{0}=g(n, q ; K), K=R, C$, or $H, n \geq q$. Let $\Sigma=\left\{v_{I_{p}}-v_{I_{p}+1} \mid 1 \leq I_{p}<I_{p+1}<q, p=1, \ldots, r\right\} \leq \Pi$. Let $\ell_{p}=I_{p}-I_{p-1}$ with $I_{0}=0, I_{r+1}=q$. Let $\Sigma^{\prime}=\Sigma \cup\{\beta\}$ where $\beta=\mathrm{v}_{\mathrm{q}}$ if $\mathrm{n}>\mathrm{q},=2 \mathrm{v}_{\mathrm{q}}$ if $\mathrm{n}=\mathrm{q}$ and $\mathrm{K}=\mathrm{C}$ or $\mathrm{H},=\mathrm{v}_{\mathrm{q}=1}+\mathrm{v}_{\mathrm{q}}$ if $\mathrm{n}=\mathrm{q}$ and $\mathrm{K}=\mathrm{R}$. Then $p_{0}(\Sigma)$ has a real admissible nilpotent element if


\begin{equation*}
\ell_{1} \leq \ell_{2} \leq \cdots \leq \ell_{r+1} . \tag{7.3.3}
\end{equation*}


Also, $\rho_{0}\left(\Sigma^{\prime}\right)$ has a real admissible nilpotent element if (7.3.3) holds and if $n-q \geq \ell_{r+1}$ when $n \neq q$. Finally, if $\sigma \gamma_{0}=g(q, q ; R)$, then $p_{0}(\Sigma)$ and $p_{0}\left(\Sigma^{\prime}\right)$ have real admissible nilpotent elements if there is an $\mathrm{R} \varepsilon \mathrm{Z}_{+}$ such that


\begin{equation*}
\ell_{1} \leq \ell_{2} \leq \cdots \leq \ell_{R}=\ell_{R+1}+1=\cdots=\ell_{r+1}+1 . \tag{7.3.3}
\end{equation*}


Moreover, if $\Sigma$ or $\Sigma^{\prime}$ is not of the above form, then\\
there are no real admissible nilpotents for $\rho_{0}(\Sigma)$ or $p_{0}{\left(\Sigma^{\prime}\right)}^{\prime}$.

Proof. If $\Sigma$ satisfies (8.3.3), write\\
$f=\sum_{p=1}^{r} f_{p}+f^{\prime}$ where $f_{p}=\sum_{\sigma=1}^{\ell} x_{v_{I_{p}+\sigma}}-v_{I_{p-1}}+\sigma$\\
and $f^{\prime}=\sum_{\sigma=1}^{\ell_{r}} X_{-\left(v_{I_{r}+\sigma}+v_{I_{r-1}+\sigma}\right)}$. If $n \neq q$ and\\
$n-q \geq \ell_{r}$, we take $f$ as above except that we redefine\\
$f^{\prime}=\sum_{\sigma=1}^{\ell} x_{-v_{I_{r}}+\sigma}^{1, \sigma}$. If $n=q_{\text {, we take }} f$ as above except\\
that we now take\\
$f^{\prime}=\sum_{\sigma=1}^{\left[\ell_{r} / 2\right]} X_{\left.-v_{I_{r}}+2 \sigma-\right]^{-} v_{I_{r}+2 \sigma}}\left(+X_{-2 v_{q}}^{\sqrt{-1}}\right.$ if $K \neq R$ and $\ell_{r}$ is odd $)$.\\
Then, if $f$ is of any of the above forms (except when $n=q, K=R$ and $\ell_{r}$ is odd) one calculates that $f$ is an admissible nilpotent element for $g \ell(n+q ; C)$. Thus, $f$ is a real admissible nilpotent element for $g(n, q ; C)$ in any of the above cases. Further, since $f \varepsilon g(n, q ; R)$, the same holds in these cases. Next, note that the same holds for $g(n, q ; H)$ since the various $f$ don't mix more than two of the basis elements of $H$ over $R$ in a standard basis (\{l,i,j,k\}).

We are thus left with the cases (7.3.3)' and $\left(n=q, K=R, \Sigma \prime, \ell_{r}\right.$ odd $)$. If $\Sigma$ satisfies (7.3.3)' we take $f$ as above except that we replace $\ell_{R}$ with $\ell_{R}-1$ in the sum defining $f_{R}$. However, in any of these cases, it is a straightforward calculation to show that the centralizer of $f$ in the corresponding nilpotent subalgebra of $g \ell(2 q, C)$ is one dimensional and doesn't intersect $\sigma y_{0}$. Thus, $f$ is admissible in this case.

Now, by Lemma 7.2, we have that $p_{0}(\Sigma)$ or $p_{0}\left(\Sigma^{\prime}\right)$ has real admissible nilpotent elements only if one of the conditions stated above applies.\\
Q.E.D.

The only remaining simple Lie algebras that we need consider are $s p(n ; R), s p(n ; C)$ and $s^{*}(2 n)$.

Proposition 7.3.3. Let $G y_{0}=s p(n ; R), s p(n ; C), s_{0} *(4 n)$, or $\operatorname{so}^{*}(4 n+2)$. Let $\Sigma=\left\{v_{I_{p}}-v_{I_{p}+1} \mid p=1, \ldots, r, 1 \leq I_{p}<I_{p+1}<n \leq \Pi\right.$. Let $\ell_{p}=I_{p}-I_{p-1}$ with $I_{0}=0, I_{r+1}=n$. Let $\Sigma^{\prime}=\Sigma U\{\beta\}$ where $\beta=\mathrm{v}_{\mathrm{q}}$ if $\gamma_{0}=\mathrm{so}^{*}(4 \mathrm{n}+2)$ and $\beta=2 \mathrm{v}_{\mathrm{q}}$ otherwise. Then $\mathcal{P}_{0}(\Sigma)$ has a real admissible nilpotent element if


\begin{equation*}
\ell_{1} \leq \ell_{2} \leq \cdots \leq \ell_{r+1} . \tag{7.3.4}
\end{equation*}


Also, $\rho_{0}\left(\Sigma^{\prime}\right)$ has a real admissible nilpotent element if (7.3.4) holds and if $\ell_{r+1}=1$ when $\gamma_{0}=s_{0} *(4 n+2)$. Moreover, if $\Sigma$ or $\Sigma^{\prime}$ is not of the above form, then there are no real admissible nilpotents for $\rho_{0}(\Sigma)$ or $p_{0}{ }^{\left(\Sigma^{\prime}\right)}$.

Proof. For so*(4n) and so* $(4 n+2)$, we take $f \varepsilon d_{-1}$ to have the same form as the corresponding f's in Proposition 7.3.2 except that we replace $x_{-\left(v_{i}+v_{j}\right)}^{1}$ with $x_{-\left(v_{i}+v_{j}\right)}^{\sqrt{-1}}$ and $x_{-2 v_{q}}^{\sqrt{-1}}$ with $x_{-2 v_{q}}^{1}$ in the appropriate\\
spots. Then, one sees on comparing the commutation relations among these bases that $f$ is admissible for so* $(4 n)$ (or so* $(4 n+2)$ ) if and only if the corresponding $f$ is admissible for $s p(n, n)$ (or $s p(n+1, n)$ ). Thus, Proposition 7.3.3 is true for these algebras.

If $\gamma_{0}=s p(n ; K)$, we take $f=\sum_{p=1}^{r} f_{p}$ where $f_{p}=\sum_{\sigma=1}^{\ell} x_{v_{I_{p}+\sigma}}-v_{I_{p-1}+\sigma}$ and we take $f^{\prime}=f+\sum_{\sigma=1}^{n-I_{r}} x_{-2 v_{I_{r}+\sigma}}$.

We claim that if (7.3.4) holds then $f$ (resp.f') is a real admissible nilpotent for $\rho_{0}(\Sigma)$ (resp. $\rho_{0}\left(\Sigma^{\prime}\right)$ ). However $p_{0}(\Sigma)^{f}$ is clearly zero by the same argument as in Proposition 7.3.1 since $f$ does not mix terms of the form\\
$\Sigma X_{v_{i}-v_{j}}$ with terms of the form $\Sigma X_{v_{i}+v_{j}}$ and is nonvanishing on these separately. On the other hand, comparing the algebra structures on $s p(n ; C)$ and $s u(n, n)$, one sees that $f^{\prime}$ is admissible for $\mathcal{p}_{0}{ }^{\left(\Sigma^{\prime}\right)}$ for the same reason that the corresponding element in $\operatorname{su}(n, n)$ is admissible.

7.4. Some comments on real admissible nilpotents in the exceptional Lie algebras

As we notices in §1.2, if $x_{0}$ is the mono-simple element of a TDS, then the parabolic subalgebra corresponding to $x_{0}$ has an admissible nilpotent element. Now Dynkin [3] has classified the mono-simple elements for the exceptional groups. Thus, any mono-simple element in the Dynkin classification which takes only the values zero and two on the simple roots gives rise to a parabolic subalgebra with admissible nilpotent elements. (Note that our $x_{0}$ is one of Dynkin's mono-simple elements divided by two.)

In particular, one can read off from the Dynkin classification the fact that any minimal parabolic of any real form of an exceptional Lie algebra has real admissible nilpotent elements.

I have shown, that the mono-simpie elements taking the values 0 and 2 only correspond to all of the parabolic subalgebras of $G_{2}$ and $F_{4}$ which have real admissible nilpotent elements. I have also shown that the maximal parabolics in $\mathrm{E}_{6}$ and $\mathrm{E}_{8}$ that don't correspond to mono-simple elements don't have admissible nilpotent elements either. (Note that all maximal parabolics in $\mathrm{E}_{7}$ correspond to mono-simple elements.) The proofs of these results, while elementary, are quite lengthy and tedious and won't be presented here.

Chapter 8\\
Examples

In this section we will study some examples which show some of the typical properties of $U_{(\bar{p})^{N}}^{N}$ (and $S(\bar{p})^{N}$ ). Recall that we have a surjection $\pi_{I}: S(\bar{p})^{N} \rightarrow \bar{p}^{f}=\gamma^{f}$ which is just the restriction of the projection $S(\bar{p}) \rightarrow S_{1}(\bar{p})$. See §1.7. This map forms a bijection between the set of generators for $S(\bar{p})$ (as a polynomial algebra) and a basis for $\bar{\rho}^{f}$. By theorem 2.3, we may find Birkhoff-Witt generators of $U(\bar{p})^{N}$ so that these map to a set of generators for $S(\bar{p})^{N}$ under $\tau(\cdot)$. In this section we will essentially be working backwards: we will easily be able to calculate of ${ }^{f}$ and we will then attempt to find a corresponding basis for $U(\bar{p})^{\mathrm{N}}$.

Recall that we have a projection map $\tilde{\pi}: U(\bar{p}) \rightarrow U(m)$ with kernel $n U(\bar{p})$ and that $\tilde{\pi}$ is an algebra isomorphism of $U(\bar{p})^{\mathrm{N}}$ onto its image in $U(m)$. See the comment following Corollary 2.3.2. In these examples it will become clear that the structure of $U(\bar{p})^{N}$ is much more complicated than just that of $\sigma^{f}$.\\
8.1. Maximal parabolic subalgebras of $g \ell(n ; K)$.

Let $K=R, C$, or $H$. We keep the same notation as in §7.2. Let $\Sigma \subseteq \Pi$ consist of the single root $\left\{v_{I}-v_{I+1}\right\}$.

Clearly, we may assume that $2 I \leq n+1$. It is easy to see that $d_{1}=\lambda_{0}(\Sigma)$ is spanned by $\left\{e_{i, I=j} \mid l \leq i \leq I, l \leq j \leq n-I\right\}$ and that $d_{0}=M_{0}$ is spanned by $\left\{e_{i j} \mid 1 \leq i, j \leq I\right\}$ and $\left\{e_{I+k, I+\ell} \mid I \leq k, \ell \leq n-I\right\}$. As in §7.2, we may take $f=\sum_{\sigma=1}^{I} e_{I+\sigma, \sigma}$ as a representative real admissible nilpotent element. As one easily sees, $\sigma_{0}^{f}=\bar{p}_{0}^{f}=m_{0}^{f} \oplus \bar{n}_{0}$. (Here we have dropped any mention of $\Sigma$ since it is assumed fixed hereafter.) Further, $m_{0}^{f}$ is spanned by $\left\{e_{i j}+e_{I+i, I+j} ; e_{I+i, I+r} ; e_{I+s, I+r} \mid I \leq i, j \leq I<r, s \leq n-I\right\}$. Thus, we have the picture:

$$
\begin{aligned}
& n_{0}=n_{(1)} \oplus n_{(2)} \\
& m_{0}=\oplus m_{(i)} \\
& \bar{n}_{0}=\bar{n}_{(1)} \oplus \bar{n}_{(2)}
\end{aligned}
$$

\includegraphics[max width=\textwidth, alt={}, center]{dc0495e0-61cd-4445-9972-03b906b7ff05-155_1009_1007_1524_629}\\
where $\left[\Omega_{0}, f\right]=M_{(4)} \oplus$ the space spanned by the antisymmetric elements in $M_{(1)} \oplus M_{(2)}$ (that is, of the form $e_{i j}-e_{I+i, I+j}$ ), and $m_{0}^{f}=m_{(3)} \oplus m_{(4)} \oplus$ the space spanned by the symmetric elements in $m_{(1)} \oplus m_{(2)}$. Now, the elements in $\mathcal{M}_{0}^{f}$ are already in $U(\bar{\rho})^{N}$ (and $s(\bar{\rho})^{N}$ ) as one easily checks, so we need only compute generators $u_{i j}^{\alpha} \varepsilon U_{(2)}(\bar{\rho})^{N}$ such that $\pi_{1}{ }^{\tau}(2) u_{i j}^{\alpha}=\alpha e_{I+i, j}$. To simplify notation, let $e_{i j}^{\alpha}=\alpha e_{i j}$.

Indeed, let\\
(8.1.1) $u_{i j}^{\alpha}=e_{I+i, j}^{\alpha}-\sum_{\sigma=1}^{I} e_{\sigma j}^{\gamma \alpha} e_{I+i, I+\sigma}^{\bar{\gamma}}$

$$
+\sum_{\tau=1}^{n-2 I} e_{2 I+\tau, I+j}^{\gamma} e_{I+i, 2 I+\tau}^{\gamma \alpha}-I\left(\operatorname{dim}_{R} K\right) e_{I+i, I+j}^{\alpha}
$$

for $l \leq i \leq n-I, l \leq j \leq I$ and $\gamma$ and $\alpha$ run over a basis of $K$ over $R$. Then, for any $e_{p, I+q}^{\beta} \varepsilon \eta$,

$$
\begin{aligned}
{\left[e_{p, I+q}^{\beta}, u_{i j}^{\alpha}\right] } & =\delta_{i q} e_{p j}^{\beta \alpha}-\delta_{p j} e_{I+i, I+q}^{\alpha \beta} \\
& +\sum_{\sigma=1}^{I}\left(e_{\sigma, I+q}^{\gamma \alpha \beta} e_{I+i, I+\sigma}^{\bar{\gamma}}\right)^{\eta} \delta_{p j} \\
& -\sum_{\sigma=1}^{I} e_{\sigma j}^{\gamma \alpha}\left(e_{p, I+\sigma}^{\beta \bar{\gamma}}\right)^{\eta} \delta_{i q} \\
& +\sum_{\tau=1}^{n-2 I} e_{2 I+\tau, I+j}^{\gamma \alpha}\left(e_{p, 2 I+\tau}^{\beta \bar{\gamma}}\right)^{\eta} \delta_{i q} \\
& +\sum_{\gamma}\left(e_{p, I+j}^{\beta \gamma \alpha} e_{I+i, I+q}^{\bar{\gamma}}\right)^{\eta}(i f q>I) \\
& -I\left(d_{i m_{R}} K\right)\left(e_{p, I+j}^{\beta \alpha}\right)^{\eta} \delta_{i q}
\end{aligned}
$$

If we take the Killing form to be $B(X, Y)=\frac{1}{2} \operatorname{Re}(\operatorname{tr} X Y)$ as matrices in $g \ell(n ; K)$, then $\left(e_{r, I+s}^{\delta}\right)^{\eta}=B\left(f, e_{r, I+s}^{\delta}\right)$\\
$=\operatorname{Re}(\delta) \delta_{r s}$ as one easily checks, and

$$
\begin{aligned}
{\left[e_{p, I+q}^{\beta}, \dot{u}_{i j}^{\alpha}\right]^{\eta} } & =\delta_{i q} e_{p j}^{\beta \alpha}-\delta_{p j} e_{I+i, I+q}^{\alpha \beta} \\
& +\sum_{\sigma=1}^{I} e_{I+i, I+\sigma}^{\bar{\gamma}} \operatorname{Re}(\gamma \alpha \beta) \delta_{p j} \delta_{\sigma q} \\
& -\sum_{\sigma=1}^{I} e_{\sigma j}^{\gamma \alpha} \operatorname{Re}(\beta \bar{\gamma}) \delta_{i q} \delta_{p \sigma} \\
& +\sum_{\tau=1}^{n-2 I} e_{2 I+\tau, I+j}^{\gamma \alpha} \operatorname{Re}(\beta \bar{\gamma}) \delta_{i q} \delta_{p, I+\tau} \\
& +\sum_{\gamma} e_{I+i, I+q}^{\bar{\gamma}} \operatorname{Re}(\beta \gamma \alpha) \delta_{p j}(i f q>I) \\
& +\sum_{\sigma=1}^{I} e_{\sigma, I+\sigma}^{\gamma \alpha \beta \bar{\gamma}}+\sum_{\sigma} e_{p, I+q}^{\beta \gamma \alpha \bar{\gamma}} \delta_{i j}(i f q>I) \\
& -I\left(\operatorname{dim}_{R} K\right) \operatorname{Re}(\beta \alpha) \delta_{p j} \delta_{i q} \\
& =0 .
\end{aligned}
$$

Since $e_{p, I+q}^{\beta} \cdot u_{i j}^{\alpha}=\left[e_{p, I+q}^{\beta}, u_{i j}^{\alpha}\right]$ by Lemma 2.2.1, it follows that the $u_{i j}^{\alpha}$ provide the necessary generators in $U_{(2)}(\bar{p})^{N}$.

Even in this simple case, the commutation relations among the generators are quite complicated:

$$
\begin{aligned}
& {\left[u_{i j}^{\alpha}, e_{k \ell}^{\beta}+e_{I+k}, I+\ell\right.} \\
& {\left[u_{i j}^{\alpha}, e_{I+k, I+r}^{\beta}\right]=-u_{i \ell}^{\alpha \beta} \delta_{j k}-u_{k j}^{\beta \alpha} \delta_{i \ell}^{\beta \alpha} \delta_{i r}-\sum\left(e_{k j}^{\beta \gamma \alpha}+e_{I+k, I+j}^{\beta \gamma \alpha}\right) e_{I+i, I+r}^{\bar{\gamma}}} \\
& \quad(i f k \leq I) \\
& +\sum_{\tau, \gamma}^{e_{2 I+\tau, I+r}^{\gamma \alpha \beta} e_{I+i, 2 I+\tau}^{\bar{\gamma}} \delta_{j k}}
\end{aligned}
$$

and the commutation relations among the $u_{i j}^{\alpha}$ 's also involve a sum including the various forms of the generators. Thus, although $O y^{f}$ was useful in calculating the generators of $U(\bar{p})^{\mathrm{N}}$, it really provided very little information about the algebra structure of $U(\bar{p})^{\mathrm{N}}$.\\
8.2. Some special parabolic subalgebras of $g \ell(n ; K)$

Let $O y_{0}=g \ell(n ; K), K=R, C$, or $H$ and let $\Sigma \leq \Pi$ be of the form $\left\{v_{i}-v_{i+1} \mid i=1, \ldots, I\right\}$. Clearly, we may take $f=\sum_{i=1}^{I} e_{i+1, i}$ as a representative real admissible nilpotent element. Then, $\mathcal{G}^{f}$ is spanned by the $\sum_{i=1}^{I+l} \alpha e_{i i}$ and the $\alpha e_{I+r, I+s}$ where $\alpha \varepsilon K$ and $1 \leq r \leq n-I$, $l<s \leq n-I$. Further, $\bar{\pi}^{f}$ is spanned by the I+i-r $\sum_{i=1} \alpha e_{i+r, i}$ and the $\alpha e_{I+s, 1}$ where $\underset{I+1-r}{r=1, \ldots, I,} s=I+1, \ldots, n-I$ and $\alpha \varepsilon K$. Clearly $\sum_{i=1} \alpha e_{i+r, i} \varepsilon d_{-r}$ and $\alpha e_{I+s, 1} \varepsilon d_{I}$ so that there are $\operatorname{dim} K$ generators for\\
$U(\bar{p})^{N}$ in dimension $r$ if $I \leq r \therefore I$ and $(n-I) \operatorname{dim} K$ generators in dimension $I$. Furthermore, if $\alpha e_{I+r, I+s} \varepsilon M^{f}$, then $\left[\alpha e_{I+r, I+s}, \beta e_{I+t, 1]}=\delta_{s+} \alpha \beta e_{I+r, 1}\right.$ so that we may take the generators of $U(\bar{p})^{\mathrm{N}}$ corresponding to $\left\{\beta e_{I+t, 1}\right\}$ to be a cyclic $m^{f}$-module. (It is not, however, irreducible.)

Next, note that if $\mathrm{K}=\mathrm{H}$, then\\
$\left[\sum_{i=1}^{I+1-r} \alpha e_{i+r, i}, \sum_{j=1}^{I+1-s} \beta e_{j+s, j}\right]=\sum_{k=1}^{I+1-(r+s)}[\alpha, \beta] e_{k+r+s, k}$ so that\\
given generators corresponding to $\alpha f, \operatorname{Re}(\alpha) \neq 0$, we may construct all of the other generators coming from $\bar{n}^{f}$ having $\operatorname{Re}(\alpha)=0$.

The only remaining generators of $U(\bar{\rho})^{\mathrm{N}}$ that we need treat are those corresponding to the $\sum_{i=1}^{I+1-r} e_{i+r, i}$. However, these are just the images of the corresponding generators of the center of $U(O \gamma)$ in the appropriate dimensions. (To see this, one need only recall that $\}(g)$ for $g \ell(n ; C)$ is generated by the trace-invariants.)\\
8.3 Minimal parabolic subalgebras of so $(p, q)$ and $s u(p, q), p>q$.

We keep the notation of §7.2. Thus, $g(p, q ; K)$ has the real basis given by (7.2.2). As we showed in §7.3, we can take $f=\sum_{i=1}^{q-1} X_{v_{i+1}-v_{i}}+X_{-v_{q}}^{1,1}$ as a representative real admissible nilpotent element.

One easily sees that $m^{f}$ is spanned by the elements\\
(8.3.1) $\quad\left\{M_{\rho \sigma}^{\alpha} \mid p-q \geq \rho \geq \sigma \geq 1\right\} \quad\left\{\sum_{i=1}^{C_{\mu}} T_{i}^{\alpha}+M_{1,1}^{\alpha} \mid \alpha=\sqrt{-1}\right\}$.

One also sees that $\bar{n}^{-f}$ is spanned by\\
(8.3.2) $\quad\left\{\left.\sum_{i=1}^{q-r} Y_{v_{i+r}}^{\alpha}-v_{i}+Y_{v_{q}-r+1}^{\bar{\alpha}, 1}-\frac{1}{2} \sum_{p=0}^{r-2} Y_{-\left(v_{q-r+2+p}+v_{q-p}\right)}(-1)^{p} \right\rvert\,\right.$\\
$r=1, \ldots, 2 q$ and $\alpha=1$ if $r$ is even, $\alpha=0$ if $r$ is odd $\}$

$$
\left\{Y_{-v_{1}}^{\alpha, \rho} \mid 1<\rho \leq p-q\right\} .
$$

As in §8.2, the elements corresponding to the first set of brackets in (8.3.2) come from the center of $\mathcal{U}(g)$ as one readily checks. (As in §8.2, consider the form of elements in $\mathcal{\gamma}(\mathcal{g})$ and notice that their image under\\
$\pi_{1}{ }^{T}(\cdot)^{\rho_{\eta}}$ must have the above form. The result then follows by counting the number of generators of $z(o g)$ in each dimension.

The generators corresponding to the elements $Y_{-V_{1}}^{\alpha, \rho}$ are of special interest since $\left[Y_{-v_{1}}^{\alpha, \rho}, Y_{-v_{1}}^{\beta, \rho}\right] \varepsilon \mathbb{R} Y_{-2 v_{1}}^{\alpha \beta}$. That is, their commutator corresponds (or at least tries to) to an element in $U(\bar{\rho})^{N}$ arising from a highest weight. Thus, they appear to be some sort of generalized Pfaffian elements.

\section{References}
\begin{enumerate}
  \item Bernstein, I. N., Gelfand, I. M., Gelfand, S. I.: Structure of representations generated by highest weight vectors, Funct. Anal. i evo Prilojence, 5, 1-9, (1971) .
  \item Dixmier, J.: Algebres Enveloppantes, Cahiers scientifiques, vol. 37, Paris: Gauthier-Villars, 1974.
  \item Dynkin, E. B.: Semi-simple subalgebras of semi-simple Lie algebras, American Mathematical Society Translations, vol. 6, 111-244 (1957).
  \item Hesselink, W: Polarizations in the classical groups, Math. Zeit, 160, 217-234 (1978).
  \item Jantzen, J: Kontravariante Formen auf induzierten Darstellungen halbeinfacher Lie-algebren, Math. Ann., 266, 53-65 (1977).
  \item Kostant, B.: The principal three dimensional subgroup and the Betti numbers of a complex semi-simple Lie group, Amer. J. Math., 81, 973-1032 (1959).
  \item Kostant, B.: Lie group representations on polynomial rings, Amer. J. Math., 85, 327-404 (1963).
  \item Kostant, B.: On the existence and irreducibility of certain series of representations, in Lie Groups and Their Representations, edited by I. M. Gelfand.
\end{enumerate}

John Wiley and Sons, 231-329, 1975.\\
9. Kostant, B.: On Whittaker Vectors and Representation Theory, Inventiones Math., 48, 101-184 (1978).\\
10. Ozeki, H., Wakimoto, M.: On polarizations of certain homogeneous spaces, Hiroshima Math. J., 2, 445-482 (1972) .\\
11. Richardson, R.: Conjugacy classes in parabolic subgroups of semi-simple algebraic groups, Bull. London Math. Soc., 6, 21-24 (1974).\\
12. Speh, B., Vogan, D.: Reducibility of generalized principal series representations. To appear.\\
13. Springer, T. A., Steinberg, R.: Conjugacy classes. In: Borel, et al.: Seminar on Algebraic Groups and Related Finite Groups. Lecture Notes in Mathematics, 131. Berlin-Heidelberg-New York: Springer, 167-266, 1970.\\
14. Steinberg, R.: Conjugacy Classes in Algebraic Groups. Lecture Notes in Mathematics, 366. Berlin-HeidelbergNew York: Springer, 1974.\\
15. Vogan, D.: Gelfand-Kirillov dimension for Harish-Chandra modules, Institute for Advanced Study, Preprint.\\
16. Warner, G.: Harmonic Anałysis on Semi-Simple Lie Groups, vol. 1. Die Grundlehren der mathematischen Wissenschaften, Band 188, Berlin: Springer-Verlag, 1972.\\
17. Wallach, N. R.: Cyclic vectors and irreducibility for principal series representations, I, Trans. Amer. Math. Soc., 158, 107-113; II, ibid., 164, 389-396 (1972) .\\
18. Harish-Chandra: Harmonic analysis on real reductive groups III The Maass-Selberg relations and the Plancherel formula, Annals of Math., 104, no. 1, 117-201 (1976).

\section{Biographical note}
I was born in the naval hospital in Vallejo, California in 1953. My family spent the next 19 years around naval bases in Pennsylvania, California, Virginia, Louisiana and, finally, Georgia where I graduated from high school. I spent 1971-3 as an undergraduate at the University of Georgia and 1973-6 as a graduate student there working on stochastic differential equations. I came to M.I.T. in 1976 to work on Lie groups. Next year I'll be teaching at SUNY-Albany.


\end{document}